% Mouvement de particules chargées dans les champs électrique et magnétique uniformes — Physique Tle C — Cours signé OnBuch+
% Programme officiel MINESEC. Compiler avec : tectonic P08-particules-chargees-champs-uniformes.tex
\newcommand{\DOCMATIERE}{Physique}
\newcommand{\DOCNIVEAU}{Tle C}
\newcommand{\DOCMODULE}{Module 2 : Mouvements et interactions — évolutions temporelles des systèmes mécaniques}
\newcommand{\DOCLECON}{P08}
\newcommand{\DOCTITRE}{Mouvement de particules chargées dans les champs électrique et magnétique uniformes}
% OnBuch+ — préambule « cours signés » ENRICHI (compilé avec tectonic / XeLaTeX)
% Système visuel « Pop » d'OnBuch+ : fond crème (jamais blanc), bordures encre
% épaisses, ombres « dures » décalées, titres Archivo Black en capitales.
% Version MATHÉMATIQUES : graphes (pgfplots/tikz), boîtes pédagogiques
% (définition, propriété, méthode, attention, le savais-tu, exercice résolu,
% exercice, TP, bilan), page de garde et signature OnBuch+ ancrée en bas.
%
% Macros attendues dans main.tex AVANT \input{preamble} :
%   \DOCMATIERE  \DOCNIVEAU  \DOCMODULE  \DOCLECON (numéro)  \DOCTITRE
\documentclass[11pt]{article}

\usepackage{fontspec}
\usepackage[french]{babel}
\usepackage[a4paper,margin=2cm,top=3.4cm,bottom=2.8cm,headheight=1.4cm,footskip=1.3cm]{geometry}
\usepackage{amsmath,amssymb,amsfonts}
\usepackage{mathtools} % \xmapsto, \coloneqq... souvent utilisés par les modèles
\usepackage{cancel} % simplifications barrées (bilans de forces)
% \abs{z} (module, valeur absolue) et \norm{u} (norme), utilisés par les modèles
\providecommand{\abs}{}\let\abs\relax\DeclarePairedDelimiter{\abs}{\lvert}{\rvert}
\providecommand{\norm}{}\let\norm\relax\DeclarePairedDelimiter{\norm}{\lVert}{\rVert}
\usepackage[table]{xcolor} % [table] : fournit aussi \rowcolors (alternance de lignes)
\usepackage{pagecolor}
\usepackage{tikz}
\usetikzlibrary{arrows.meta,positioning,calc,shapes.geometric,shapes.misc,shapes.arrows,decorations.pathmorphing,decorations.markings,patterns,shadows,mindmap,trees,fit,backgrounds,matrix,intersections}
\usepackage{pgfplots}
\usepackage[version=4]{mhchem} % équations nucléaires \ce{^{235}_{92}U}
\usepackage[european,siunitx]{circuitikz} % schémas électriques
\pgfplotsset{compat=1.18}
\usepgfplotslibrary{fillbetween} % aires entre courbes (calcul intégral)
\usepackage{enumitem}
\usepackage{fancyhdr}
% [most] charge toutes les bibliothèques tcolorbox (listings, minted, raster,
% documentation...) et gonfle inutilement la mémoire TeX sur un document long
% (~300 boîtes) : on ne charge que ce qui est réellement utilisé.
\usepackage[skins,breakable]{tcolorbox}
\usepackage{graphicx}
\usepackage{colortbl}
\usepackage{booktabs}
\usepackage{tabularx}
\usepackage{diagbox} % case d'en-tête diagonale des tableaux à double entrée
% Colonnes extensibles centrée / gauche / droite (C, L, R), souvent utilisées par les modèles
\newcolumntype{C}{>{\centering\arraybackslash}X}
\newcolumntype{L}{>{\raggedright\arraybackslash}X}
\newcolumntype{R}{>{\raggedleft\arraybackslash}X}
\usepackage{array}
\usepackage{multicol}
\usepackage{siunitx}
% parse-numbers=false : accepte aussi \SI{2,5 \times 10^{-3}}{...}
\sisetup{locale=FR,output-decimal-marker={,},group-minimum-digits=4,parse-numbers=false}
\DeclareSIUnit\division{div} % réglages d'oscilloscope (ms/div, V/div)
\usepackage{qrcode}
% \degree : commande fréquemment utilisée par les modèles pour un angle ou une
% température en dehors de \SI{}{} (non fournie par les paquets chargés).
\providecommand{\degree}{^{\circ}}
% \qed (fin de démonstration), utilisé par les modèles sans amsthm
\providecommand{\qed}{\hfill$\square$}
% \barre : double barre des valeurs interdites dans un tableau de variations (tkz-tab)
\providecommand{\barre}{\ensuremath{\|}}
% environnement proof (amsthm non chargé) utilisé par les modèles
\ifdefined\proof\else\newenvironment{proof}[1][Démonstration]{\par\noindent\textit{#1.}\ }{\qed\par}\fi
\usepackage{lastpage}
\usepackage{hyperref}
\hypersetup{hidelinks}

% ============================================================
% Palette « Pop » OnBuch+ — mêmes hex que la classe P (plus_ui.dart)
% ============================================================
\definecolor{poporange}{HTML}{F59321}
\definecolor{popdark}{HTML}{E07A0C}
\definecolor{popcream}{HTML}{FFF6E8}
\definecolor{popink}{HTML}{1C1714}
\definecolor{poppurple}{HTML}{7A5AE0}
\definecolor{popgreen}{HTML}{2E9E6A}
\definecolor{poppink}{HTML}{E0457B}
\definecolor{popblue}{HTML}{2F7DE1}
\definecolor{popgold}{HTML}{C98A12}
\definecolor{popmuted}{HTML}{8A7F76}
\definecolor{popline}{HTML}{EADFD2}
\definecolor{popgray}{HTML}{8A7F76}
\colorlet{popgrey}{popgray}
% Alias tolérants (le modèle nomme parfois les couleurs différemment)
\colorlet{popwhite}{white}
\colorlet{popblack}{popink}
\colorlet{popred}{poppink}
\colorlet{popyellow}{popgold}
% Teintes claires pour fonds de tableaux / zones
\colorlet{popblueL}{popblue!10!white}
\colorlet{poporangeL}{poporange!14!white}
\colorlet{popgreenL}{popgreen!12!white}
\colorlet{poppurpleL}{poppurple!12!white}
\colorlet{poppinkL}{poppink!12!white}
\colorlet{popgoldL}{popgold!14!white}

\pagecolor{popcream}

% ============================================================
% Typographie
% ============================================================
\defaultfontfeatures{Ligatures=TeX}
\setmainfont{PlusJakartaSans-Medium.ttf}[Path=../../../fonts/,BoldFont=PlusJakartaSans-Bold.ttf]
% Maths en sans-serif (Fira Math) avec chiffres et lettres droites de Plus
% Jakarta, pour que les formules se fondent dans le texte.
\usepackage[math-style=ISO,bold-style=ISO]{unicode-math}
\setmathfont{FiraMath-Regular.otf}[Path=../../../fonts/]
\setmathfont{PlusJakartaSans-Medium.ttf}[Path=../../../fonts/,range={up/num,up/{Latin,latin},\mathup}]
\usepackage{icomma} % virgule décimale sans espace en mode math
% Caractères mathématiques Unicode absents des polices de texte (Plus Jakarta,
% Archivo Black) : rendus par leur équivalent mathématique, en texte comme en
% formule — sinon ils s'affichent en carré vide (ex. « Arithmétique dans ℤ »).
\usepackage{newunicodechar}
\newunicodechar{ℝ}{\ensuremath{\mathbb{R}}}
\newunicodechar{ℤ}{\ensuremath{\mathbb{Z}}}
\newunicodechar{ℂ}{\ensuremath{\mathbb{C}}}
\newunicodechar{ℕ}{\ensuremath{\mathbb{N}}}
\newunicodechar{ℚ}{\ensuremath{\mathbb{Q}}}
\newunicodechar{𝔻}{\ensuremath{\mathbb{D}}}
\newunicodechar{α}{\ensuremath{\alpha}}
\newunicodechar{β}{\ensuremath{\beta}}
\newunicodechar{γ}{\ensuremath{\gamma}}
\newunicodechar{δ}{\ensuremath{\delta}}
\newunicodechar{ε}{\ensuremath{\varepsilon}}
\newunicodechar{θ}{\ensuremath{\theta}}
\newunicodechar{λ}{\ensuremath{\lambda}}
\newunicodechar{μ}{\ensuremath{\mu}}
\newunicodechar{σ}{\ensuremath{\sigma}}
\newunicodechar{τ}{\ensuremath{\tau}}
\newunicodechar{φ}{\ensuremath{\varphi}}
\newunicodechar{ψ}{\ensuremath{\psi}}
\newunicodechar{ω}{\ensuremath{\omega}}
\newunicodechar{Σ}{\ensuremath{\Sigma}}
% majuscules produites par \MakeUppercase dans les titres (α-aminés -> Α)
\newunicodechar{Α}{\ensuremath{\alpha}}
\newunicodechar{Β}{\ensuremath{\beta}}
\newunicodechar{Γ}{\ensuremath{\Gamma}}
\newunicodechar{Δ}{\ensuremath{\Delta}}
\newunicodechar{Ε}{\ensuremath{\varepsilon}}
\newunicodechar{Θ}{\ensuremath{\Theta}}
\newunicodechar{Λ}{\ensuremath{\Lambda}}
\newunicodechar{Μ}{\ensuremath{\mu}}
\newunicodechar{Τ}{\ensuremath{\tau}}
\newunicodechar{Φ}{\ensuremath{\Phi}}
\newunicodechar{Ψ}{\ensuremath{\Psi}}
\newunicodechar{Ω}{\ensuremath{\Omega}}
\newunicodechar{↦}{\ensuremath{\mapsto}}
\newunicodechar{∈}{\ensuremath{\in}}
\newunicodechar{∉}{\ensuremath{\notin}}
\newunicodechar{∧}{\ensuremath{\wedge}}
\newunicodechar{⇔}{\ensuremath{\Leftrightarrow}}
\newunicodechar{⟺}{\ensuremath{\Longleftrightarrow}}
\newunicodechar{⇒}{\ensuremath{\Rightarrow}}
\newunicodechar{⟹}{\ensuremath{\Longrightarrow}}
\newunicodechar{∘}{\ensuremath{\circ}}
\newunicodechar{∪}{\ensuremath{\cup}}
\newunicodechar{∩}{\ensuremath{\cap}}
\newunicodechar{≡}{\ensuremath{\equiv}}
\newunicodechar{⊂}{\ensuremath{\subset}}
\newunicodechar{⊥}{\ensuremath{\perp}}
\newunicodechar{∃}{\ensuremath{\exists}}
\newunicodechar{∀}{\ensuremath{\forall}}
\newunicodechar{∛}{\ensuremath{\sqrt[3]{\,}}}
\newunicodechar{∜}{\ensuremath{\sqrt[4]{\,}}}
\newunicodechar{⃗}{\ensuremath{{}^{\to}}}
\newunicodechar{ⁿ}{\ifmmode^{n}\else\textsuperscript{\ensuremath{n}}\fi}
\newunicodechar{ˣ}{\ifmmode^{x}\else\textsuperscript{\ensuremath{x}}\fi}
\newunicodechar{ᵇ}{\ifmmode^{b}\else\textsuperscript{\ensuremath{b}}\fi}
\newunicodechar{ʳ}{\ifmmode^{r}\else\textsuperscript{\ensuremath{r}}\fi}
\newunicodechar{ᵘ}{\ifmmode^{u}\else\textsuperscript{\ensuremath{u}}\fi}
\newunicodechar{ʸ}{\ifmmode^{y}\else\textsuperscript{\ensuremath{y}}\fi}
\newunicodechar{ᵃ}{\ifmmode^{a}\else\textsuperscript{\ensuremath{a}}\fi}
\newunicodechar{ᵉ}{\ifmmode^{e}\else\textsuperscript{\ensuremath{e}}\fi}
\newunicodechar{ᵏ}{\ifmmode^{k}\else\textsuperscript{\ensuremath{k}}\fi}
\newunicodechar{ᵖ}{\ifmmode^{p}\else\textsuperscript{\ensuremath{p}}\fi}
\newunicodechar{ᴺ}{\ifmmode^{N}\else\textsuperscript{\ensuremath{N}}\fi}
\newunicodechar{ᶜ}{\ifmmode^{c}\else\textsuperscript{\ensuremath{c}}\fi}
\newunicodechar{ʰ}{\ifmmode^{h}\else\textsuperscript{\ensuremath{h}}\fi}
\newunicodechar{ᵠ}{\ifmmode^{q}\else\textsuperscript{\ensuremath{q}}\fi}
\newunicodechar{ₙ}{\ifmmode_{n}\else\textsubscript{\ensuremath{n}}\fi}
\newunicodechar{ₐ}{\ifmmode_{a}\else\textsubscript{\ensuremath{a}}\fi}
\newunicodechar{ᵢ}{\ifmmode_{i}\else\textsubscript{\ensuremath{i}}\fi}
\newunicodechar{ᵣ}{\ifmmode_{r}\else\textsubscript{\ensuremath{r}}\fi}
\newunicodechar{ₖ}{\ifmmode_{k}\else\textsubscript{\ensuremath{k}}\fi}
\newunicodechar{ᵦ}{\ifmmode_{\beta}\else\textsubscript{\ensuremath{\beta}}\fi}
\newunicodechar{ₓ}{\ifmmode_{x}\else\textsubscript{\ensuremath{x}}\fi}
\newfontfamily\popdisplay{ArchivoBlack-Regular.ttf}[Path=../../../fonts/]
\newfontfamily\poplogo{SpaceGrotesk-Bold.ttf}[Path=../../../fonts/]
\newfontfamily\popsemi{PlusJakartaSans-SemiBold.ttf}[Path=../../../fonts/]
\newfontfamily\popxbold{PlusJakartaSans-ExtraBold.ttf}[Path=../../../fonts/]
\newfontfamily\popxboldls{PlusJakartaSans-ExtraBold.ttf}[Path=../../../fonts/,LetterSpace=3.0]

\setlist[enumerate]{leftmargin=1.7em,itemsep=5pt,topsep=4pt}
\setlist[itemize]{leftmargin=1.7em,itemsep=4pt,topsep=4pt,label={\color{poporange}\textbullet}}
\setlength{\parindent}{0pt}
\setlength{\parskip}{5pt}
\renewcommand{\arraystretch}{1.25}

% pgfplots : style Pop par défaut
\pgfplotsset{
  every axis/.append style={axis line style={popink,line width=1pt},tick style={popink},
    grid style={popline,line width=0.5pt},label style={font=\small},tick label style={font=\footnotesize}},
  popcurve/.style={poporange,line width=1.8pt,no markers,smooth},
}
% Même style disponible dans un \draw TikZ simple (hors pgfplots)
\tikzset{popcurve/.style={poporange,line width=1.8pt}}
\tikzset{popgrid/.style={popline,very thin}}

% ============================================================
% Pill / squiggle
% ============================================================
\newcommand{\poppill}[3]{%
  \tikz[baseline=-0.55ex]\node[draw=popink,line width=1.2pt,rounded corners=2.6pt,%
    fill=#2,inner xsep=6.5pt,inner ysep=3pt]{%
    \popxboldls\fontsize{7}{7}\selectfont\color{#3}\MakeUppercase{#1}};%
}
\newcommand{\popsquiggle}[1]{%
  \tikz[baseline=-0.4ex]{
    \draw[#1,line width=2.6pt,line cap=round]
      plot [smooth] coordinates {(0,0.09) (0.34,0.01) (0.68,0.21) (1.02,0.01) (1.36,0.10)};
  }%
}
% Mot-clé surligné (marqueur orange sous le texte)
\newcommand{\cle}[1]{{\popxbold\color{popdark}#1}}

% ============================================================
% En-tête / pied
% ============================================================
\pagestyle{fancy}
\fancyhf{}
\renewcommand{\headrulewidth}{0pt}
\renewcommand{\footrulewidth}{0pt}
\lhead{\raisebox{-0.15ex}{\poplogo\fontsize{13}{13}\selectfont\color{popink}OnBuch\color{poporange}+}}
\rhead{\poppill{\DOCMATIERE\ \textbullet\ \DOCNIVEAU\ \textbullet\ Leçon \DOCLECON}{white}{popink}}
\lfoot{\popxbold\fontsize{7.6}{7.6}\selectfont\color{popmuted}\MakeUppercase{Cours signé OnBuch+}\ \textbullet\ {\popsemi\DOCTITRE}}
\rfoot{\tikz[baseline=-0.6ex]\node[rounded corners=8.5pt,fill=popink,minimum height=17pt,minimum width=17pt,inner xsep=5pt,inner sep=0]{\popxbold\fontsize{8}{8}\selectfont\color{popcream}\thepage\,/\,\pageref*{LastPage}};}

% ============================================================
% Titres
% ============================================================
\newcommand{\chaptitre}[2]{%
  \begin{tcolorbox}[enhanced,colback=poporange,colframe=popink,boxrule=2.6pt,arc=11pt,
      drop shadow=popink,left=15pt,right=15pt,top=12pt,bottom=11pt,before skip=3pt,after skip=18pt]
    \poppill{#2}{popink}{popcream}\par\vspace{9pt}
    {\raggedright\hyphenpenalty=10000\exhyphenpenalty=10000\popdisplay\fontsize{22}{24}\selectfont\color{popcream}\MakeUppercase{#1}\par}
  \end{tcolorbox}
}
% Section numérotée : pastille orange avec le numéro + titre Archivo
\newcounter{popsec}
\newcommand{\coursec}[1]{%
  \stepcounter{popsec}\setcounter{popsub}{0}%
  \par\vspace{16pt}\noindent
  \begin{minipage}{\linewidth}
  \tikz[baseline=-0.7ex]\node[draw=popink,line width=1.6pt,fill=poporange,rounded corners=4pt,minimum size=22pt,inner sep=0]{\popdisplay\fontsize{12}{12}\selectfont\color{popcream}\thepopsec};%
  \hspace{8pt}{\popdisplay\fontsize{15}{17}\selectfont\color{popink}\MakeUppercase{#1}}\par
  \vspace{4pt}\noindent{\color{poporange}\rule{36pt}{3.6pt}}
  \end{minipage}\par\nopagebreak\vspace{8pt}
}
\newcounter{popsub}
\newcommand{\courssub}[1]{%
  \stepcounter{popsub}%
  \par\vspace{9pt}\noindent{\popxbold\fontsize{11.5}{13}\selectfont\color{popink}{\color{poporange}\thepopsec.\thepopsub}\hspace{6pt}#1}\par\nopagebreak\vspace{4pt}
}

% ============================================================
% Boîtes pédagogiques — même recette : fond blanc, bordure épaisse colorée,
% ombre dure encre, badge plein. Titre optionnel [#1].
% ============================================================
\tcbset{popbase/.style={breakable,enhanced,colback=white,boxrule=2.3pt,arc=9pt,
  shadow={3pt}{-3pt}{0pt}{popink},left=14pt,right=14pt,top=11pt,bottom=11pt,before skip=11pt,after skip=15pt,
  fontupper=\popsemi\fontsize{10.6}{14.5}\selectfont\color{popink},
  fontlower=\popsemi\fontsize{10.6}{14.5}\selectfont\color{popink}}}
% Coche dessinée (le glyphe ✓ n'existe pas dans les polices du document)
\newcommand{\popcheck}[1]{\tikz[baseline=-0.5ex]\draw[#1,line width=1.6pt,line cap=round,line join=round](-0.13,0.0)--(-0.04,-0.09)--(0.13,0.1);}
\providecommand{\coche}{\popcheck{popgreen}}
\newcommand{\popboxhead}[3]{\poppill{#1}{#2}{white}\ifx\relax#3\relax\else\hspace{6pt}{\popxbold\fontsize{10.6}{12}\selectfont\color{popink}#3}\fi\par\vspace{7pt}}

\newtcolorbox{aretenir}[1][]{popbase,colframe=popgreen,before upper={\popboxhead{À retenir}{popgreen}{#1}}}
\newtcolorbox{exemplebox}[1][]{popbase,colframe=poporange,before upper={\popboxhead{Exemple}{poporange}{#1}}}
\newtcolorbox{definition}[1][]{popbase,colframe=popblue,before upper={\popboxhead{Définition}{popblue}{#1}}}
\newtcolorbox{propriete}[1][]{popbase,colframe=poppurple,before upper={\popboxhead{Propriété}{poppurple}{#1}}}
\newtcolorbox{methode}[1][]{popbase,colframe=popgold,before upper={\popboxhead{Méthode}{popgold}{#1}}}
\newtcolorbox{attention}[1][]{popbase,colframe=poppink,before upper={\popboxhead{Attention}{poppink}{#1}}}
\newtcolorbox{experience}[1][]{popbase,colframe=popblue,colback=popblueL,before upper={\popboxhead{Expérience}{popblue}{#1}}}
% « Le savais-tu ? » : carte violette pleine (contexte, culture, Cameroun)
\newtcolorbox{savaistu}[1][]{popbase,colback=poppurple,colframe=popink,
  fontupper=\popsemi\fontsize{10.4}{14.2}\selectfont\color{white},
  before upper={\poppill{Le savais-tu\,?}{white}{poppurple}\ifx\relax#1\relax\else\hspace{6pt}{\popxbold\color{white}#1}\fi\par\vspace{7pt}}}
% Exercice résolu : énoncé en haut, \tcblower puis la solution sur fond crème
\newcounter{popexo}\newcounter{popexr}
\newtcolorbox{exoresolu}[1][]{popbase,colframe=poporange,colbacklower=popcream,
  segmentation style={popink,line width=1.2pt,dashed},
  before upper={\stepcounter{popexr}\popboxhead{Exercice résolu \thepopexr}{poporange}{#1}},
  before lower={\poppill{Solution}{popink}{popcream}\par\vspace{6pt}}}
% Exercice (non résolu en place) — la correction est dans la partie Corrigés
\newtcolorbox{exercice}[1][]{popbase,colframe=popink,
  before upper={\stepcounter{popexo}\popboxhead{Exercice \thepopexo}{popink}{#1}}}
% Corrigé
\newtcolorbox{corrige}[1][]{popbase,colframe=popgreen,colback=popgreenL,
  before upper={\popboxhead{Corrigé}{popgreen}{#1}}}
% Objectifs / prérequis / bilan
\newtcolorbox{objectifs}{popbase,colframe=popink,colback=poporangeL,
  before upper={\popboxhead{Objectifs de la leçon}{popink}{}}}
\newtcolorbox{prerequis}{popbase,colframe=popink,colback=popblueL,
  before upper={\popboxhead{Prérequis}{popblue}{}}}
\newtcolorbox{bilan}{popbase,colframe=popink,colback=white,boxrule=2.8pt,
  before upper={\popboxhead{Fiche bilan}{poporange}{L'essentiel à connaître pour l'examen}}}
% Figure encadrée avec légende
\newtcolorbox{popfigure}[1][]{enhanced,breakable=false,colback=white,colframe=popline,boxrule=1.4pt,arc=8pt,
  left=10pt,right=10pt,top=10pt,bottom=8pt,before skip=10pt,after skip=12pt,halign=center,
  lower separated=false,halign lower=center,
  fontlower=\popsemi\fontsize{9}{11.5}\selectfont\color{popmuted},
  #1}
\newcommand{\legende}[1]{\par\vspace{6pt}{\popsemi\fontsize{9}{11.5}\selectfont\color{popmuted}#1}}

% ============================================================
% Signature OnBuch+
% ============================================================
\newcommand{\poplogoinline}[1]{{\poplogo\fontsize{#1}{#1}\selectfont\color{popink}OnBuch\color{poporange}+}}
\newcommand{\signatureonbuch}{%
  \begin{tcolorbox}[enhanced,colback=popink,colframe=popink,boxrule=2.2pt,arc=10pt,
      drop shadow=poporange,left=14pt,right=14pt,top=10pt,bottom=10pt,before skip=0pt,after skip=0pt]
    \tikz[baseline=-0.7ex]\node[circle,fill=popgreen,draw=popcream,line width=1.3pt,minimum size=22pt,inner sep=0]{\popcheck{white}};%
    \hspace{9pt}\begin{minipage}[c]{0.62\linewidth}
      {\popxbold\fontsize{12}{13}\selectfont\color{popcream}Signé\ \textbullet\ {\poplogo OnBuch\color{poporange}+}}\par\vspace{2pt}
      {\raggedright\popsemi\fontsize{8}{10.5}\selectfont\color{popline}\MakeUppercase{Équipe pédagogique OnBuch+}\\\MakeUppercase{Conforme au programme officiel MINESEC}\par}
    \end{minipage}\hfill
    {\poplogo\fontsize{22}{22}\selectfont\color{popcream}OnBuch\color{poporange}+}
  \end{tcolorbox}%
}

% ============================================================
% Page de garde
% #1 titre · #2 matière · #3 niveau · #4 module · #5 n° leçon · #6 durée ·
% #7 description / objectifs (texte ou itemize)
% ============================================================
\newcommand{\pagegarde}[7]{%
  \thispagestyle{empty}
  \newgeometry{a4paper,margin=1.7cm,top=1.6cm,bottom=1.4cm}%
  \begin{tikzpicture}[remember picture,overlay]
    \fill[poporange!18] ([xshift=-3.2cm,yshift=-3.6cm]current page.north east) circle (4.6cm);
    \fill[poppurple!14] ([xshift=2.4cm,yshift=7.6cm]current page.south west) circle (3.8cm);
  \end{tikzpicture}%
  {\poplogo\fontsize{30}{30}\selectfont\color{popink}OnBuch\color{poporange}+}\hfill
  \raisebox{0.6ex}{\poppill{Cours signé \textbullet\ Premium}{popink}{popcream}}\par
  \vspace{1pt}\popsquiggle{poppurple}\par
  \vspace{8pt}
  \begin{tcolorbox}[enhanced,colback=poporange,colframe=popink,boxrule=2.8pt,arc=13pt,
      drop shadow=popink,left=18pt,right=18pt,top=12pt,bottom=13pt,after skip=10pt]
    \poppill{#2 \textbullet\ #3}{popink}{popcream}\hspace{5pt}\poppill{#4}{white}{popink}\par\vspace{8pt}
    {\popdisplay\fontsize{14}{14}\selectfont\color{popink}LEÇON #5}\par\vspace{4pt}
    {\raggedright\hyphenpenalty=10000\exhyphenpenalty=10000\popdisplay\fontsize{25}{27}\selectfont\color{popcream}\MakeUppercase{#1}\par}
    \vspace{8pt}
    {\popxbold\fontsize{9.5}{9.5}\selectfont\color{popink}\MakeUppercase{Durée conseillée : #6}}
  \end{tcolorbox}
  \begin{tcolorbox}[enhanced,colback=white,colframe=popink,boxrule=2.2pt,arc=10pt,
      drop shadow=popink,left=16pt,right=16pt,top=10pt,bottom=10pt,after skip=8pt]
    \poppill{Dans cette leçon}{poppurple}{white}\par\vspace{5pt}
    \popsemi\fontsize{9.3}{12.6}\selectfont\color{popink}#7
  \end{tcolorbox}
  \noindent\begin{minipage}[t]{0.487\linewidth}
    \begin{tcolorbox}[enhanced,colback=popgreen,colframe=popink,boxrule=2.2pt,arc=10pt,
        drop shadow=popink,left=11pt,right=11pt,top=10pt,bottom=10pt,height=3.1cm,valign=center]
      \tcbox[colback=white,colframe=popink,boxrule=1.3pt,arc=5pt,boxsep=0pt,nobeforeafter,
        left=3pt,right=3pt,top=3pt,bottom=3pt,valign=center]{\qrcode[height=1.9cm]{https://play.google.com/store/apps/details?id=com.luvvix.onbuch&hl=fr}}%
      \hspace{8pt}\begin{minipage}[c]{0.5\linewidth}
        {\popdisplay\fontsize{11}{12.5}\selectfont\color{white}\MakeUppercase{Télécharge OnBuch}}\par\vspace{4pt}
        {\raggedright\popsemi\fontsize{8.4}{11}\selectfont\color{white}Gratuit \textbullet\ Google Play\\Tuteur IA, annales, concours, résultats\par}
      \end{minipage}
    \end{tcolorbox}
  \end{minipage}\hfill
  \begin{minipage}[t]{0.487\linewidth}
    \begin{tcolorbox}[enhanced,colback=poppurple,colframe=popink,boxrule=2.2pt,arc=10pt,
        drop shadow=popink,left=11pt,right=11pt,top=10pt,bottom=10pt,height=3.1cm,valign=center]
      \tikz[baseline=-0.7ex]\node[circle,fill=white,draw=popink,line width=1.3pt,minimum size=1.6cm,inner sep=0]{\popdisplay\fontsize{24}{24}\selectfont\color{poppurple}+};%
      \hspace{8pt}\begin{minipage}[c]{0.6\linewidth}
        {\popdisplay\fontsize{11}{12.5}\selectfont\color{white}\MakeUppercase{Passe à OnBuch+}}\par\vspace{4pt}
        {\raggedright\popsemi\fontsize{8.4}{11}\selectfont\color{white}Tous les cours signés, Tuteur IA illimité, fiches PDF — onglet OnBuch+ de l'app\par}
      \end{minipage}
    \end{tcolorbox}
  \end{minipage}
  \vspace{8pt}
  \popcontenu
  \vfill
  \signatureonbuch
  \restoregeometry
  \newpage
}

% Rangée « Ce cours contient » (page de garde)
\newcommand{\poptile}[3]{% #1 couleur #2 gros texte #3 légende
  \begin{tcolorbox}[enhanced,colback=#1,colframe=popink,boxrule=2pt,arc=9pt,shadow={2.5pt}{-2.5pt}{0pt}{popink},
      left=6pt,right=6pt,top=9pt,bottom=9pt,halign=center,valign=center,height=1.75cm,width=\linewidth,nobeforeafter]
    {\popdisplay\fontsize{12.5}{13}\selectfont\color{white}\MakeUppercase{#2}}\par\vspace{3pt}
    {\centering\popsemi\fontsize{7.8}{9.5}\selectfont\color{white}#3\par}
  \end{tcolorbox}}
\newcommand{\popcontenu}{%
  \par\noindent{\popxboldls\fontsize{7.5}{7.5}\selectfont\color{popmuted}CE COURS SIGNÉ CONTIENT}\par\vspace{7pt}
  \noindent\begin{minipage}[t]{0.235\linewidth}\poptile{popblue}{Cours}{complet, illustré, pas à pas}\end{minipage}\hfill
  \begin{minipage}[t]{0.235\linewidth}\poptile{poporange}{Méthodes}{et exercices résolus}\end{minipage}\hfill
  \begin{minipage}[t]{0.235\linewidth}\poptile{poppink}{Exercices}{type examen + corrigés}\end{minipage}\hfill
  \begin{minipage}[t]{0.235\linewidth}\poptile{popgreen}{Fiche bilan}{l'essentiel à retenir}\end{minipage}\par}

% Sommaire
\newtcolorbox{sommaire}{enhanced,colback=white,colframe=popink,boxrule=2.2pt,arc=10pt,
  drop shadow=popink,left=18pt,right=18pt,top=13pt,bottom=13pt,before skip=6pt,after skip=20pt,
  before upper={\poppill{Sommaire}{poppurple}{white}\par\vspace{9pt}\popsemi\fontsize{10.6}{14.5}\selectfont\color{popink}}}

% Signature de fin de cours — poussée tout en bas de la dernière page
\newcommand{\findecours}{%
  \par\vfill\nopagebreak
  \begin{center}\popsquiggle{poporange}\end{center}\vspace{4pt}
  \signatureonbuch
}

%% FIGFIX-BEGIN v4 — correctifs globaux des figures (docs/onbuch-generation/tools/figfix)
\usepackage{adjustbox}\usepackage{etoolbox}
% toute figure TikZ/pgfplots plus large que la ligne est réduite à la largeur disponible
\BeforeBeginEnvironment{tikzpicture}{\begin{adjustbox}{max width=\linewidth}}
\AfterEndEnvironment{tikzpicture}{\end{adjustbox}}
% légendes pgfplots lisibles par-dessus les courbes
\pgfplotsset{every axis/.append style={legend style={fill=white,fill opacity=0.92,text opacity=1,draw=black!15,inner sep=3pt}}}
% étiquettes d'arêtes (syntaxe edge["..."]) avec fond blanc
\tikzset{every edge quotes/.append style={fill=white,inner sep=1.2pt}}
%% FIGFIX-END
\begin{document}

\pagegarde{Mouvement de particules chargées dans les champs électrique et magnétique uniformes}{Physique}{Tle C}{Module 2 : Mouvements et interactions — évolutions temporelles des systèmes mécaniques}{P08}{6 h}{Cette leçon étudie le mouvement d'une particule chargée dans un champ électrique uniforme (accélération, déviation parabolique) puis dans un champ magnétique uniforme (mouvement circulaire uniforme). Elle débouche sur des applications majeures du programme : oscilloscope, spectrographe de masse, cyclotron et filtre de Wien.
\vspace{4pt}
\begin{multicols}{2}\small
\begin{itemize}
\item À la fin de cette leçon, je sais appliquer le théorème de l'énergie cinétique pour calculer la vitesse acquise par une particule chargée accélérée par une tension (canon à électrons)
\item À la fin de cette leçon, je sais établir les équations horaires du mouvement d'une particule chargée dans un champ électrique uniforme
\item À la fin de cette leçon, je sais montrer que la trajectoire dans un condensateur plan est parabolique et calculer la déflexion sur un écran
\item À la fin de cette leçon, je sais démontrer qu'une particule chargée animée d'une vitesse perpendiculaire à un champ magnétique uniforme décrit un mouvement circulaire uniforme
\item À la fin de cette leçon, je sais calculer le rayon R = mv/(|q|B) et la période du mouvement circulaire
\item À la fin de cette leçon, je sais expliquer et exploiter quantitativement le principe du spectrographe de masse
\end{itemize}\end{multicols}}

\begin{sommaire}
\begin{multicols}{2}
\begin{enumerate}
\item Accélération d'une particule chargée par un champ électrique uniforme
\item Déviation électrique : trajectoire parabolique
\item Déflexion sur un écran : l'oscilloscope
\item Mouvement dans un champ magnétique uniforme
\item Applications : spectrographe de masse et cyclotron
\item Filtre de Wien : champs électrique et magnétique croisés
\item Activité d'intégration
\item Méthodes et exercices résolus
\item Exercices
\item Corrigés des exercices
\item Fiche bilan
\end{enumerate}
\end{multicols}
\end{sommaire}

\begin{objectifs}
\begin{itemize}
\item À la fin de cette leçon, je sais appliquer le théorème de l'énergie cinétique pour calculer la vitesse acquise par une particule chargée accélérée par une tension (canon à électrons)
\item À la fin de cette leçon, je sais établir les équations horaires du mouvement d'une particule chargée dans un champ électrique uniforme
\item À la fin de cette leçon, je sais montrer que la trajectoire dans un condensateur plan est parabolique et calculer la déflexion sur un écran
\item À la fin de cette leçon, je sais démontrer qu'une particule chargée animée d'une vitesse perpendiculaire à un champ magnétique uniforme décrit un mouvement circulaire uniforme
\item À la fin de cette leçon, je sais calculer le rayon R = mv/(|q|B) et la période du mouvement circulaire
\item À la fin de cette leçon, je sais expliquer et exploiter quantitativement le principe du spectrographe de masse
\item À la fin de cette leçon, je sais expliquer le fonctionnement du cyclotron et calculer la fréquence cyclotron
\item À la fin de cette leçon, je sais déterminer la condition de sélection d'un filtre de Wien (champs croisés)
\end{itemize}
\end{objectifs}

\begin{prerequis}
\begin{itemize}
\item Force électrostatique F = qE et champ électrique uniforme entre deux armatures (Première)
\item Force de Lorentz magnétique F = qv∧B et règle d'orientation (Première)
\item Théorème de l'énergie cinétique et travail d'une force (Première)
\item Mouvement circulaire uniforme : accélération normale a = v²/R, période (Première)
\item Deuxième loi de Newton et projections sur des axes (début de Terminale)
\end{itemize}
\end{prerequis}

\newpage

\coursec{Accélération d'une particule chargée par un champ électrique uniforme}

À l'hôpital régional de Douala, un technicien règle un appareil de radiographie : des électrons sont accélérés puis dirigés avec une précision extrême vers une cible de tungstène, dont le choc produit les rayons X qui traverseront le thorax du patient. Dans les vieux téléviseurs à tube cathodique encore utilisés dans certains quartiers de Yaoundé, un faisceau d'électrons balaye l'écran pour dessiner l'image, image après image. À l'entrée de ces deux machines se trouve le même dispositif : un \cle{canon à électrons}.

Mais comment peut-on accélérer, dévier et contrôler des particules invisibles à l'œil nu, mille fois plus légères qu'un grain de poussière ? On ne peut pas les pousser avec une baguette, ni les lancer avec une fronde. La réponse tient en deux outils que la nature met à notre disposition : le \cle{champ électrique} et le \cle{champ magnétique} uniformes. Dans cette première section, nous apprenons à accélérer une particule chargée grâce à un champ électrique uniforme, et nous établissons la formule capitale qui donne sa vitesse de sortie : $v = \sqrt{\dfrac{2|q|U}{m}}$.

\courssub{Champ électrique uniforme entre deux armatures planes}

\textbf{Situation concrète.}\par
Prenons deux plaques métalliques planes et parallèles, distantes de $d$, reliées aux bornes d'un générateur de tension continue $U$. C'est un \emph{condensateur plan}. Entre les armatures règne un champ électrique $\vec{E}$ dont les propriétés sont remarquablement simples.

\begin{definition}[Champ électrique uniforme]
Entre deux armatures planes et parallèles, séparées d'une distance $d$ et soumises à une tension $U$, le champ électrique est \cle{uniforme} : le vecteur $\vec{E}$ a même direction, même sens et même norme en tout point de l'espace entre les plaques. Il est dirigé perpendiculairement aux armatures, de la plaque positive vers la plaque négative, et sa norme vaut :
\[ E = \frac{U}{d} \]
avec $E$ en volt par mètre ($\mathrm{V\,m^{-1}}$), $U$ en volt (V) et $d$ en mètre (m).
\end{definition}

\begin{attention}[Ne pas confondre $U$ et $E$]
C'est l'erreur la plus fréquente au Baccalauréat. La \textbf{tension} $U$ se mesure en \textbf{volts} (V) ; le \textbf{champ} $E$ se mesure en \textbf{volts par mètre} ($\mathrm{V\,m^{-1}}$). Ils sont reliés par $E = U/d$ : si tu doubles la distance entre les plaques à tension constante, le champ est divisé par deux. Vérifie toujours l'homogénéité : $U/d$ a bien la dimension d'une tension divisée par une longueur.
\end{attention}

Une particule de charge $q$ placée dans ce champ subit la force électrique :
\[ \vec{F} = q\vec{E} \]
Comme $\vec{E}$ est uniforme et que $q$ est constante, la force $\vec{F}$ est \cle{constante} en norme, direction et sens pendant toute la traversée. C'est précisément cette constance qui rend le problème soluble élégamment : nous avons affaire à un mouvement à force constante, comme la chute libre, mais avec une « gravité » électrique que nous contrôlons à volonté grâce au bouton de tension.

\courssub{Le canon à électrons : description du dispositif}

\begin{definition}[Canon à électrons]
Un \cle{canon à électrons} est un dispositif placé dans une enceinte vide, qui produit un faisceau d'électrons dont la vitesse est contrôlée par une \cle{tension accélératrice} $U$. Il comprend :
\begin{itemize}
\item une \textbf{cathode} (électrode négative), filament chauffé qui émet des électrons par effet thermoélectronique, avec une vitesse initiale quasi nulle ;
\item une \textbf{anode} (électrode positive), plaque percée d'un trou, portée au potentiel $+U$ par rapport à la cathode ;
\item la \textbf{tension accélératrice} $U$ appliquée entre cathode et anode.
\end{itemize}
Les électrons émis par la cathode sont attirés par l'anode, accélérés dans le champ $\vec{E}$ qui règne entre les deux électrodes, puis traversent le trou de l'anode en formant un faisceau fin de vitesse $v$ bien déterminée.
\end{definition}

\begin{popfigure}
\begin{center}
\begin{tikzpicture}[scale=1.0, >=Stealth]
% Cathode
\fill[poppink] (-0.15,-1.4) rectangle (0.15,1.4);
\node[below, poppink!70!black] at (0,-1.5) {\small Cathode (potentiel $0$)};
% Anode percée
\fill[popblue] (3.4,-1.4) rectangle (3.7,-0.18);
\fill[popblue] (3.4,0.18) rectangle (3.7,1.4);
\node[below, popblue!70!black] at (3.55,-1.5) {\small Anode percée (potentiel $+U$)};
% Champ E (vers la gauche, de + vers -)
\foreach \y in {-1.0,-0.5,0.5,1.0} {
  \draw[->, popgreen!70!black, thick] (2.9,\y) -- (0.7,\y);
}
\node[popgreen!70!black] at (1.8,1.25) {$\vec{E}$};
% Trajectoire de l'électron
\draw[->, poporange, very thick] (0.25,0) -- (5.3,0);
\node[poporange!80!black, above] at (4.6,0.05) {\small faisceau d'électrons};
% électron
\fill[popdark] (1.6,0) circle (0.09);
\node[below, popdark] at (1.6,-0.12) {\small $e^-$};
% vitesse
\draw[->, poporange!80!black, very thick] (1.6,0.35) -- (2.6,0.35);
\node[above, poporange!80!black] at (2.1,0.35) {$\vec{v}$};
% tension U
\draw[<->, poppurple, thick] (0,-1.9) -- (3.55,-1.9);
\node[below, poppurple] at (1.78,-1.95) {\small tension accélératrice $U$};
\end{tikzpicture}
\end{center}
\legende{Canon à électrons : la cathode émet des électrons (vitesse initiale négligeable) ; le champ $\vec{E}$, dirigé de l'anode vers la cathode, exerce sur chaque électron une force $\vec{F} = q\vec{E}$ dirigée vers l'anode (car $q = -e < 0$). L'électron ressort par le trou de l'anode avec la vitesse $\vec{v}$.}
\end{popfigure}

\begin{attention}[Sens de la force pour un électron]
La force $\vec{F} = q\vec{E}$ est colinéaire à $\vec{E}$, mais pour un électron $q = -e < 0$ : la force est donc de \textbf{sens opposé} au champ. C'est bien ce qu'il faut : le champ $\vec{E}$ pointe de l'anode ($+$) vers la cathode ($-$), tandis que l'électron, négatif, est attiré vers l'anode ($+$). Retiens : \emph{les charges négatives « remontent » le champ}.
\end{attention}

\courssub{Vitesse de sortie : application du théorème de l'énergie cinétique}

\textbf{Intuition.}\par
Pourquoi utiliser l'énergie plutôt que les lois de Newton ? On pourrait écrire $\vec{F} = m\vec{a}$, trouver une accélération constante, puis intégrer deux fois. Cela fonctionne, mais le théorème de l'énergie cinétique donne le résultat en une seule ligne, sans même connaître la distance $d$ entre les électrodes. C'est un outil puissant : il relie directement la variation d'énergie cinétique au travail des forces.

\begin{propriete}[Vitesse acquise par une particule chargée — démonstration exigible]
Une particule de charge $q$ et de masse $m$, partant avec une vitesse initiale négligeable de la cathode, est accélérée par une tension $U$ jusqu'à l'anode. Sa vitesse à la sortie est :
\[ \boxed{\; v = \sqrt{\frac{2|q|U}{m}} \;} \]
\textbf{Démonstration guidée} (à savoir refaire au Bac) :
\begin{enumerate}
\item \textbf{Bilan des forces.} La particule subit la force électrique $\vec{F} = q\vec{E}$ et son poids $\vec{P} = m\vec{g}$. Nous montrerons plus bas que le poids est négligeable devant la force électrique.
\item \textbf{Théorème de l'énergie cinétique} entre la cathode C et l'anode A :
\[ E_c(A) - E_c(C) = W_{C\to A}(\vec{F}) \]
\item \textbf{Travail de la force électrique.} Pour une charge $q$ se déplaçant d'un point C au potentiel $V_C$ vers un point A au potentiel $V_A$ :
\[ W_{C\to A}(\vec{F}) = q\,(V_C - V_A) = q\,U_{CA} \]
Pour un électron ($q = -e$) allant de la cathode ($V_C = 0$) vers l'anode ($V_A = +U$) : $W = (-e)(0 - U) = eU > 0$. Le travail est bien \emph{moteur}.
\item \textbf{Énergie cinétique initiale nulle} (électrons émis « sans vitesse ») : $E_c(C) \approx 0$.
\item \textbf{Conclusion} :
\[ \frac{1}{2}mv^2 - 0 = |q|U \quad\Longrightarrow\quad v = \sqrt{\frac{2|q|U}{m}} \]
\end{enumerate}
\textbf{Vérification dimensionnelle} : $[qU] = \mathrm{C \times V} = \mathrm{J} = \mathrm{kg\,m^2\,s^{-2}}$, donc $\left[\sqrt{qU/m}\right] = \sqrt{\mathrm{m^2\,s^{-2}}} = \mathrm{m\,s^{-1}}$. La formule est homogène.
\end{propriete}

\begin{attention}[La valeur absolue sur $q$]
Écrire $v = \sqrt{2qU/m}$ sans valeur absolue est une erreur grave : pour un électron, $q = -e < 0$, et la racine carrée d'un nombre négatif n'existe pas dans $\mathbb{R}$ ! Physiquement, c'est $|q|U$ qui représente l'énergie gagnée. Pour une charge positive accélérée, on la fait partir de l'anode vers la cathode : dans tous les cas, le travail moteur vaut $|q|U$. Retiens : \textbf{énergie acquise $= |q|U$, toujours positive}.
\end{attention}

\begin{methode}[Avant tout calcul : vérifier que le poids est négligeable]
On néglige systématiquement le poids devant la force électrique dans ce chapitre, mais il faut savoir le \emph{justifier}, car un correcteur peut le demander. Procède ainsi :
\begin{enumerate}
\item Calcule la norme du poids : $P = mg$.
\item Calcule la norme de la force électrique : $F = |q|E = |q|\dfrac{U}{d}$.
\item Compare les deux (rapport $F/P$). Si $F \gg P$ (typiquement un rapport supérieur à $10^{10}$ pour un électron), le poids est négligeable et tu peux l'écrire explicitement dans ta copie.
\end{enumerate}
\end{methode}

\begin{exemplebox}[Le poids est-il vraiment négligeable ?]
Pour un électron ($m = \num{9,11e-31}$ kg, $|q| = e = \num{1,60e-19}$ C) accéléré entre deux plaques distantes de $d = \SI{5,0}{cm}$ sous $U = \SI{500}{V}$ :
\begin{align*}
P &= mg = \num{9,11e-31} \times \num{9,81} \approx \num{8,9e-30} \ \mathrm{N} \\
F &= e\,\frac{U}{d} = \num{1,60e-19} \times \frac{500}{5,0 \times 10^{-2}} = \num{1,6e-15} \ \mathrm{N}
\end{align*}
Le rapport vaut $\dfrac{F}{P} \approx \num{1,8e14}$ : la force électrique est environ \textbf{cent mille milliards de fois} plus intense que le poids. Négliger $\vec{P}$ est donc parfaitement légitime.
\end{exemplebox}

\courssub{Exemple chiffré et ordres de grandeur}

\begin{exoresolu}[Électron accéléré sous 500 V]
\textbf{Énoncé.} Dans le tube d'un appareil de radiographie, un électron est émis par la cathode avec une vitesse négligeable, puis accéléré par une tension $U = \SI{500}{V}$. Données : $e = \num{1,60e-19}$ C ; $m = \num{9,11e-31}$ kg.
\begin{enumerate}
\item Calculer le travail de la force électrique entre cathode et anode.
\item En déduire la vitesse $v$ de l'électron à la sortie du canon.
\item Calculer son énergie cinétique en joules, puis en électron-volts.
\end{enumerate}
\tcblower
\textbf{Solution détaillée.}
\begin{enumerate}
\item Le travail de la force électrique est moteur et vaut :
\[ W = |q|U = eU = \num{1,60e-19} \times 500 = \num{8,0e-17} \ \mathrm{J} \]
\item Théorème de l'énergie cinétique (vitesse initiale nulle, poids négligeable) :
\[ \frac{1}{2}mv^2 = eU \quad\Longrightarrow\quad v = \sqrt{\frac{2eU}{m}} = \sqrt{\frac{2 \times \num{1,60e-19} \times 500}{\num{9,11e-31}}} \]
\[ v = \sqrt{1,756 \times 10^{14}} \approx \num{1,3e7} \ \mathrm{m\,s^{-1}} \]
Soit environ \SI{13 000}{km/s}, plus de 4 \% de la vitesse de la lumière, atteinte en quelques centimètres !
\item Énergie cinétique :
\[ E_c = \frac{1}{2}mv^2 = eU = \num{8,0e-17} \ \mathrm{J} \]
En électron-volts : un électron accéléré sous 1 volt gagne 1 eV ; sous 500 V, il gagne donc :
\[ E_c = 500 \ \mathrm{eV} \]
La conversion est immédiate, c'est tout l'intérêt de cette unité.
\end{enumerate}
\end{exoresolu}

\begin{popfigure}
\begin{center}
\begin{tikzpicture}
\begin{axis}[
  width=0.85\linewidth, height=6.5cm,
  xlabel={$\sqrt{U}$ (en $\mathrm{V^{1/2}}$)},
  ylabel={$v$ (en $10^{7}\ \mathrm{m\,s^{-1}}$)},
  xmin=0, xmax=32, ymin=0, ymax=2.1,
  grid=major, grid style={popline!40},
  legend pos=north west,
]
\addplot[popcurve, domain=0:32, samples=100] {0.0593*x};
\addlegendentry{$v = \sqrt{\dfrac{2e}{m}}\times\sqrt{U}$}
\addplot[only marks, mark=*, poporange, mark size=2pt] coordinates {(22.36,1.326)};
\node[poporange!80!black, anchor=south west] at (axis cs:22.36,1.326) {\small $U = 500$ V};
\end{axis}
\end{tikzpicture}
\end{center}
\legende{Vitesse d'un électron en fonction de $\sqrt{U}$ : la relation $v = \sqrt{2eU/m}$ est une fonction linéaire de $\sqrt{U}$, de pente $\sqrt{2e/m} \approx \num{5,93e5}$ $\mathrm{m\,s^{-1}\,V^{-1/2}}$. Le point marqué correspond à $U = 500$ V, soit $v \approx \num{1,3e7}$ m/s.}
\end{popfigure}

\textbf{Ordres de grandeur à retenir.}\par
Sous quelques centaines de volts, un électron atteint des vitesses de l'ordre de $10^{7}$ m/s. Sous les dizaines de kilovolts d'un tube à rayons X hospitalier, on approche $10^{8}$ m/s, et la mécanique classique commence à montrer ses limites (il faudrait la relativité). Au programme de Terminale C, on reste dans le domaine où $v \ll c$ et où notre formule s'applique sans restriction.

\courssub{L'électron-volt : l'unité d'énergie des physiciens}

\begin{aretenir}[Énergie acquise et électron-volt]
L'énergie cinétique acquise par une particule de charge $|q|$ accélérée sous une tension $U$ est :
\[ E_c = |q|\,U \]
Pour toute particule de charge $|q| = e$ (électron, proton, ion simplement chargé) accélérée sous une tension $U$ exprimée en volts, l'énergie acquise en électron-volts est \emph{numériquement égale} à $U$ :
\[ U = 500\ \mathrm{V} \quad\Longrightarrow\quad E_c = 500\ \mathrm{eV} \]
Conversion : $1\ \mathrm{eV} = \num{1,60e-19}$ J.
\end{aretenir}

\begin{savaistu}[L'électron-volt, monnaie courante de la physique moderne]
Le joule est une unité adaptée à notre échelle (une pomme qui tombe), mais démesurée pour un électron : $10^{-17}$ J, ce n'est pas parlant. Les physiciens utilisent donc l'\cle{électron-volt} : $1\ \mathrm{eV} = \num{1,60e-19}$ J, énergie gagnée par un électron accéléré sous 1 volt. Ses multiples rythment la physique moderne : le keV (rayons X des hôpitaux de Douala ou Yaoundé, quelques dizaines de keV), le MeV (radiothérapie des cancers, accélérateurs de particules), le GeV (grand collisionneur du CERN). Au Cameroun, les appareils de radiothérapie utilisés en oncologie accélèrent des électrons jusqu'à plusieurs MeV : c'est exactement le principe du canon à électrons que tu viens d'étudier, poussé à l'extrême.
\end{savaistu}

\begin{exercice}[Pour t'entraîner]
Un proton ($q = +e = \num{1,60e-19}$ C, $m_p = \num{1,67e-27}$ kg), initialement au repos, est accéléré sous une tension $U = \SI{2,0}{kV}$.
\begin{enumerate}
\item Justifier que le poids du proton est négligeable devant la force électrique si les plaques sont distantes de $d = \SI{10}{cm}$.
\item Calculer sa vitesse finale et son énergie cinétique en eV.
\item Comparer avec la vitesse qu'atteindrait un électron sous la même tension. Expliquer l'écart.
\end{enumerate}
\end{exercice}

\begin{corrige}[Exercice]
\begin{enumerate}
\item $P = m_p g = \num{1,67e-27} \times \num{9,81} \approx \num{1,6e-26}$ N ; $F = eU/d = \num{1,60e-19} \times 2000/0{,}10 = \num{3,2e-15}$ N. Rapport $F/P \approx \num{2e11} \gg 1$ : le poids est négligeable.
\item $\dfrac{1}{2}m_p v^2 = eU$ donne $v = \sqrt{\dfrac{2 \times \num{1,60e-19} \times 2000}{\num{1,67e-27}}} \approx \num{6,2e5}$ m/s. Énergie : $E_c = eU = \num{2,0}{keV} = \num{3,2e-16}$ J.
\item Pour l'électron : $v_e = \sqrt{2eU/m_e} \approx \num{2,7e7}$ m/s. À énergie égale ($eU$ identique), la vitesse varie en $1/\sqrt{m}$ : le proton, 1836 fois plus lourd, va $\sqrt{1836} \approx 43$ fois moins vite. C'est pourquoi les canons à particules légères (électrons) produisent des vitesses bien plus grandes.
\end{enumerate}
\end{corrige}

\begin{aretenir}[L'essentiel de la section]
\begin{itemize}
\item Entre deux armatures planes sous tension $U$ : champ uniforme $E = U/d$, force constante $\vec{F} = q\vec{E}$.
\item Le canon à électrons convertit l'énergie électrique $|q|U$ en énergie cinétique : $\dfrac{1}{2}mv^2 = |q|U$, d'où $v = \sqrt{\dfrac{2|q|U}{m}}$.
\item Le poids est négligeable devant la force électrique (à justifier par comparaison numérique).
\item L'électron-volt : $1\ \mathrm{eV} = \num{1,60e-19}$ J ; une charge $e$ sous $U$ volts gagne $U$ eV.
\item Pièges : valeur absolue de $q$, distinction $U$ (V) / $E$ (V/m), vitesse initiale supposée nulle.
\end{itemize}
Maintenant que nos électrons sont lancés à grande vitesse, la question suivante se pose naturellement : comment les \emph{dévier} pour dessiner une image sur un écran ? C'est l'objet de la section suivante, où la trajectoire rectiligne laissera place à une belle parabole.
\end{aretenir}

\coursec{Déviation électrique : trajectoire parabolique}

Dans la section précédente, nous avons vu qu'un champ électrique uniforme \emph{colinéaire} à la vitesse initiale d'une particule chargée l'accélère ou la ralentit en ligne droite : c'est le principe du canon à électrons, et le théorème de l'énergie cinétique nous a donné la vitesse acquise. Mais que se passe-t-il si la particule arrive avec une vitesse $\vec{v}_0$ \emph{perpendiculaire} au champ électrique ? Le champ ne peut plus la « pousser » dans sa direction de marche : il la \cle{dévie} latéralement. C'est exactement la situation d'une balle lancée horizontalement dans le champ de pesanteur : la trajectoire se courbe. Nous allons montrer rigoureusement que, comme pour la chute libre horizontale, la trajectoire est une \cle{parabole}.

\courssub{Situation d'étude et mise en place du problème}

\begin{experience}[Le condensateur plan déviateur]
\textbf{Dispositif.} Un faisceau de particules chargées (produit par exemple par un canon à électrons) pénètre en $O$, avec une vitesse $\vec{v}_0$ horizontale, entre les deux armatures planes et parallèles d'un condensateur. Les plaques, de longueur $L$, sont distantes de $d$ et soumises à une tension $U = V_A - V_B$. Le champ électrique $\vec{E}$ régnant entre les armatures est uniforme, vertical. On suppose que la plaque du bas ($V_A$) est positive et la plaque du haut ($V_B$) négative ; le champ est donc dirigé \textbf{vers le haut} : $\vec{E} = E\,\vec{\jmath}$ avec $E = U/d > 0$.
\textbf{Observation.} À la sortie des plaques, le faisceau frappe un écran fluorescent en un point décalé par rapport à l'axe initial : la particule a été \emph{déviée}. En inversant la polarité de la tension, la déviation s'inverse. En augmentant $U$, la déviation augmente.
\textbf{Question.} Quelle est la nature mathématique de la trajectoire, et de quels paramètres dépend la déviation ?
\end{experience}

\begin{popfigure}
\begin{tikzpicture}[scale=1.05]
% Armatures : BAS = positive (+), HAUT = negative (-)
\fill[popblue!25] (0,-1.4) rectangle (5.5,-1.7); % Bas
\fill[poporange!25] (0,1.4) rectangle (5.5,1.7);  % Haut
\node[popblue] at (5.1,-1.55) {\small $+$};
\node[popblue] at (4.6,-1.55) {\small $+$};
\node[popblue] at (4.1,-1.55) {\small $+$};
\node[poporange] at (5.1,1.55) {\small $-$};
\node[poporange] at (4.6,1.55) {\small $-$};
\node[poporange] at (4.1,1.55) {\small $-$};
\node[below, popblue] at (2.75,-1.7) {\small plaque positive ($V_A$)};
\node[above, poporange] at (2.75,1.7) {\small plaque négative ($V_B$)};
% Repère : Oy vers le HAUT
\draw[->, thick] (0,0) -- (0.9,0) node[below left] {\small $x$};
\draw[->, thick] (0,0) -- (0,0.9) node[left] {\small $y$};
\fill (0,0) circle (1.5pt) node[below left] {\small $O$};
% Champ E : vers le haut (de + vers -)
\foreach \xx in {0.7,1.9,3.1,4.3} {\draw[->, popgreen, thick] (\xx,-0.9) -- (\xx,-0.25);}
\node[popgreen] at (4.3,-1.1) {\small $\vec{E}$};
% Trajectoire parabolique (electron : q<0, force vers le bas, y<0)
\draw[very thick, poppurple, domain=-1.2:0, samples=30] plot (\x,0);
\draw[very thick, poppurple, domain=0:5.2, samples=60] plot (\x,{-0.045*\x*\x});
% Vecteur v0
\draw[->, very thick, popink] (-1.2,0) -- (-0.2,0) node[midway, above] {\small $\vec{v}_0$};
% Force sur la trajectoire (vers le bas)
\draw[->, very thick, poporange] (3.5,-0.55) -- (3.5,-1.25) node[right] {\small $\vec{F} = q\vec{E}$};
% Sortie
\draw[dashed, popmuted] (5.2,-1.4) -- (5.2,1.4);
\node[popmuted, rotate=90, anchor=south] at (5.35,0) {\small sortie $x = L$};
% Cotes
\draw[<->, popmuted] (6.0,-1.4) -- (6.0,1.4) node[midway, right] {\small $d$};
\end{tikzpicture}
\legende{Particule chargée (électron, $q<0$) entrant avec $\vec{v}_0$ horizontal entre les armatures. La plaque du bas est positive ($V_A$), celle du haut négative ($V_B$) : $\vec{E}$ est vertical vers le haut. La force $\vec{F}=q\vec{E}$ est vers le bas, la trajectoire est une parabole à concavité vers le bas.}
\end{popfigure}

\courssub{Choix du repère et hypothèses}

On choisit un repère \textbf{fixe} $(O\,;\,\vec{\imath},\,\vec{\jmath})$ où :
\begin{itemize}
\item $O$ est le point d'entrée de la particule entre les plaques (milieu de l'entrée) ;
\item l'axe $Ox$ est horizontal, dirigé dans le sens de $\vec{v}_0$ ;
\item l'axe $Oy$ est vertical, dirigé \textbf{vers le haut}.
\end{itemize}

Les conditions initiales sont donc :
\[ \overrightarrow{OM}_0 = \vec{0} \qquad \text{et} \qquad \vec{v}_0 = v_0\,\vec{\imath} \quad (v_{0x} = v_0,\ v_{0y} = 0). \]

\textbf{Forces appliquées à la particule :}
\begin{itemize}
\item le \cle{poids} $\vec{P} = m\vec{g}$ (vertical, vers le bas) ;
\item la \cle{force électrique} $\vec{F}_e = q\vec{E}$.
\end{itemize}

\begin{attention}[Pourquoi négliger le poids ?]
Pour un électron ($m = \num{9,11}\times 10^{-31}$ kg) dans un champ $E = \num{5,0}\times 10^{3}$ V/m :
\[ \frac{F_e}{P} = \frac{|q|E}{mg} = \frac{\num{1,60}\times 10^{-19}\times \num{5,0}\times 10^{3}}{\num{9,11}\times 10^{-31}\times 9,8} \approx \num{9}\times 10^{13}. \]
La force électrique est environ $10^{14}$ fois plus intense que le poids ! Le poids est donc \textbf{toujours négligé} devant la force électrique pour les particules élémentaires (électrons, protons, ions). Mais attention : cette hypothèse doit être \emph{justifiée} dans ta copie, au moins par une phrase ou un ordre de grandeur.
\end{attention}

\courssub{Bilan des forces et deuxième loi de Newton}

Le champ $\vec{E}$ étant uniforme et la particule évoluant dans le vide, la force résultante est $\vec{F}_e = q\vec{E}$, constante. La deuxième loi de Newton s'écrit :
\[ q\vec{E} = m\vec{a} \qquad \Longrightarrow \qquad \boxed{\vec{a} = \frac{q\vec{E}}{m}}. \]

L'accélération est donc \textbf{constante} (en norme, direction et sens) pendant toute la traversée des plaques. Notons $E_y$ la composante algébrique de $\vec{E}$ sur $Oy$ ($\vec{E} = E_y\,\vec{\jmath}$). La projection sur les axes donne :
\[ \begin{cases} a_x = 0 \\[4pt] a_y = \dfrac{q E_y}{m} \end{cases} \]

\begin{aretenir}[Idée physique essentielle]
Le mouvement se décompose en deux mouvements indépendants :
\begin{itemize}
\item selon $Ox$ : \textbf{mouvement uniforme} (aucune force horizontale, $a_x=0$) ;
\item selon $Oy$ : \textbf{mouvement uniformément accéléré} (force électrique constante, $a_y = qE_y/m$ constant).
\end{itemize}
C'est exactement la structure du mouvement d'un projectile dans le champ de pesanteur, avec $qE_y/m$ jouant le rôle de l'accélération de la pesanteur $g$ (orientée vers le bas, donc $g_y = -g$).
\end{aretenir}

\courssub{Équations horaires du mouvement}

\begin{methode}[Établir les équations horaires et la trajectoire — la démarche à maîtriser]
\begin{enumerate}
\item \textbf{Poser le repère} $(O\,;\,x,\,y)$ adapté : origine au point d'entrée, $Ox$ selon $\vec{v}_0$, $Oy$ vertical vers le haut.
\item \textbf{Faire le bilan des forces} et justifier que le poids est négligeable devant $q\vec{E}$.
\item \textbf{Projeter la deuxième loi de Newton} sur chaque axe : $a_x = 0$, $a_y = qE_y/m$.
\item \textbf{Intégrer une première fois} pour obtenir les composantes de la vitesse, en utilisant $v_x(0) = v_0$ et $v_y(0) = 0$.
\item \textbf{Intégrer une seconde fois} pour obtenir $x(t)$ et $y(t)$, en utilisant $x(0) = y(0) = 0$.
\item \textbf{Éliminer le temps} : exprimer $t = x/v_0$ et reporter dans $y(t)$ pour obtenir $y(x)$.
\end{enumerate}
\end{methode}

Appliquons cette méthode. \textbf{Première intégration} (composantes de la vitesse) :
\[ v_x(t) = v_0 \qquad ; \qquad v_y(t) = \frac{q E_y}{m}\,t. \]

\textbf{Seconde intégration} (coordonnées de la position) :
\[ \boxed{x(t) = v_0\,t} \qquad ; \qquad \boxed{y(t) = \frac{q E_y}{2m}\,t^2}. \]

Vérifions l'homogénéité de $y(t)$ : $[qE_y/m] = \dfrac{\text{C}\times\text{V}\cdot\text{m}^{-1}}{\text{kg}} = \dfrac{\text{J}\cdot\text{m}^{-1}}{\text{kg}} = \dfrac{\text{N}}{\text{kg}} = \text{m}\cdot\text{s}^{-2}$, donc $[y] = \text{m}\cdot\text{s}^{-2}\times\text{s}^{2} = \text{m}$. La relation est bien homogène.

\courssub{Équation de la trajectoire : une parabole}

\begin{propriete}[Trajectoire parabolique — démonstration exigible]
Une particule de charge $q$ et de masse $m$, pénétrant avec une vitesse $\vec{v}_0$ perpendiculaire à un champ électrique uniforme $\vec{E} = E_y\,\vec{\jmath}$, décrit entre les plaques une \textbf{parabole} d'équation :
\[ y = \frac{q E_y}{2mv_0^2}\,x^2. \]
\textbf{Démonstration.} De l'équation horaire $x(t) = v_0 t$, on tire $t = \dfrac{x}{v_0}$. En reportant dans $y(t) = \dfrac{q E_y}{2m}t^2$ :
\[ y = \frac{q E_y}{2m}\left(\frac{x}{v_0}\right)^2 = \frac{q E_y}{2mv_0^2}\,x^2. \]
C'est une équation de la forme $y = Kx^2$ avec $K = \dfrac{q E_y}{2mv_0^2}$ constante : la trajectoire est bien une \cle{parabole} de sommet $O$, d'axe $Oy$. $\blacksquare$
\end{propriete}

\begin{attention}[Le signe de $q$ et la direction de $\vec{E}$ déterminent le sens de la déviation]
Le coefficient $K = \dfrac{q E_y}{2mv_0^2}$ porte le signe du produit $q E_y$ :
\begin{itemize}
\item Si $q E_y > 0$ (ex: proton $q>0$ dans un champ vers le haut $E_y>0$, ou électron $q<0$ dans un champ vers le bas $E_y<0$) : $y > 0$, la particule est déviée \textbf{vers le haut} (dans le sens de $\vec{E}$ si $q>0$, opposé si $q<0$).
\item Si $q E_y < 0$ (ex: électron $q<0$ dans un champ vers le haut $E_y>0$, comme sur notre figure) : $y < 0$, la particule est déviée \textbf{vers le bas}.
\end{itemize}
Erreur fréquente au Bac : oublier le signe de $q$ dans l'équation, ou affirmer qu'un électron est attiré par la plaque négative. Retiens : la force $\vec{F} = q\vec{E}$ est \emph{colinéaire et de même sens} que $\vec{E}$ pour $q>0$, \emph{colinéaire et de sens opposé} pour $q<0$. En pratique, on travaille souvent avec les valeurs absolues pour les grandeurs et on précise le sens de la déviation sur le schéma.
\end{attention}

\begin{savaistu}[Des tubes cathodiques aux oscilloscopes]
C'est exactement ce phénomène qui déviait le faisceau d'électrons dans les \textbf{téléviseurs à tube cathodique} (CRT), encore très répandus au Cameroun il y a une quinzaine d'années : deux paires de plaques (l'une horizontale, l'autre verticale) alimentées par des tensions variables balayaient l'écran à grande vitesse pour « dessiner » l'image ligne par ligne. Le même principe équipe l'\textbf{oscilloscope} que tu utilises en TP : les plaques $Y$ (verticales) reçoivent la tension à étudier, les plaques $X$ la tension de balayage. C'est aussi par déviation électrostatique que J.J. Thomson, en 1897, a mesuré le rapport $e/m$ de l'électron et prouvé qu'il s'agissait d'une particule.
\end{savaistu}

\courssub{Condition de sortie des plaques}

La particule ne peut exploiter la déviation que si elle \emph{ressort} du condensateur, c'est-à-dire si elle n'a pas heurté l'armature vers laquelle elle est attirée. L'origine $O$ étant au centre de l'entrée, les plaques se trouvent aux ordonnées $y = +d/2$ (haut) et $y = -d/2$ (bas). À la sortie, $x = L$, et l'ordonnée de sortie vaut :
\[ y_S = y(L) = \frac{q E_y}{2mv_0^2}\,L^2. \]

\begin{propriete}[Condition de non-collision avec l'armature]
La particule sort des plaques sans toucher l'armature si et seulement si :
\[ |y_S| < \frac{d}{2} \qquad \Longleftrightarrow \qquad \frac{|q|\,|E_y|\,L^2}{2mv_0^2} < \frac{d}{2}. \]
Si cette condition n'est pas respectée, la particule est captée par la plaque : il n'y a ni sortie, ni déflexion observable sur l'écran.
\end{propriete}

Cette condition montre les leviers du physicien : pour augmenter la déviation, on peut augmenter $U$ (donc $|E_y|$), allonger les plaques ($L$), ou diminuer la vitesse d'injection $v_0$ — c'est pourquoi les oscilloscopes utilisent des électrons pas trop rapides.

\begin{exoresolu}[Déviation d'un électron entre deux plaques]
\textbf{Énoncé.} Un électron (charge $q = -e = \num{-1,60}\times 10^{-19}$ C, masse $m = \num{9,11}\times 10^{-31}$ kg) pénètre en $O$ avec la vitesse $v_0 = \num{2,0}\times 10^{7}$ m/s entre les armatures horizontales d'un condensateur plan de longueur $L = \SI{5,0}{cm}$, distantes de $d = \SI{2,0}{cm}$, soumises à la tension $U = \SI{100}{V}$. La plaque du bas est positive ($V_A$), la plaque du haut est négative ($V_B$).
\begin{enumerate}
\item Déterminer le champ $\vec{E}$ et l'accélération $\vec{a}$ de l'électron.
\item Établir l'équation de la trajectoire $y(x)$.
\item Calculer l'ordonnée de sortie $y_S$ et vérifier que l'électron sort des plaques.
\item Calculer la durée de traversée et la composante verticale de la vitesse à la sortie. En déduire l'angle de déviation à la sortie.
\end{enumerate}
\tcblower
\textbf{Solution.}
\textbf{1.} Champ électrique : $\vec{E}$ est dirigé de la plaque positive (bas) vers la plaque négative (haut), donc \textbf{vers le haut}. Dans notre repère ($Oy$ vers le haut), $\vec{E} = E\,\vec{\jmath}$ avec
\[ E = \frac{U}{d} = \frac{100}{\num{2,0}\times 10^{-2}} = \num{5,0}\times 10^{3}\ \text{V/m}. \]
Accélération (vectorielle) :
\[ \vec{a} = \frac{q\vec{E}}{m} = \frac{-e\,E\,\vec{\jmath}}{m} = -\frac{eE}{m}\,\vec{\jmath}. \]
Sa norme vaut :
\[ a = \frac{eE}{m} = \frac{\num{1,60}\times 10^{-19}\times \num{5,0}\times 10^{3}}{\num{9,11}\times 10^{-31}} \approx \num{8,8}\times 10^{14}\ \text{m/s}^2. \]
L'électron, de charge négative, subit une force \textbf{opposée} à $\vec{E}$ : il est dévié \textbf{vers le bas} (vers la plaque positive du bas).

\textbf{2.} Composante $y$ de l'accélération : $a_y = -\num{8,8}\times 10^{14}\ \text{m/s}^2$.
L'équation de la trajectoire est :
\[ y = \frac{a_y}{2v_0^2}\,x^2 = -\frac{\num{8,8}\times 10^{14}}{2\times (\num{2,0}\times 10^{7})^{2}}\,x^2 \approx -\num{1,1}\,x^2 \quad (\text{en mètres, } x \text{ en mètres}). \]
C'est une parabole à concavité vers le bas ($y \le 0$).

\textbf{3.} Ordonnée de sortie à $x = L = \num{5,0}\times 10^{-2}$ m :
\[ y_S = -\num{1,1}\times (\num{5,0}\times 10^{-2})^2 = -\num{2,75}\times 10^{-3}\ \text{m} \approx \SI{-2,7}{mm}. \]
Condition de sortie : $|y_S| = \SI{2,7}{mm} < \dfrac{d}{2} = \SI{10}{mm}$. L'électron \textbf{sort bien} des plaques (il ne touche pas la plaque du bas).

\textbf{4.} Durée de traversée :
\[ t_S = \frac{L}{v_0} = \frac{\num{5,0}\times 10^{-2}}{\num{2,0}\times 10^{7}} = \num{2,5}\times 10^{-9}\ \text{s} = \SI{2,5}{ns}. \]
Composante verticale de la vitesse à la sortie :
\[ v_y(t_S) = a_y\,t_S = -\num{8,8}\times 10^{14}\times \num{2,5}\times 10^{-9} \approx -\num{2,2}\times 10^{6}\ \text{m/s}. \]
L'angle de déviation $\alpha$ (orienté, négatif vers le bas) vérifie :
\[ \tan\alpha = \frac{v_y}{v_x} = \frac{-\num{2,2}\times 10^{6}}{\num{2,0}\times 10^{7}} = -0{,}11 \quad \Rightarrow \quad \alpha \approx -6{,}3^\circ. \]
La déviation reste modérée (environ $11\,\%$ de la vitesse initiale en composante verticale).
\end{exoresolu}

\begin{exemplebox}[Influence de la tension sur la trajectoire]
Reprenons les données précédentes et traçons $y(x)$ pour deux tensions, $U_1 = \SI{50}{V}$ et $U_2 = \SI{100}{V}$. Comme $y \propto E_y \propto U$, doubler la tension double l'ordonnée de sortie : $y_{S,1} \approx \SI{-1,4}{mm}$ et $y_{S,2} \approx \SI{-2,7}{mm}$. La déviation est donc un excellent moyen de \emph{mesurer} une tension : c'est le cœur du fonctionnement de l'oscilloscope.
\end{exemplebox}

\begin{popfigure}
\begin{tikzpicture}
\begin{axis}[
width=0.85\linewidth, height=6cm,
xlabel={$x$ (cm)}, ylabel={$y$ (mm)},
xmin=0, xmax=5, ymin=-3.2, ymax=0,
grid=both, grid style={popline!40},
legend pos=south west,
legend style={font=\small},
samples=100, domain=0:5]
\addplot[popcurve, popblue] {-0.55*x*x/25*2.5};
\addlegendentry{$U_1 = 50$ V}
\addplot[popcurve, poporange] {-1.1*x*x/25*2.5};
\addlegendentry{$U_2 = 100$ V}
\draw[dashed, popmuted] (axis cs:5,-3.2) -- (axis cs:5,0);
\node[popmuted, rotate=90, anchor=north east, font=\small] at (axis cs:5,-3.0) {sortie $x=L$};
\end{axis}
\end{tikzpicture}
\legende{Trajectoires paraboliques $y(x)$ de l'électron entre les plaques pour deux tensions ($L = \SI{5}{cm}$, $d = \SI{2}{cm}$, $v_0 = \num{2,0}\times 10^{7}$ m/s, plaque du bas positive). L'ordonnée de sortie (négative) est proportionnelle à $U$.}
\end{popfigure}

\courssub{Analogie avec la chute libre horizontale}

Tu as déjà étudié en Première le mouvement d'un projectile lancé horizontalement dans le champ de pesanteur (repère $Oy$ vers le haut) : $\vec{a} = -g\,\vec{\jmath}$, $x(t) = v_0 t$, $y(t) = -\frac{1}{2}gt^2$, trajectoire $y = -\dfrac{g}{2v_0^2}x^2$. La correspondance est totale :

\begin{center}
\begin{tabularx}{\linewidth}{|l|X|X|}
\hline
\rowcolor{popblueL} \textbf{Grandeur} & \textbf{Projectile (pesanteur)} & \textbf{Particule chargée} \\
\hline
Force & $\vec{P} = m\vec{g} = -mg\,\vec{\jmath}$ & $\vec{F}_e = q\vec{E} = qE_y\,\vec{\jmath}$ \\
\hline
Accélération & $\vec{g} = -g\,\vec{\jmath}$ (constante) & $\vec{a} = \dfrac{qE_y}{m}\,\vec{\jmath}$ (constante) \\
\hline
Trajectoire & $y = -\dfrac{g}{2v_0^2}x^2$ & $y = \dfrac{qE_y}{2mv_0^2}x^2$ \\
\hline
Nature & parabole (concavité vers le bas) & parabole (concavité selon signe de $qE_y$) \\
\hline
\end{tabularx}
\end{center}

\begin{aretenir}[L'essentiel de la section]
\begin{itemize}
\item Entre les plaques, l'accélération $\vec{a} = q\vec{E}/m$ est \textbf{constante} : mouvement uniforme selon $Ox$, uniformément accéléré selon $Oy$.
\item Équations horaires : $x(t) = v_0 t$ et $y(t) = \dfrac{qE_y}{2m}t^2$.
\item Trajectoire : \textbf{parabole} $y = \dfrac{qE_y}{2mv_0^2}x^2$.
\item Le \textbf{signe de $qE_y$} fixe le sens de la déviation (concavité de la parabole) ; l'électron ($q<0$) dans un champ vers le haut ($E_y>0$) est dévié vers le bas (vers la plaque \textbf{positive}).
\item Condition de sortie : $|y(L)| < d/2$.
\item À la sortie ($x=L$), la vitesse fait un angle $\alpha$ tel que $\tan\alpha = \dfrac{qE_y L}{mv_0^2}$.
\end{itemize}
\end{aretenir}

À la sortie des plaques, la particule n'est plus soumise à aucune force (champ nul à l'extérieur) : elle poursuit en ligne droite, tangente à la parabole, jusqu'à l'écran où l'on mesure la déflexion. C'est précisément l'objet de la section suivante.

\coursec{Déflexion sur un écran : l'oscilloscope}

Dans la section précédente, nous avons établi qu'une particule chargée pénétrant avec une vitesse $\vec{v}_0$ perpendiculaire à un champ électrique uniforme $\vec{E}$ décrit, entre les plaques du condensateur, une trajectoire \cle{parabolique} d'équation :
\[ y = \frac{qE}{2 m v_0^2}\, x^2. \]
Les équations horaires associées (rappel de la section 2) sont $x(t) = v_0 t$ et $y(t) = \frac{1}{2}\frac{qE}{m}t^2$. Mais à quoi sert cette courbure ? La réponse tient en un mot : l'\cle{oscilloscope}. Cet appareil, présent dans tous les laboratoires de physique des lycées du Cameroun et dans tous les ateliers de réparation d'appareils électroniques de Douala ou Yaoundé, exploite la déviation électrique d'un faisceau d'électrons pour \emph{rendre visible} une tension électrique. Dans cette section, nous allons suivre l'électron après sa sortie des plaques, montrer qu'il frappe un écran fluorescent en un point dont la position est \textbf{proportionnelle à la tension appliquée}, et en déduire le principe de la mesure des tensions.

\courssub{Le tube à déviation électrostatique : description}

Un oscilloscope classique (à tube cathodique) contient un \cle{tube à déviation électrostatique} dans lequel règne un vide poussé. On y distingue trois parties :

\begin{itemize}
\item Le \cle{canon à électrons} : un filament chauffé émet des électrons (effet thermoélectronique) qui sont accélérés par une tension $U_{\text{acc}}$ de quelques centaines à quelques milliers de volts. À la sortie du canon, chaque électron possède une vitesse $v_0$ horizontale telle que (théorème de l'énergie cinétique, section 1) :
\[ \frac{1}{2} m v_0^2 = |q|\, U_{\text{acc}} \quad\Longrightarrow\quad v_0 = \sqrt{\frac{2 |q| U_{\text{acc}}}{m}}. \]
\item Les \cle{plaques de déviation} : deux paires de plaques planes, perpendiculaires entre elles. Les plaques verticales $Y_1Y_2$ créent un champ $\vec{E}$ vertical qui dévie le faisceau verticalement (déviation $Y$) ; les plaques $X_1X_2$ créent un champ horizontal qui dévie le faisceau horizontalement (déviation $X$). En appliquant des tensions variables sur ces plaques, on dirige le spot à volonté sur l'écran.
\item L'\cle{écran fluorescent} : une couche de substance luminescente déposée sur la face interne du tube. Là où les électrons frappent l'écran, un point lumineux (le \cle{spot}) apparaît.
\end{itemize}

\begin{popfigure}
\begin{center}
\begin{tikzpicture}[scale=1,>=Stealth]
% blocs
\draw[fill=popblueL,draw=popblue] (0,0) rectangle (1.6,1.4);
\node[align=center,font=\small] at (0.8,0.7) {Canon à\\électrons};
\draw[fill=poporangeL,draw=poporange] (2.2,0) rectangle (3.8,1.4);
\node[align=center,font=\small] at (3,0.7) {Plaques\\$Y_1Y_2$};
\draw[fill=popgreenL,draw=popgreen] (4.4,0) rectangle (6,1.4);
\node[align=center,font=\small] at (5.2,0.7) {Plaques\\$X_1X_2$};
\draw[fill=poppurpleL,draw=poppurple] (6.6,0) rectangle (8.2,1.4);
\node[align=center,font=\small] at (7.4,0.7) {Écran\\fluorescent};
% flèches
\draw[->,thick,popdark] (1.6,0.7) -- (2.2,0.7);
\draw[->,thick,popdark] (3.8,0.7) -- (4.4,0.7);
\draw[->,thick,popdark] (6,0.7) -- (6.6,0.7);
% faisceau
\draw[->,very thick,popink] (0.8,1.5) -- (7.4,1.5) node[midway,above,font=\small]{faisceau d'électrons};
% tension
\node[font=\small,popmuted,align=center] at (0.8,-0.5) {$U_{\text{acc}}$};
\node[font=\small,popmuted,align=center] at (3,-0.5) {$U_Y$};
\node[font=\small,popmuted,align=center] at (5.2,-0.5) {$U_X$};
\node[font=\small,popmuted,align=center] at (7.4,-0.5) {spot};
\end{tikzpicture}
\end{center}
\legende{Schéma simplifié d'un oscilloscope à tube cathodique : les électrons accélérés par le canon traversent deux paires de plaques de déviation avant de frapper l'écran fluorescent.}
\end{popfigure}

\begin{savaistu}[L'oscilloscope, l'œil de l'électronicien]
L'oscilloscope permet de \emph{voir} des tensions variables : tension du secteur ($\SI{220}{V}$, $\SI{50}{Hz}$ distribués par ENEO), signaux audio d'un microphone, signaux cardiaques en électrocardiographie dans les hôpitaux... Au Cameroun, tout technicien en électronique qui répare téléviseurs, amplificateurs ou onduleurs utilise un oscilloscope pour diagnostiquer les pannes. La grandeur affichée verticalement est une tension ; l'axe horizontal est le temps, grâce à une tension en « dents de scie » appliquée aux plaques $X$ (la \emph{base de temps}). C'est l'aboutissement technologique direct de la physique que tu étudies dans cette leçon !
\end{savaistu}

\courssub{Mouvement après la sortie des plaques}

À la sortie des plaques (point $S$ d'abscisse $x = L$), l'électron quitte la zone où règne le champ $\vec{E}$. Que devient son mouvement ?

\begin{propriete}[Mouvement après les plaques]
Après la sortie des plaques, la particule n'est plus soumise à aucune force (poids négligeable, champ nul). D'après la première loi de Newton (principe d'inertie), son mouvement est \cle{rectiligne uniforme}, suivant la \textbf{tangente à la parabole au point de sortie} $S$, avec la vitesse $\vec{v}_S$ acquise en $S$.
\end{propriete}

C'est une conséquence directe du principe d'inertie : $\sum \vec{F} = \vec{0} \Rightarrow \vec{v} = \overrightarrow{\text{cste}}$. Le faisceau se dirige donc en ligne droite vers l'écran, placé à la distance $D$ du milieu des plaques.

\courssub{Propriété géométrique : la tangente passe par le milieu des plaques}

Avant de calculer la déflexion sur l'écran, établissons un résultat géométrique remarquable qui simplifie tous les calculs.

\begin{experience}[Démonstration guidée : la tangente en $S$ passe par le point $I$, milieu des plaques]
\textbf{Étape 1 — Coordonnées du point de sortie $S$.} La trajectoire entre les plaques est $y = \dfrac{qE}{2 m v_0^2} x^2$. Posons $k = \dfrac{qE}{2 m v_0^2}$, de sorte que $y = kx^2$. En $x = L$ :
\[ y_S = k L^2 = \frac{q E L^2}{2 m v_0^2}. \]
\textbf{Étape 2 — Pente de la tangente en $S$.} La pente de la tangente est la dérivée :
\[ \tan\alpha = \left(\frac{dy}{dx}\right)_{x=L} = 2kL = \frac{qEL}{m v_0^2}. \]
\textbf{Étape 3 — Équation de la tangente en $S$.} C'est la droite passant par $S(L\,;\,kL^2)$ de pente $2kL$ :
\[ y - kL^2 = 2kL\,(x - L) \quad\Longrightarrow\quad y = 2kL\,x - kL^2. \]
\textbf{Étape 4 — Intersection avec l'axe $Ox$.} Cherchons le point $I$ où cette tangente coupe l'axe des $x$ ($y = 0$) :
\[ 0 = 2kL\,x_I - kL^2 \quad\Longrightarrow\quad x_I = \frac{L}{2}. \]
\textbf{Conclusion.} La tangente à la trajectoire au point de sortie $S$ coupe l'axe $Ox$ au point $I$ d'abscisse $L/2$, c'est-à-dire au \textbf{milieu des plaques}. Tout se passe donc comme si, après les plaques, l'électron provenait en ligne droite du point $I$.
\end{experience}

\begin{aretenir}
La déviation angulaire $\alpha$ du faisceau à la sortie des plaques vérifie :
\[ \tan\alpha = \frac{q E L}{m v_0^2} = \frac{q U L}{m v_0^2\, d}, \]
où $U$ est la tension appliquée entre les plaques, $d$ leur écartement, $L$ leur longueur. Le prolongement de la trajectoire rectiligne après les plaques passe par le milieu $I$ des plaques.
\end{aretenir}

\begin{popfigure}
\begin{center}
\begin{tikzpicture}[scale=1.05,>=Stealth]
% axe
\draw[->,popmuted] (-0.5,0) -- (10.5,0) node[below left]{$x$};
% plaques
\draw[very thick,popblue] (1,0.9) -- (4,0.9) node[pos=0,above left,font=\small]{$Y_1$};
\draw[very thick,popblue] (1,-0.9) -- (4,-0.9) node[pos=0,below left,font=\small]{$Y_2$};
% champ E
\foreach \x in {1.4,2.2,3.0,3.8}{\draw[->,popgreen] (\x,0.6) -- (\x,-0.6);}
\node[popgreen,font=\small] at (3.4,0.45) {$\vec{E}$};
% cotes L et d
\draw[<->,popmuted] (1,-1.25) -- (4,-1.25) node[midway,below,font=\small]{$L$};
\draw[<->,popmuted] (0.7,-0.9) -- (0.7,0.9) node[midway,left,font=\small]{$d$};
% trajectoire parabolique (schéma qualitatif)
\draw[very thick,popink,domain=0:3,samples=60] plot (\x,{0.04*\x*\x});
% point S
\fill[popink] (3,0.36) circle (1.5pt) node[above right,font=\small]{$S$};
% tangente prolongée jusqu'à l'écran
\draw[very thick,poporange] (3,0.36) -- (9,2.52);
% prolongement pointillé vers I
\draw[dashed,poporange] (3,0.36) -- (2.5,0);
\fill[poporange] (2.5,0) circle (1.5pt) node[below,font=\small]{$I$};
% écran
\draw[very thick,poppurple] (9,-0.5) -- (9,3) node[above,font=\small]{écran};
\fill[poppurple] (9,2.52) circle (2pt) node[right,font=\small]{spot};
% déflexion Y
\draw[<->,popdark] (9.35,0) -- (9.35,2.52) node[midway,right,font=\small]{$Y$};
% distance D
\draw[<->,popmuted] (2.5,-0.35) -- (9,-0.35) node[midway,below,font=\small]{$D$};
% angle alpha
\draw[popdark] (5.5,0) arc (0:20:3.5);
\node[popdark,font=\small] at (6.2,0.3) {$\alpha$};
\draw[dashed,popmuted] (3,0.36) -- (6.5,0.36);
% vitesse v0
\draw[->,very thick,popdark] (-0.4,0) -- (0.7,0) node[midway,above,font=\small]{$\vec{v}_0$};
\end{tikzpicture}
\end{center}
\legende{Trajectoire complète de l'électron : parabole entre les plaques, puis droite tangente jusqu'à l'écran. La tangente passe par $I$, milieu des plaques ; la déflexion sur l'écran est $Y = D\tan\alpha$. Schéma qualitatif.}
\end{popfigure}

\courssub{Calcul de la déflexion sur l'écran}

\begin{methode}[Décomposer le problème en deux phases]
Pour déterminer la déflexion $Y$ sur l'écran, procède toujours en deux phases distinctes :
\begin{enumerate}
\item \textbf{Entre les plaques} ($0 \le x \le L$) : mouvement parabolique. Calcule la déviation à la sortie $y_S = \dfrac{qEL^2}{2mv_0^2}$ et l'angle de sortie $\tan\alpha = \dfrac{qEL}{mv_0^2}$.
\item \textbf{Après les plaques} : mouvement rectiligne uniforme. Utilise la propriété géométrique : le faisceau semble provenir du point $I$ (milieu des plaques), situé à la distance $D$ de l'écran. La déflexion s'obtient alors par une simple relation trigonométrique dans le triangle rectangle.
\end{enumerate}
\end{methode}

Dans le triangle rectangle de sommet $I$, de base $D$ (distance du milieu des plaques à l'écran) et de hauteur $Y$ :
\[ \tan\alpha = \frac{Y}{D} \quad\Longrightarrow\quad Y = D\,\tan\alpha. \]
En remplaçant $\tan\alpha$ par son expression :

\begin{propriete}[Déflexion électrostatique sur l'écran — démonstration exigible]
La déflexion du spot sur l'écran est :
\[ \boxed{\ Y = \frac{q\,U\,L\,D}{m\,v_0^2\,d}\ } \]
où $U$ est la tension appliquée entre les plaques de déviation, $L$ la longueur des plaques, $d$ leur écartement, $D$ la distance du milieu des plaques à l'écran, $q$ et $m$ la charge et la masse de la particule, $v_0$ sa vitesse d'entrée. La déflexion $Y$ est \textbf{proportionnelle à la tension $U$} : c'est le principe de fonctionnement du voltmètre à oscilloscope.
\end{propriete}

\textbf{Vérification dimensionnelle.} $[qU] = \text{J}$ (énergie), $[mv_0^2] = \text{J}$, donc le rapport $\dfrac{qU}{mv_0^2}$ est sans dimension ; il reste $\dfrac{L D}{d}$, homogène à une longueur. La formule est cohérente.

\textbf{Lien avec la tension d'accélération.} Si l'électron a été accéléré par une tension $U_{\text{acc}}$, alors $\dfrac{1}{2}mv_0^2 = |q|U_{\text{acc}}$, soit $mv_0^2 = 2|q|U_{\text{acc}}$. La déflexion s'écrit alors, en valeur absolue :
\[ Y = \frac{U\,L\,D}{2\,U_{\text{acc}}\,d}. \]
Remarque élégante : la charge et la masse de la particule ont disparu ! La déflexion ne dépend que des tensions et de la géométrie du tube.

\courssub{Sensibilité de l'oscilloscope}

\begin{definition}[Sensibilité verticale]
La \cle{sensibilité} $s$ de l'oscilloscope est le quotient de la déflexion par la tension appliquée :
\[ s = \frac{Y}{U} = \frac{qLD}{mv_0^2\,d} = \frac{LD}{2\,U_{\text{acc}}\,d}. \]
Elle s'exprime en $\text{m}\cdot\text{V}^{-1}$, ou plus couramment en $\text{cm}/\text{V}$. Sur un oscilloscope réel, on utilise plutôt son inverse, le \textbf{calibre} (en $\text{V}/\text{cm}$ ou $\text{V}/\text{div}$) : tension correspondant à un centimètre (ou une division) de déflexion.
\end{definition}

Pour mesurer une tension inconnue $U$, on mesure la déflexion $Y$ du spot sur l'écran gradué et on applique :
\[ U = \frac{Y}{s} = Y \times (\text{calibre}). \]

\begin{exoresolu}[Déflexion d'un faisceau d'électrons]
\textbf{Énoncé.} Dans un tube d'oscilloscope, des électrons (charge $q = -e = -\num{1,60e-19}\ \text{C}$, masse $m = \num{9,11e-31}\ \text{kg}$) sont accélérés sous la tension $U_{\text{acc}} = \SI{1000}{V}$, puis traversent des plaques de déviation de longueur $L = \SI{4,0}{cm}$, distantes de $d = \SI{1,0}{cm}$, soumises à la tension $U = \SI{80}{V}$. L'écran est situé à $D = \SI{20}{cm}$ du milieu des plaques.
\begin{enumerate}
\item Calculer la vitesse $v_0$ des électrons à l'entrée des plaques.
\item Calculer $\tan\alpha$ et l'angle $\alpha$.
\item En déduire la déflexion $Y$ sur l'écran, puis la sensibilité $s$ du tube.
\end{enumerate}
\tcblower
\textbf{Solution détaillée.}
\begin{enumerate}
\item Théorème de l'énergie cinétique entre le filament et l'anode du canon :
\[ v_0 = \sqrt{\frac{2 e U_{\text{acc}}}{m}} = \sqrt{\frac{2 \times \num{1,60e-19} \times 1000}{\num{9,11e-31}}} = \sqrt{\num{3,51e14}} \approx \num{1,87e7}\ \text{m}\cdot\text{s}^{-1}. \]
\item Déviation angulaire :
\[ \tan\alpha = \frac{e U L}{m v_0^2\, d} = \frac{e U L}{2 e U_{\text{acc}}\, d} = \frac{U L}{2 U_{\text{acc}}\, d} = \frac{80 \times \num{4,0e-2}}{2 \times 1000 \times \num{1,0e-2}} = \num{0,16}. \]
D'où $\alpha = \arctan(\num{0,16}) \approx 9,1^{\circ}$.
\item Déflexion sur l'écran :
\[ Y = D\tan\alpha = 0,20 \times \num{0,16} = \num{3,2e-2}\ \text{m} = 3,2\ \text{cm}. \]
Sensibilité : $s = \dfrac{Y}{U} = \dfrac{3,2\ \text{cm}}{80\ \text{V}} = \num{0,040}\ \text{cm}/\text{V}$, soit un calibre de $\SI{25}{V/cm}$.
\end{enumerate}
\textbf{Vérification par la formule complète} : $Y = \dfrac{ULD}{2U_{\text{acc}}d} = \dfrac{80 \times 4,0 \times 20}{2 \times 1000 \times 1,0}\ \text{cm} = 3,2\ \text{cm}$. Cohérent.
\end{exoresolu}

\begin{attention}[Ne confonds pas déviation à la sortie et déflexion sur l'écran !]
L'erreur la plus fréquente au Baccalauréat est de donner comme réponse finale la déviation à la sortie des plaques :
\[ y_S = \frac{qUL^2}{2mv_0^2\,d}, \]
alors que l'énoncé demande la déflexion $Y$ \textbf{sur l'écran}. En utilisant $mv_0^2 = 2|q|U_{\text{acc}}$, l'expression correcte de $y_S$ en fonction de $U_{\text{acc}}$ est :
\[ y_S = \frac{U L^2}{4 U_{\text{acc}} d}. \]
Avec les données de l'exercice résolu, $y_S = \dfrac{80 \times (4,0)^2}{4 \times 1000 \times 1,0}\ \text{cm} = 0,32\ \text{cm}$, bien différent de $Y = 3,2\ \text{cm}$ ! Retiens bien :
\begin{itemize}
\item $y_S$ : ordonnée du point de sortie $S$, à l'extrémité des plaques ;
\item $Y$ : ordonnée du spot sur l'écran, après le trajet rectiligne sur la distance $D$.
\end{itemize}
On a d'ailleurs $Y = D\tan\alpha = \dfrac{2D}{L}\,y_S$ : la déflexion sur l'écran est toujours bien supérieure à la déviation à la sortie (ici $\dfrac{2D}{L} = 10$, donc $Y = 10 \times 0,32 = 3,2\ \text{cm}$). Autre piège : la distance $D$ se mesure depuis le \textbf{milieu} des plaques, pas depuis leur extrémité ; si l'énoncé donne la distance entre l'extrémité des plaques et l'écran, ajoute $L/2$.
\end{attention}

\begin{exercice}[Mesure d'une tension à l'oscilloscope]
Le tube de l'exercice résolu précédent ($L = \SI{4,0}{cm}$, $d = \SI{1,0}{cm}$, $D = \SI{20}{cm}$, $U_{\text{acc}} = \SI{1000}{V}$) est utilisé pour mesurer une tension inconnue $U'$. On observe une déflexion du spot de $Y' = \SI{2,0}{cm}$ vers le bas.
\begin{enumerate}
\item Déterminer la valeur de $U'$.
\item Que vaut la déflexion si l'on double la tension d'accélération, la tension $U'$ restant inchangée ? Commenter.
\end{enumerate}
\end{exercice}

\begin{corrige}[Exercice : mesure d'une tension]
\begin{enumerate}
\item La déflexion est proportionnelle à la tension : $U' = \dfrac{Y'}{s}$ avec $s = \num{0,040}\ \text{cm}/\text{V}$. Donc $U' = \dfrac{2,0}{\num{0,040}} = \SI{50}{V}$. Le signe négatif de la déflexion (vers le bas) indique que la polarité de $U'$ est opposée à celle de l'exercice résolu.
\item La formule $Y = \dfrac{ULD}{2U_{\text{acc}}d}$ montre que $Y$ est inversement proportionnelle à $U_{\text{acc}}$ : en doublant $U_{\text{acc}}$, la déflexion est divisée par deux : $Y'' = \SI{1,0}{cm}$. Interprétation : des électrons plus rapides passent moins de temps entre les plaques, ils sont donc moins déviés. La sensibilité du tube diminue quand on accélère davantage les électrons.
\end{enumerate}
\end{corrige}

\begin{aretenir}[L'essentiel de la section]
\begin{itemize}
\item Après les plaques, le mouvement est \textbf{rectiligne uniforme}, tangent à la parabole ; la tangente passe par le milieu $I$ des plaques.
\item Déflexion sur l'écran : $Y = D\tan\alpha = \dfrac{qULD}{mv_0^2\,d} = \dfrac{ULD}{2U_{\text{acc}}d}$.
\item $Y$ est proportionnelle à $U$ : l'oscilloscope est un voltmètre capable de suivre des tensions variables rapides.
\item Sensibilité : $s = \dfrac{Y}{U}$ (en cm/V) ; calibre = inverse de la sensibilité (en V/cm).
\end{itemize}
\end{aretenir}

Nous savons désormais dévier un faisceau d'électrons à volonté grâce à un champ électrique. Mais un champ \emph{magnétique} peut aussi courber la trajectoire d'une particule chargée, d'une manière radicalement différente : c'est ce que nous étudierons dans la section suivante, où nous verrons le mouvement circulaire uniforme et ses spectaculaires applications (spectrographe de masse, cyclotron).

\coursec{Mouvement dans un champ magnétique uniforme}

Dans la section précédente, nous avons vu comment un \cle{champ électrique} uniforme dévie un faisceau d'électrons et permet de tracer des courbes sur l'écran d'un oscilloscope. Mais l'électricité n'est pas la seule à pouvoir piloter des particules chargées : le \cle{champ magnétique} en est tout aussi capable, et de manière radicalement différente. Là où le champ électrique accélère ou freine, le champ magnétique \emph{courbe sans jamais accélérer}. Cette propriété, en apparence modeste, est à la base d'instruments fondamentaux de la physique moderne : le spectrographe de masse, le cyclotron, les anneaux du CERN, et même la protection que le champ magnétique terrestre offre au Cameroun et à toute la planète contre le vent solaire. C'est l'objet de cette section.

\courssub{Situation d'étude et force de Lorentz magnétique}

Considérons une particule ponctuelle de \cle{charge} $q$, de \cle{masse} $m$, pénétrant avec une vitesse initiale $\vec{v_0}$ dans une région de l'espace où règne un champ magnétique $\vec{B}$ \cle{uniforme} (même direction, même sens, même norme en tout point). On se place dans le cas fondamental où $\vec{v_0}$ est \textbf{perpendiculaire} à $\vec{B}$. On néglige le poids de la particule devant la force magnétique (hypothèse excellente : pour un proton dans un champ de $\SI{0,1}{T}$ à $\SI{1,0e6}{m.s^{-1}}$, la force magnétique vaut environ $10^{10}$ fois son poids !).

\begin{definition}[Force magnétique de Lorentz]
Une particule de charge $q$ animée d'une vitesse $\vec{v}$ dans un champ magnétique $\vec{B}$ subit la \cle{force magnétique} (ou force de Lorentz magnétique) :
\[ \vec{F} = q\,\vec{v} \wedge \vec{B} \]
Sa norme vaut $F = |q|\,v\,B\,|\sin\alpha|$, où $\alpha$ est l'angle entre $\vec{v}$ et $\vec{B}$. Elle est \textbf{perpendiculaire au plan} $(\vec{v}, \vec{B})$, donc en particulier perpendiculaire à la vitesse.
\end{definition}

\begin{attention}[Conditions d'application]
La force magnétique exige \textbf{deux conditions} : la particule doit être \cle{chargée} ($q \neq 0$) \textbf{et} \cle{en mouvement} ($v \neq 0$). Une charge immobile dans un champ magnétique ne subit \emph{aucune} force magnétique. De plus, si $\vec{v}$ est parallèle à $\vec{B}$ ($\alpha = 0$ ou $\pi$), alors $\vec{F} = \vec{0}$ : la particule poursuit son mouvement rectiligne uniforme.
\end{attention}

\courssub{Conséquence énergétique : la force magnétique ne travaille pas}

C'est le point de départ de tout le raisonnement. À chaque instant, $\vec{F} \perp \vec{v}$, donc la puissance de la force magnétique est nulle :
\[ \mathcal{P} = \vec{F} \cdot \vec{v} = \left(q\,\vec{v} \wedge \vec{B}\right) \cdot \vec{v} = 0 \]
car un produit vectoriel est toujours orthogonal à chacun de ses deux facteurs.

\begin{propriete}[La force magnétique ne travaille jamais]
La puissance de la force magnétique étant nulle à tout instant, son travail est nul sur tout déplacement. D'après le théorème de l'énergie cinétique, $\Delta E_c = W(\vec{F}) = 0$ : l'\cle{énergie cinétique} de la particule est \textbf{conservée}, donc la \textbf{norme de la vitesse est constante}. La force magnétique courbe la trajectoire sans jamais modifier la valeur de la vitesse.
\end{propriete}

C'est une différence essentielle avec le champ électrique : le champ électrique travaille (canon à électrons, accélération), le champ magnétique, lui, se contente de \emph{diriger}. C'est un « guide » et non un « moteur ».

\courssub{Démonstration : mouvement circulaire uniforme}

\begin{propriete}[Théorème — démonstration exigible au Bac]
Une particule chargée pénétrant dans un champ magnétique uniforme $\vec{B}$ avec une vitesse $\vec{v_0}$ \textbf{perpendiculaire} à $\vec{B}$ décrit, dans le plan perpendiculaire à $\vec{B}$, un \cle{mouvement circulaire uniforme} de rayon
\[ R = \dfrac{m\,v_0}{|q|\,B} \]
et de période
\[ T = \dfrac{2\pi m}{|q|\,B} \]
indépendante de la vitesse de la particule.
\end{propriete}

\begin{experience}[Démonstration complète (à savoir refaire)]
\textbf{Étape 1 : le mouvement est plan.} Appliquons la deuxième loi de Newton dans le référentiel galiléen du laboratoire :
\[ m\vec{a} = q\,\vec{v} \wedge \vec{B} \]
L'accélération $\vec{a}$ est à tout instant perpendiculaire à $\vec{B}$. La composante de la vitesse selon $\vec{B}$ est donc constante ; or elle est nulle initialement ($\vec{v_0} \perp \vec{B}$), donc elle reste nulle : le mouvement s'effectue entièrement dans le \textbf{plan perpendiculaire à} $\vec{B}$ passant par le point d'entrée.

\textbf{Étape 2 : le mouvement est uniforme (base de Frenet).} Dans la \cle{base de Frenet} $(\vec{\tau}, \vec{n})$, l'accélération s'écrit :
\[ \vec{a} = \dfrac{dv}{dt}\,\vec{\tau} + \dfrac{v^2}{\rho}\,\vec{n} \]
où $\rho$ est le rayon de courbure de la trajectoire. La force magnétique étant perpendiculaire à $\vec{v}$, elle n'a \textbf{aucune composante tangentielle} : $\vec{F}\cdot\vec{\tau} = 0$. La projection de la loi de Newton sur $\vec{\tau}$ donne :
\[ m\,\dfrac{dv}{dt} = 0 \quad\Longrightarrow\quad v = \text{constante} = v_0 \]
Le mouvement est \textbf{uniforme} (ce qui confirme le bilan énergétique précédent).

\textbf{Étape 3 : le rayon de courbure est constant.} Comme $\vec{v} \perp \vec{B}$, la norme de la force vaut $F = |q|\,v_0\,B$. La projection de la loi de Newton sur la normale $\vec{n}$ (la force, perpendiculaire à $\vec{v}$ dans le plan de la trajectoire, est dirigée vers le centre de courbure) donne :
\[ m\,\dfrac{v_0^2}{\rho} = |q|\,v_0\,B \quad\Longrightarrow\quad \rho = \dfrac{m\,v_0}{|q|\,B} = \text{constante} \]
car $m$, $v_0$, $|q|$ et $B$ sont constants. Une trajectoire plane dont le rayon de courbure est constant est un \textbf{cercle} de rayon $R = \dfrac{m v_0}{|q| B}$.

\textbf{Étape 4 : la période.} La particule parcourt la circonférence $2\pi R$ à la vitesse constante $v_0$ :
\[ T = \dfrac{2\pi R}{v_0} = \dfrac{2\pi}{v_0}\times\dfrac{m v_0}{|q| B} = \dfrac{2\pi m}{|q| B} \]
La vitesse disparaît de l'expression : la période ne dépend que de $m$, $|q|$ et $B$. \textbf{C.Q.F.D.}

\textbf{Analyse dimensionnelle (vérification) :}
$[R] = \dfrac{\text{kg}\cdot\text{m}\cdot\text{s}^{-1}}{\text{C}\cdot\text{T}} = \dfrac{\text{kg}\cdot\text{m}\cdot\text{s}^{-1}}{\text{C}\cdot(\text{N}\cdot\text{A}^{-1}\cdot\text{m}^{-1})} = \dfrac{\text{kg}\cdot\text{m}^2\cdot\text{s}^{-1}\cdot\text{A}}{\text{C}\cdot\text{N}} = \dfrac{\text{kg}\cdot\text{m}^2\cdot\text{s}^{-1}\cdot(\text{C}\cdot\text{s}^{-1})}{\text{C}\cdot\text{kg}\cdot\text{m}\cdot\text{s}^{-2}} = \text{m}$.
$[T] = \dfrac{\text{kg}}{\text{C}\cdot\text{T}} = \dfrac{\text{kg}}{\text{C}\cdot(\text{N}\cdot\text{A}^{-1}\cdot\text{m}^{-1})} = \dfrac{\text{kg}\cdot\text{m}\cdot\text{A}}{\text{C}\cdot\text{N}} = \dfrac{\text{kg}\cdot\text{m}\cdot(\text{C}\cdot\text{s}^{-1})}{\text{C}\cdot\text{kg}\cdot\text{m}\cdot\text{s}^{-2}} = \text{s}$.
\end{experience}

\begin{attention}[Piège classique : la période ne dépend pas de la vitesse]
Beaucoup d'élèves pensent qu'une particule plus rapide « tourne plus vite ». C'est faux ! Si $v_0$ double, le rayon $R$ double aussi, et la circonférence $2\pi R$ double : le temps de parcours $T = 2\pi m/(|q|B)$ reste \textbf{identique}. Cette remarquable propriété — dite \emph{isochronisme} du mouvement cyclotron — est la clé de fonctionnement du cyclotron, que nous étudierons dans la section suivante. De même, la pulsation $\omega = \dfrac{v_0}{R} = \dfrac{|q|B}{m}$ (appelée \cle{pulsation cyclotron}) est indépendante de $v_0$.
\end{attention}

\begin{popfigure}
\begin{center}
\begin{tikzpicture}[scale=1.1]
% Champ B sortant : points (représentation standard : pointe de flèche = point)
\foreach \x in {-3,-2,-1,1,2,3}{
  \foreach \y in {-2.5,-1.5,-0.5,0.5,1.5,2.5}{
    \fill[popblue] (\x,\y) circle (0.045);
    \draw[popblue] (\x,\y) circle (0.09);
  }
}
\node[popblue] at (3.4,2.5) {$\vec{B}$ \textbf{(sortant)}};
% Trajectoire circulaire (sens trigonométrique pour q>0, B sortant)
\draw[poporange, very thick] (0,0) circle (1.6);
\fill (0,0) circle (0.04) node[below right=-2pt] {$O$};
\draw[popmuted, dashed] (0,0) -- (1.6,0) node[midway, below] {$R$};
% Points cardinaux
\coordinate (P1) at (1.6,0);   % Est
\coordinate (P2) at (0,1.6);   % Nord
\coordinate (P3) at (-1.6,0);  % Ouest
\coordinate (P4) at (0,-1.6);  % Sud
\foreach \P in {P1,P2,P3,P4} \fill[popdark] (\P) circle (0.06);
% Vecteurs vitesse (vert) et force (rose) pour q>0, B sortant -> rotation trigonométrique
% En P1 (Est) : v vers le Sud (bas), F vers l'Ouest (gauche, centripète)
\draw[-{Stealth[length=3mm]}, popgreen, very thick] (P1) -- ++(0,-1.1) node[right] {$\vec{v}$};
\draw[-{Stealth[length=3mm]}, poppink, very thick] (P1) -- ++(-1.0,0) node[above] {$\vec{F}$};
% En P2 (Nord) : v vers l'Est (droite), F vers le Sud (bas, centripète)
\draw[-{Stealth[length=3mm]}, popgreen, very thick] (P2) -- ++(1.1,0) node[above] {$\vec{v}$};
\draw[-{Stealth[length=3mm]}, poppink, very thick] (P2) -- ++(0,-1.0) node[right] {$\vec{F}$};
% En P3 (Ouest) : v vers le Nord (haut), F vers l'Est (droite, centripète)
\draw[-{Stealth[length=3mm]}, popgreen, very thick] (P3) -- ++(0,1.1) node[left] {$\vec{v}$};
\draw[-{Stealth[length=3mm]}, poppink, very thick] (P3) -- ++(1.0,0) node[below] {$\vec{F}$};
% En P4 (Sud) : v vers l'Ouest (gauche), F vers le Nord (haut, centripète)
\draw[-{Stealth[length=3mm]}, popgreen, very thick] (P4) -- ++(-1.1,0) node[below] {$\vec{v}$};
\draw[-{Stealth[length=3mm]}, poppink, very thick] (P4) -- ++(0,1.0) node[left] {$\vec{F}$};
% Flèche de rotation globale
\draw[poporange, thick, ->] (2.2, 1.8) arc (45:135:2.5) node[above] {\small sens trigonométrique};
\end{tikzpicture}
\end{center}
\legende{Particule de charge $q>0$ dans un champ $\vec{B}$ uniforme \textbf{sortant} (points). La force $\vec{F} = q\,\vec{v}\wedge\vec{B}$ est centripète : la trajectoire est un cercle parcouru à vitesse constante dans le sens trigonométrique.}
\end{popfigure}

\courssub{Sens de rotation : le rôle du signe de la charge}

Le produit vectoriel $\vec{v} \wedge \vec{B}$ donne le sens de la force pour une charge \textbf{positive}. Pour une charge \textbf{négative} (électron, ion négatif), le facteur $q<0$ \textbf{inverse} le sens de la force, donc le sens de rotation.

\begin{methode}[Déterminer le sens de courbure]
\begin{enumerate}
\item Représenter $\vec{v}$ et $\vec{B}$ au point considéré.
\item Orienter le produit vectoriel $\vec{v} \wedge \vec{B}$ (règle du tire-bouchon ou des trois doigts de la main droite : index $\vec{v}$, majeur $\vec{B}$, pouce $\vec{v}\wedge\vec{B}$).
\item Si $q > 0$ : $\vec{F}$ est dans le sens de $\vec{v} \wedge \vec{B}$. Si $q < 0$ : $\vec{F}$ est dans le \textbf{sens opposé}.
\item La force pointe toujours vers le \textbf{centre} du cercle : la trajectoire s'enroule dans le sens imposé par $\vec{F}$.
\end{enumerate}
Exemple : avec $\vec{B}$ sortant de la feuille, une charge positive tourne dans le sens trigonométrique (figure ci-dessus), un électron tourne dans le sens horaire.
\end{methode}

\courssub{Exemple chiffré : un proton dans un champ magnétique}

\begin{exemplebox}[Proton dans un champ de \SI{0,1}{T}]
Un proton ($m = \num{1,67e-27}$ kg, $q = +e = \num{1,60e-19}$ C) pénètre avec une vitesse $v_0 = \SI{1,0e6}{m.s^{-1}}$, perpendiculaire à un champ magnétique uniforme $B = \SI{0,10}{T}$.

\textbf{Rayon de la trajectoire :}
\[ R = \dfrac{m v_0}{|q| B} = \dfrac{\num{1,67e-27}\times\num{1,0e6}}{\num{1,60e-19}\times 0,10} = \dfrac{\num{1,67e-21}}{\num{1,60e-20}} \approx \SI{0,104}{m} \approx \SI{10,4}{cm} \]

\textbf{Période de révolution :}
\[ T = \dfrac{2\pi m}{|q| B} = \dfrac{2\pi\times\num{1,67e-27}}{\num{1,60e-19}\times 0,10} \approx \SI{6,6e-7}{s} \]
soit environ $0,66\ \mu$s par tour, c'est-à-dire plus d'\textbf{un million et demi de tours par seconde} !

\textbf{Vérification énergétique :} l'énergie cinétique $E_c = \tfrac{1}{2}mv_0^2 \approx \num{8,4e-16}$ J $\approx 5,2$ keV reste constante durant tout le mouvement : le champ magnétique n'a fourni ni prélevé aucune énergie.
\end{exemplebox}

\begin{exoresolu}[Électron dans le champ magnétique terrestre]
Un électron ($m = \num{9,11e-31}$ kg, $q = -e$) est émis avec une vitesse $v_0 = \SI{2,0e6}{m.s^{-1}}$ perpendiculairement à un champ magnétique uniforme de norme $B = \SI{5,0e-5}{T}$ (ordre de grandeur du champ magnétique terrestre). Déterminer la nature du mouvement, le rayon et la période de la trajectoire.
\tcblower
\textbf{Nature du mouvement :} $\vec{v_0} \perp \vec{B}$ et $\vec{B}$ uniforme : le mouvement est \textbf{circulaire uniforme} (l'électron, de charge négative, tourne en sens inverse d'une charge positive).

\textbf{Rayon :}
\[ R = \dfrac{m v_0}{|q| B} = \dfrac{\num{9,11e-31}\times\num{2,0e6}}{\num{1,60e-19}\times\num{5,0e-5}} = \dfrac{\num{1,82e-24}}{\num{8,0e-24}} \approx \SI{0,23}{m} \]

\textbf{Période :}
\[ T = \dfrac{2\pi m}{|q| B} = \dfrac{2\pi\times\num{9,11e-31}}{\num{1,60e-19}\times\num{5,0e-5}} \approx \SI{7,2e-7}{s} \]
Remarque : si l'électron avait une vitesse dix fois plus grande, $R$ vaudrait $\SI{2,3}{m}$ mais $T$ resterait égale à $\SI{7,2e-7}{s}$.
\end{exoresolu}

\courssub{Complément : vitesse initiale non perpendiculaire à B — trajectoire hélicoïdale}

Que se passe-t-il si $\vec{v_0}$ n'est \textbf{pas} perpendiculaire à $\vec{B}$ ? On décompose alors la vitesse initiale en deux composantes :
\begin{itemize}
\item la composante \textbf{parallèle} $\vec{v_{0/\!/}}$ (selon $\vec{B}$) : elle ne subit aucune force ($\vec{v_{0/\!/}} \wedge \vec{B} = \vec{0}$), donc elle reste \textbf{constante} : mouvement rectiligne uniforme le long de l'axe du champ ;
\item la composante \textbf{perpendiculaire} $\vec{v_{0\perp}}$ : elle engendre un mouvement circulaire uniforme de rayon $R = \dfrac{m v_{0\perp}}{|q| B}$, comme démontré précédemment.
\end{itemize}
La combinaison des deux mouvements donne une \cle{trajectoire hélicoïdale} : la particule s'enroule autour des lignes de champ tout en avançant le long de celles-ci. Le \cle{pas} de l'hélice (distance parcourue parallèlement à $\vec{B}$ pendant un tour) vaut :
\[ p = v_{0/\!/}\times T = \dfrac{2\pi m\, v_{0/\!/}}{|q| B} \]
Ce résultat est \textbf{admis} au programme, mais sa compréhension est un excellent entraînement : il illustre le principe d'indépendance des mouvements selon des directions perpendiculaires.

\begin{popfigure}
\begin{center}
\begin{tikzpicture}[scale=0.9]
% Axe du champ
\draw[-{Stealth[length=3mm]}, popblue, very thick] (-0.5,0) -- (9.5,0) node[right] {$\vec{B}$};
% Hélice (paramétrisation : x=t, y=R cos(wt), z=R sin(wt) -> projection xOy)
\draw[poporange, very thick, domain=0:8.4, samples=300, variable=\t]
  plot ({\t}, {1.1*sin(260*\t)});
% Cylindre en pointillés (guides visuels)
\draw[popmuted, dashed] (0,1.1) -- (8.4,1.1);
\draw[popmuted, dashed] (0,-1.1) -- (8.4,-1.1);
\draw[popmuted, dashed] (8.4,0) ellipse (0.35 and 1.1);
% Point de départ et v0
\fill[popdark] (0,0) circle (0.07);
\draw[-{Stealth[length=3mm]}, popgreen, very thick] (0,0) -- (1.5,0.85) node[above right] {$\vec{v_0}$};
\draw[-{Stealth[length=3mm]}, popgreen!60, thick] (0,0) -- (1.5,0) node[below right=-1pt] {$\vec{v_{0/\!/}}$};
\draw[-{Stealth[length=3mm]}, popgreen!60, thick] (0,0) -- (0,0.85) node[left] {$\vec{v_{0\perp}}$};
\draw[popgreen!60, dashed] (1.5,0) -- (1.5,0.85);
\draw[popgreen!60, dashed] (0,0.85) -- (1.5,0.85);
% Pas de l'hélice
\draw[{Stealth[length=2.5mm]}-{Stealth[length=2.5mm]}, poppurple, thick] (2.42,-1.5) -- (4.85,-1.5) node[midway, below] {$p$};
\draw[popmuted, dotted] (2.42,0) -- (2.42,-1.6);
\draw[popmuted, dotted] (4.85,0) -- (4.85,-1.6);
\end{tikzpicture}
\end{center}
\legende{Trajectoire hélicoïdale lorsque $\vec{v_0}$ n'est pas perpendiculaire à $\vec{B}$ : mouvement circulaire uniforme dans le plan perpendiculaire à $\vec{B}$ combiné à une translation uniforme le long de $\vec{B}$.}
\end{popfigure}

\begin{savaistu}[Les aurores polaires : des hélices dans le ciel]
Le Soleil éjecte en permanence un « vent » de particules chargées (protons, électrons). En approchant de la Terre, ces particules rencontrent le \cle{champ magnétique terrestre} et s'enroulent en \textbf{hélices} le long de ses lignes de champ — exactement le mouvement que nous venons d'étudier ! Comme les lignes de champ convergent vers les pôles magnétiques, les particules y sont canalisées et excitent les atomes de la haute atmosphère, qui émettent alors de la lumière : ce sont les \textbf{aurores polaires} (boréales dans l'hémisphère nord, australes au sud). Sans ce bouclier magnétique, le vent solaire bombarderait directement la surface de la planète. Au Cameroun, proche de l'équateur magnétique, le champ terrestre nous protège silencieusement, jour et nuit.
\end{savaistu}

\begin{aretenir}[L'essentiel de la section]
\begin{itemize}
\item Force magnétique : $\vec{F} = q\,\vec{v}\wedge\vec{B}$, perpendiculaire à $\vec{v}$ : elle \textbf{ne travaille pas}, la norme de la vitesse est conservée.
\item Si $\vec{v_0} \perp \vec{B}$ : \textbf{mouvement circulaire uniforme} de rayon $R = \dfrac{m v_0}{|q| B}$ et de période $T = \dfrac{2\pi m}{|q| B}$, \textbf{indépendante de la vitesse}.
\item Le sens de rotation dépend du \textbf{signe de} $q$ : pour $q<0$, inverser le sens donné par $\vec{v}\wedge\vec{B}$.
\item Si $\vec{v_0}$ a une composante parallèle à $\vec{B}$ : trajectoire \textbf{hélicoïdale} (admis).
\end{itemize}
\end{aretenir}

\coursec{Applications : spectrographe de masse et cyclotron}

Dans la section précédente, nous avons établi un résultat capital : une particule chargée pénétrant dans un champ magnétique uniforme $\vec{B}$ avec une vitesse $\vec{v}$ perpendiculaire à $\vec{B}$ décrit un cercle de rayon
\[ R = \frac{mv}{|q|B} \]
parcouru à vitesse constante. Ce rayon dépend de la masse $m$, de la charge $q$ et de la vitesse $v$ de la particule. C'est précisément cette dépendance que les physiciens exploitent dans deux dispositifs emblématiques : le \cle{spectrographe de masse}, qui trie les particules selon leur masse, et le \cle{cyclotron}, qui communique aux particules des énergies considérables. Ces deux machines illustrent magnifiquement la complémentarité des deux champs : le champ électrique \emph{accélère} (il travaille), le champ magnétique \emph{dévie} (il ne travaille pas).

\courssub{Le spectrographe de masse}

\textbf{Principe général.} Comment peser un atome ? Impossible avec une balance, même de laboratoire : la masse d'un atome de chlore vaut environ \SI{6e-26}{kg}. L'astuce consiste à transformer la mesure de masse en mesure de \emph{longueur} : on ionise les atomes, on les accélère tous avec la même tension, puis on les laisse décrire des demi-cercles dans un champ magnétique. Les plus lourds, plus difficiles à courber, décrivent des cercles plus grands. Il suffit alors de mesurer la position des impacts sur une plaque photographique.

\begin{definition}[Spectrographe de masse]
Un \cle{spectrographe de masse} est un appareil qui sépare des ions de même charge mais de masses différentes, en exploitant la déviation magnétique de leur trajectoire. Il comporte quatre zones :
\begin{enumerate}
\item la \textbf{chambre d'ionisation} : les atomes (ou molécules) sont transformés en ions de charge $q$ (souvent $q = +e$) ;
\item la \textbf{chambre d'accélération} : les ions traversent une tension $U$ qui leur communique une vitesse $v$ ;
\item la \textbf{chambre de déviation} : règne un champ magnétique uniforme $\vec{B}$ perpendiculaire à la vitesse ; les ions décrivent un demi-cercle ;
\item le \textbf{détecteur} (plaque photographique ou capteur) : il enregistre la position des impacts.
\end{enumerate}
\end{definition}

\begin{popfigure}
\begin{center}
\begin{tikzpicture}[scale=1.0, >=Stealth]
% Chambre d'ionisation
\draw[fill=popgoldL] (0,0.4) rectangle (1.2,1.6);
\node[align=center, font=\small] at (0.6,1.0) {source\\ions};
% Fente F1
\draw[thick] (1.5,0.3) -- (1.5,0.85); \draw[thick] (1.5,1.15) -- (1.5,1.7);
\node[below, font=\small] at (1.5,0.3) {$F_1$};
% Plaques accélératrices
\draw[very thick, popblue] (2.6,0.3) -- (2.6,1.7);
\draw[very thick, poporange] (3.6,0.3) -- (3.6,1.7);
\node[above, font=\small, popblue] at (2.6,1.7) {$P_1$};
\node[above, font=\small, poporange] at (3.6,1.7) {$P_2$};
\node[below, font=\small] at (3.1,0.3) {$U$};
\draw[->, popblue, thick] (2.75,1.0) -- (3.45,1.0) node[midway, above, font=\small]{$\vec{E}$};
% Faisceau accéléré
\draw[->, thick, popdark] (1.6,1.0) -- (4.6,1.0);
% Zone de champ B
\draw[fill=popblueL, draw=popblue] (4.6,-0.2) rectangle (9.4,2.4);
\foreach \x in {5.1,5.9,6.7,7.5,8.3,9.1}
  \foreach \y in {0.2,1.0,1.8}
    \node[popblue, font=\small] at (\x,\y) {$\times$};
\node[popblue, font=\small] at (8.9,2.2) {$\vec{B}$};
% Demi-cercles
\draw[very thick, popink] (4.6,1.0) arc (180:0:1.35);
\draw[very thick, popgreen] (4.6,1.0) arc (180:0:1.75);
% Plaque photographique
\draw[very thick, popdark] (4.6,-0.2) -- (9.4,-0.2);
\fill[popink] (7.3,-0.2) circle (2pt) node[below, font=\small, popink]{$M_1$};
\fill[popgreen] (8.1,-0.2) circle (2pt) node[below, font=\small, popgreen]{$M_2$};
% Annotations rayons
\draw[<->, popink] (4.6,1.0) -- (5.95,1.0) node[midway, above, font=\small]{$R_1$};
\draw[<->, popgreen] (4.6,0.7) -- (6.35,0.7) node[midway, below, font=\small]{$R_2$};
\node[font=\small, popink] at (5.6,2.15) {$^{35}$Cl$^+$};
\node[font=\small, popgreen] at (6.6,2.35) {$^{37}$Cl$^+$};
\end{tikzpicture}
\end{center}
\legende{Schéma de principe du spectrographe de masse. Les ions légers ($^{35}$Cl$^+$) décrivent un demi-cercle plus petit que les ions lourds ($^{37}$Cl$^+$) ; les impacts $M_1$ et $M_2$ sur la plaque sont séparés.}
\end{popfigure}

\textbf{Étape 1 : accélération des ions.} Entre les plaques $P_1$ et $P_2$ règne un champ électrique uniforme. L'ion, supposé parti avec une vitesse négligeable en $F_1$, est accéléré sous la tension $U = V_{P_1} - V_{P_2} > 0$. Appliquons le théorème de l'énergie cinétique entre $P_1$ et $P_2$ ; seule la force électrique travaille :
\[ \frac{1}{2}mv^2 - 0 = |q|U \qquad \Longrightarrow \qquad v = \sqrt{\frac{2|q|U}{m}} \]
Ce résultat, déjà rencontré avec le canon à électrons, est la première clé du dispositif : tous les ions de même charge reçoivent la même \emph{énergie cinétique} $|q|U$, mais pas la même vitesse — les plus massifs sont plus lents.

\textbf{Étape 2 : déviation magnétique.} L'ion pénètre ensuite dans la chambre de déviation avec la vitesse $\vec{v}$ perpendiculaire à $\vec{B}$. D'après l'étude de la section précédente, la force de Lorentz $\vec{F} = q\,\vec{v}\wedge\vec{B}$, perpendiculaire à la vitesse, ne travaille pas : le mouvement est circulaire uniforme, de rayon
\[ R = \frac{mv}{|q|B} \]
L'ion décrit un \emph{demi-cercle} (car il entre par le bord de la zone de champ) et frappe la plaque photographique à une distance $2R$ de la fente d'entrée.

\begin{propriete}[Formule fondamentale du spectrographe de masse]
Dans un spectrographe de masse, le rayon de la trajectoire d'un ion de masse $m$, de charge $q$, accéléré sous la tension $U$ puis dévié par le champ $B$, est :
\[ R = \frac{1}{B}\sqrt{\frac{2mU}{|q|}} \]
Les ions sont donc séparés selon le rapport $\dfrac{m}{|q|}$ : à $U$, $B$ et $|q|$ fixés, \textbf{le rayon croît comme la racine carrée de la masse}.
\end{propriete}

\begin{experience}[Démonstration guidée de la formule $R = \dfrac{1}{B}\sqrt{\dfrac{2mU}{|q|}}$]
Il suffit d'enchaîner les deux étapes.
\begin{enumerate}
\item \textbf{Accélération} : le théorème de l'énergie cinétique donne $\dfrac{1}{2}mv^2 = |q|U$, d'où $v = \sqrt{\dfrac{2|q|U}{m}}$.
\item \textbf{Déviation} : dans le champ magnétique, $R = \dfrac{mv}{|q|B}$.
\item \textbf{Combinaison} : on injecte l'expression de $v$ dans celle de $R$ :
\[ R = \frac{m}{|q|B}\sqrt{\frac{2|q|U}{m}} = \frac{1}{B}\sqrt{\frac{m^2}{q^2}\times\frac{2|q|U}{m}} = \frac{1}{B}\sqrt{\frac{2mU}{|q|}} \]
\end{enumerate}
Vérification dimensionnelle : $\sqrt{\dfrac{\text{kg}\cdot\text{V}}{\text{C}}} = \sqrt{\dfrac{\text{kg}\cdot\text{kg}\cdot\text{m}^2\cdot\text{s}^{-3}\cdot\text{A}^{-1}}{\text{A}\cdot\text{s}}} = \text{kg}\cdot\text{m}\cdot\text{s}^{-2}\cdot\text{A}^{-1}$, et en divisant par le tesla ($\text{kg}\cdot\text{s}^{-2}\cdot\text{A}^{-1}$), on obtient bien des mètres. \textbf{Cette démonstration est exigible au Baccalauréat.}
\end{experience}

\begin{methode}[Résoudre un problème de spectrographe de masse]
Face à tout exercice de spectrographe, enchaîne toujours les deux étapes dans cet ordre :
\begin{enumerate}
\item \textbf{Énergie cinétique (accélération)} : écris $\dfrac{1}{2}mv^2 = |q|U$ et tire $v$ (ou vérifie l'homogénéité si $v$ est donnée) ;
\item \textbf{Rayon (déviation)} : écris $R = \dfrac{mv}{|q|B}$ et conclus ;
\item si l'énoncé demande directement le rayon en fonction de $m$, $U$, $q$, $B$, utilise la formule combinée $R = \dfrac{1}{B}\sqrt{\dfrac{2mU}{|q|}}$ ;
\item la position de l'impact sur la plaque est à la distance $d = 2R$ de la fente d'entrée (diamètre du demi-cercle, pas rayon !).
\end{enumerate}
\end{methode}

\begin{exoresolu}[Séparation des isotopes du chlore $^{35}$Cl et $^{37}$Cl]
Le chlore naturel est un mélange de deux isotopes : $^{35}$Cl (environ 75\,\%) et $^{37}$Cl (environ 25\,\%). Dans un spectrographe de masse, des ions Cl$^+$ ($q = e = \num{1,60e-19}$ C) sont accélérés sous $U = \SI{5,0}{kV}$ puis déviés par un champ $B = \SI{0,50}{T}$. On donne $m(^{35}\text{Cl}) = \SI{5,81e-26}{kg}$ et $m(^{37}\text{Cl}) = \SI{6,14e-26}{kg}$.
\begin{enumerate}
\item Calculer les rayons $R_1$ et $R_2$ des trajectoires des deux isotopes.
\item En déduire la distance $\Delta d$ séparant les deux impacts sur la plaque.
\end{enumerate}
\tcblower
\textbf{1. Rayons des trajectoires.} Appliquons la formule combinée $R = \dfrac{1}{B}\sqrt{\dfrac{2mU}{|q|}}$.
\begin{align*}
R_1 &= \frac{1}{0{,}50}\sqrt{\frac{2\times \num{5,81e-26}\times \num{5,0e3}}{\num{1,60e-19}}}
= 2\times\sqrt{\num{3,63e-3}} = 2\times \num{6,02e-2} \approx \SI{0,120}{m} \\[4pt]
R_2 &= \frac{1}{0{,}50}\sqrt{\frac{2\times \num{6,14e-26}\times \num{5,0e3}}{\num{1,60e-19}}}
= 2\times\sqrt{\num{3,84e-3}} = 2\times \num{6,20e-2} \approx \SI{0,124}{m}
\end{align*}
Soit $R_1 \approx \SI{12,0}{cm}$ et $R_2 \approx \SI{12,4}{cm}$. On vérifie au passage que $R_2/R_1 = \sqrt{37/35} \approx \num{1,028}$, conforme à la loi en $\sqrt{m}$.

\textbf{2. Distance entre les impacts.} Chaque impact se situe à une distance $2R$ de la fente d'entrée :
\[ \Delta d = 2R_2 - 2R_1 = 2\times(12{,}4 - 12{,}0) \approx \SI{0,68}{cm} \approx \SI{7}{mm} \]
Les deux taches sont séparées d'environ \SI{7}{mm} sur la plaque : largement de quoi les distinguer à l'œil nu. L'isotope le plus léger ($^{35}$Cl) frappe le plus près de la fente.
\end{exoresolu}

\begin{attention}[Pièges fréquents au Bac]
\begin{itemize}
\item Ne confonds pas \textbf{rayon} $R$ et \textbf{distance de l'impact} à la fente, qui vaut $2R$ (l'ion décrit un \emph{demi}-cercle).
\item La vitesse d'entrée dans le champ magnétique n'est \emph{pas} la même pour deux isotopes : c'est l'\textbf{énergie cinétique} $|q|U$ qui est commune, pas la vitesse.
\item Dans $R = \dfrac{1}{B}\sqrt{\dfrac{2mU}{|q|}}$, la charge $q$ est la valeur \emph{algébrique} en valeur absolue si l'on ne s'intéresse qu'au rayon : utilise $|q|$.
\item Le champ magnétique ne modifie jamais la norme de la vitesse : inutile de recalculer $v$ après la déviation.
\end{itemize}
\end{attention}

\begin{savaistu}[Dater les roches et contrôler les médicaments]
La spectrométrie de masse est partout. Les géologues s'en servent pour \textbf{dater les roches} : en mesurant le rapport des isotopes du plomb ou de l'argon dans un échantillon, on détermine son âge — c'est ainsi qu'on date les formations volcaniques de la ligne du Cameroun (mont Cameroun, monts Bambouto). Les laboratoires pharmaceutiques l'utilisent pour \textbf{contrôler la pureté des médicaments} : chaque impureté, même à l'état de traces, laisse sa signature en masse. Les douanes et la police scientifique détectent ainsi des traces d'explosifs ou de drogues, et les contrôles antidopage des sportifs reposent sur la même technique.
\end{savaistu}

\courssub{Le cyclotron}

\textbf{Le problème à résoudre.} Pour obtenir des particules de très haute énergie (par exemple des protons destinés à produire des isotopes médicaux ou à bombarder des noyaux), il faudrait une tension d'accélération de plusieurs millions de volts : techniquement impossible à maintenir en continu. L'idée géniale d'Ernest Lawrence (1930) : utiliser une tension \emph{modeste} (quelques dizaines de kilovolts), mais la faire travailler \emph{de nombreuses fois} en faisant repasser la particule encore et encore dans la zone d'accélération. Pour cela, on replie la trajectoire sur elle-même grâce à un champ magnétique.

\begin{definition}[Cyclotron]
Un \cle{cyclotron} est un accélérateur de particules constitué de :
\begin{itemize}
\item deux demi-cylindres creux appelés \textbf{« dés »} (en anglais \emph{dees}), $D_1$ et $D_2$, séparés par un étroit intervalle ;
\item un \textbf{champ magnétique uniforme} $\vec{B}$ perpendiculaire au plan des dés, qui courbe la trajectoire en demi-cercles ;
\item une \textbf{tension alternative} $u(t)$ appliquée entre les deux dés, qui accélère la particule à chaque traversée de l'intervalle.
\end{itemize}
La particule, injectée au centre, décrit une spirale dont le rayon augmente à chaque passage entre les dés, jusqu'à l'extraction en périphérie.
\end{definition}

\begin{popfigure}
\begin{center}
\begin{tikzpicture}[scale=0.9, >=Stealth]
% Les deux dés
\draw[very thick, popblue, fill=popblueL] (2.2,0) arc (0:180:2.2) -- cycle;
\draw[very thick, poporange, fill=poporangeL] (-2.2,0) arc (180:360:2.2) -- cycle;
\node[popblue, font=\small] at (1.2,1.5) {$D_1$};
\node[poporange, font=\small] at (-1.2,-1.5) {$D_2$};
% Champ B
\foreach \p in {(1.2,0.6),(0.6,1.6),(-1.2,0.6),(-0.6,-1.6),(1.2,-0.8),(-1.4,-0.6)}
  \node[popblue, font=\small] at \p {$\times$};
\node[popblue, font=\small] at (2.0,2.35) {$\vec{B}$};
% Spirale (demi-cercles croissants)
\draw[very thick, popink] (0.25,0) arc (0:180:0.25);
\draw[very thick, popink] (-0.5,0) arc (180:360:0.5);
\draw[very thick, popink] (0.85,0) arc (0:180:0.85);
\draw[very thick, popink] (-1.25,0) arc (180:360:1.25);
\draw[very thick, popink, ->] (1.65,0) arc (0:150:1.65);
% Source
\fill[popdark] (0,0) circle (2pt) node[above right=1pt, font=\small]{source};
% Intervalle accélérateur
\draw[<->, popgreen, thick] (0.06,2.35) -- (-0.06,2.35);
\node[popgreen, font=\small, align=center] at (0,2.75) {intervalle\\accélérateur};
% Tension alternative
\draw[popdark, thick] (2.3,-1.2) -- (3.2,-1.2);
\draw[popdark, thick] (-2.3,-1.2) -- (-3.2,-1.2);
\draw[popdark, thick] (3.2,-1.2) -- (3.2,-2.0) -- (-3.2,-2.0) -- (-3.2,-1.2);
\draw[popdark, domain=-2.6:2.6, samples=100, smooth, thick] plot (\x, {-2.6 + 0.35*sin(180*\x)});
\node[font=\small, popdark] at (0,-3.25) {$u(t) = U_m\cos(2\pi f t)$};
\end{tikzpicture}
\end{center}
\legende{Le cyclotron vu de dessus. À l'intérieur des dés, la trajectoire est un demi-cercle (champ $\vec{B}$ seul) ; l'accélération n'a lieu que dans l'intervalle entre les dés, où la tension alternative $u(t)$ crée un champ électrique.}
\end{popfigure}

\textbf{Pourquoi ça marche : la condition de résonance.} À chaque demi-tour, la particule traverse l'intervalle entre les dés. Pour qu'elle soit \emph{accélérée} (et non freinée) à chaque passage, il faut que la tension $u(t)$ ait changé de signe exactement pendant le demi-tour : la demi-période de la tension doit égaler la durée d'un demi-tour. Or cette durée est remarquable.

\begin{propriete}[Fréquence cyclotron]
La période du mouvement circulaire d'une particule dans un champ magnétique uniforme est
\[ T = \frac{2\pi m}{|q|B} \]
Elle est \textbf{indépendante de la vitesse} (et donc du rayon). La fréquence associée, dite \cle{fréquence cyclotron}, vaut :
\[ f = \frac{|q|B}{2\pi m} \]
\textbf{Conséquence fondamentale :} même si la particule accélère et que le rayon de sa trajectoire augmente, elle met toujours le \emph{même temps} pour effectuer un demi-tour. Une tension alternative de fréquence fixe $f = \dfrac{|q|B}{2\pi m}$ reste donc synchronisée avec la particule pendant toute l'accélération : c'est la \textbf{condition de résonance}, et c'est elle qui rend le cyclotron possible.
\end{propriete}

\begin{experience}[Démonstration guidée : période indépendante de la vitesse]
\begin{enumerate}
\item Le rayon du cercle est $R = \dfrac{mv}{|q|B}$.
\item La période est la durée d'un tour complet : $T = \dfrac{2\pi R}{v}$ (circonférence divisée par la vitesse, constante en norme).
\item On remplace $R$ : $T = \dfrac{2\pi}{v}\times\dfrac{mv}{|q|B} = \dfrac{2\pi m}{|q|B}$. La vitesse $v$ s'élimine : la période ne dépend que de $m$, $|q|$ et $B$.
\end{enumerate}
\end{experience}

\textbf{Énergie maximale atteinte.} À chaque traversée de l'intervalle, la particule gagne l'énergie $|q|U_m$ (si elle traverse au moment où la tension vaut $U_m$). Son rayon augmente donc régulièrement, jusqu'au rayon maximal $R$ imposé par la taille des dés, où on l'extrait.

\begin{propriete}[Énergie cinétique maximale en sortie de cyclotron]
Pour un cyclotron de rayon $R$ (rayon des dés), de champ magnétique $B$, l'énergie cinétique maximale d'une particule de masse $m$ et de charge $q$ est :
\[ E_c = \frac{q^2B^2R^2}{2m} \]
Remarquablement, cette énergie \textbf{ne dépend pas de la tension} $U_m$ : une tension plus grande fait simplement gagner l'énergie en moins de tours.
\end{propriete}

\begin{experience}[Démonstration guidée de $E_c = \dfrac{q^2B^2R^2}{2m}$]
\begin{enumerate}
\item En sortie, la particule décrit un cercle de rayon maximal $R$ à la vitesse $v_{\max}$ : $R = \dfrac{mv_{\max}}{|q|B}$, d'où $v_{\max} = \dfrac{|q|BR}{m}$.
\item L'énergie cinétique vaut alors :
\[ E_c = \frac{1}{2}mv_{\max}^2 = \frac{1}{2}m\left(\frac{|q|BR}{m}\right)^2 = \frac{q^2B^2R^2}{2m} \]
\end{enumerate}
Vérification dimensionnelle : $\dfrac{\text{C}^2\cdot\text{T}^2\cdot\text{m}^2}{\text{kg}} = \dfrac{\text{C}^2}{\text{kg}}\times\dfrac{\text{kg}^2}{\text{C}^2\cdot\text{s}^2}\times\text{m}^2 = \text{kg}\cdot\text{m}^2\cdot\text{s}^{-2}$, c'est-à-dire des joules.
\end{experience}

\begin{exemplebox}[Un cyclotron à protons]
Un cyclotron de rayon $R = \SI{0,50}{m}$ fonctionne avec un champ $B = \SI{1,2}{T}$ pour accélérer des protons ($m = \num{1,67e-27}$ kg, $q = e = \num{1,60e-19}$ C).
\begin{itemize}
\item \textbf{Fréquence de la tension} : $f = \dfrac{eB}{2\pi m} = \dfrac{\num{1,60e-19}\times 1{,}2}{2\pi\times \num{1,67e-27}} \approx \SI{1,83e7}{Hz} \approx \SI{18}{MHz}$ (ondes radio !).
\item \textbf{Énergie maximale} : $E_c = \dfrac{e^2B^2R^2}{2m} = \dfrac{(\num{1,60e-19})^2\times(1{,}2)^2\times(0{,}50)^2}{2\times \num{1,67e-27}} \approx \SI{2,8e-12}{J} \approx \SI{17}{MeV}$.
\end{itemize}
Une telle énergie exigerait une tension continue de 17 millions de volts ; le cyclotron l'obtient avec quelques kilovolts répétés des centaines de fois.
\end{exemplebox}

\begin{attention}[Où a lieu l'accélération ?]
Erreur classique au Bac : croire que la particule est accélérée en permanence. C'est faux !
\begin{itemize}
\item \textbf{À l'intérieur des dés}, il n'y a que le champ magnétique $\vec{B}$ (les dés métalliques creux écrantent le champ électrique) : la force de Lorentz ne travaille pas, la vitesse garde une norme \emph{constante}, la trajectoire est un demi-cercle.
\item \textbf{L'accélération a lieu uniquement dans l'intervalle entre les dés}, où la tension alternative crée un champ électrique qui travaille : $\Delta E_c = |q|\,|u| > 0$.
\item Autre piège : la fréquence de la tension alternative est $f = \dfrac{|q|B}{2\pi m}$, et non $\dfrac{|q|B}{\pi m}$ — la tension fait un \emph{demi}-cycle pendant chaque demi-tour, donc sa période complète égale la période de révolution.
\end{itemize}
\end{attention}

\begin{savaistu}[Des cyclotrons à l'hôpital]
Les cyclotrons modernes équipent de grands centres hospitaliers : ils produisent les isotopes radioactifs à vie courte (comme le fluor-18) utilisés en imagerie médicale TEP (\emph{tomographie par émission de positons}) pour détecter les cancers. L'Afrique s'équipe progressivement de tels centres de médecine nucléaire. À plus grande échelle, les descendants du cyclotron (synchrotrons) comme le LHC du CERN propulsent des protons à des énergies des milliards de fois supérieures, toujours en combinant champs électriques accélérateurs et champs magnétiques guidants — exactement le principe étudié dans cette leçon.
\end{savaistu}

\begin{aretenir}[L'essentiel de la section]
\begin{itemize}
\item \textbf{Spectrographe de masse} : accélération sous $U$ ($\frac{1}{2}mv^2 = |q|U$) puis déviation en demi-cercle dans $\vec{B}$ ; le rayon $R = \dfrac{1}{B}\sqrt{\dfrac{2mU}{|q|}}$ sépare les ions selon $\dfrac{m}{|q|}$ ; l'impact est à la distance $2R$ de la fente.
\item \textbf{Cyclotron} : deux dés sous champ $\vec{B}$, tension alternative entre les dés ; l'accélération n'a lieu que dans l'intervalle.
\item \textbf{Condition de résonance} : $f = \dfrac{|q|B}{2\pi m}$, indépendante de la vitesse — c'est la clé du fonctionnement.
\item \textbf{Énergie maximale} : $E_c = \dfrac{q^2B^2R^2}{2m}$, fixée par le rayon des dés et le champ, indépendante de la tension.
\item Méthode reine : \emph{énergie cinétique d'abord, rayon ensuite}.
\end{itemize}
\end{aretenir}

\coursec{Filtre de Wien : champs électrique et magnétique croisés}

Dans la section précédente, nous avons vu comment un champ magnétique uniforme sépare les particules selon leur rapport $q/m$ (spectrographe de masse) et comment il les accélère par cycles dans le cyclotron. Mais une difficulté pratique demeure : dans un spectrographe de masse, les ions produits par la source n'ont \emph{pas tous la même vitesse}. Or le rayon de la trajectoire circulaire $R = \dfrac{mv}{|q|B}$ dépend de $v$ : si les vitesses sont dispersées, les impacts sur la plaque photographique s'étalent et la mesure des masses devient imprécise. Il faut donc, \emph{en amont} du spectrographe, un dispositif capable de ne laisser passer que les particules ayant une vitesse bien précise. Ce dispositif, c'est le \cle{filtre de Wien}, aussi appelé \cle{sélecteur de vitesse}, qui combine un champ électrique et un champ magnétique \emph{croisés}. C'est l'aboutissement de cette leçon : nous allons enfin faire travailler $\vec{E}$ et $\vec{B}$ \emph{ensemble}.

\courssub{Description du dispositif et bilan des forces}

\textbf{Situation physique.} Une particule chargée, de charge $q$ et de masse $m$, pénètre dans une zone de l'espace où règnent simultanément :
\begin{itemize}
\item un champ électrique uniforme $\vec{E}$, créé par exemple par deux plaques planes parallèles verticales reliées à un générateur de tension ;
\item un champ magnétique uniforme $\vec{B}$, perpendiculaire à $\vec{E}$ ;
\item la particule arrive avec une vitesse initiale $\vec{v}_0$ perpendiculaire à la fois à $\vec{E}$ et à $\vec{B}$.
\end{itemize}
Les trois vecteurs $\vec{v}_0$, $\vec{E}$ et $\vec{B}$ forment donc un trièdre : c'est la configuration dite des \emph{champs croisés}.

\textbf{Bilan des forces.} Dans la zone des champs croisés, la particule est soumise à deux forces (son poids est négligeable devant les forces électromagnétiques, comme toujours pour des particules élémentaires) :
\begin{itemize}
\item la \cle{force électrique} $\vec{F}_e = q\vec{E}$, de direction fixe (celle de $\vec{E}$ si $q>0$, opposée si $q<0$), indépendante de la vitesse ;
\item la \cle{force magnétique} (force de Lorentz) $\vec{F}_m = q\,\vec{v}\wedge\vec{B}$, perpendiculaire à $\vec{v}$ et à $\vec{B}$, dont la norme vaut $F_m = |q|\,v\,B$ puisque $\vec{v}\perp\vec{B}$.
\end{itemize}

\textbf{Orientation relative.} Les sens de $\vec{E}$ et de $\vec{B}$ sont choisis de telle sorte que $\vec{F}_e$ et $\vec{F}_m$ soient \emph{colinéaires et de sens opposés}. La particule est alors « tirée » d'un côté par le champ électrique et de l'autre par le champ magnétique. Tout le jeu du filtre de Wien tient dans ce duel de forces.

\begin{popfigure}
\begin{tikzpicture}[scale=1.05, >=Stealth]
% Plaques du condensateur
\fill[popblueL] (-0.3,2.0) rectangle (7.3,2.3);
\fill[popblueL] (-0.3,-2.3) rectangle (7.3,-2.0);
\node[popblue, above] at (7.0,2.3) {\small $+$};
\node[popblue, below] at (7.0,-2.3) {\small $-$};
% Champ E (vers le bas)
\foreach \x in {1.2,2.7,4.2,5.7}{
  \draw[->, popblue, thick] (\x,1.7) -- (\x,0.9);
}
\node[popblue] at (6.5,1.35) {$\vec{E}$};
% Champ B (croix = rentrant)
\foreach \x in {1.2,2.7,4.2,5.7}{
  \foreach \y in {-1.4,-0.5,0.4}{
    \node[popgreen] at (\x,\y) {\small $\times$};
  }
}
\node[popgreen] at (6.5,-1.4) {$\vec{B}$ (rentrant)};
% Trajectoire rectiligne
\draw[popdark, very thick, dashed] (-0.3,0) -- (7.3,0);
% Particule
\fill[poporange] (3.5,0) circle (0.12);
\node[poporange, above right] at (3.55,0.05) {\small $q>0$};
% Vitesse
\draw[->, popdark, very thick] (3.5,0) -- (5.0,0) node[midway, below] {$\vec{v}$};
% Forces
\draw[->, popblue, very thick] (3.5,0) -- (3.5,-1.6) node[right, pos=0.6] {$\vec{F}_e = q\vec{E}$};
\draw[->, popgreen, very thick] (3.5,0) -- (3.5,1.6) node[right, pos=0.6] {$\vec{F}_m = q\vec{v}\wedge\vec{B}$};
% Fente de sortie
\fill[popcream] (7.3,-0.35) rectangle (7.6,0.35);
\draw[popdark, thick] (7.3,-0.35) -- (7.6,-0.35) (7.3,0.35) -- (7.6,0.35);
\node[popdark, right] at (7.6,0) {\small fente};
\end{tikzpicture}
\legende{Filtre de Wien : $\vec{E}$ dirigé vers le bas, $\vec{B}$ rentrant (croix), $\vec{v}$ horizontale. Pour une particule positive, $\vec{F}_e$ est vers le bas et $\vec{F}_m$ vers le haut. Si $v = E/B$, les deux forces se compensent exactement et la trajectoire reste rectiligne jusqu'à la fente de sortie.}
\end{popfigure}

\courssub{Condition de trajectoire rectiligne : la vitesse sélectionnée}

\begin{propriete}[Condition de non-déviation dans le filtre de Wien (démonstration exigible)]
Dans une zone où règnent un champ électrique uniforme $\vec{E}$ et un champ magnétique uniforme $\vec{B}$ perpendiculaires entre eux et perpendiculaires à la vitesse initiale $\vec{v}_0$ de la particule, orientés de sorte que les forces électrique et magnétique soient opposées, la particule traverse la zone \textbf{en ligne droite et à vitesse constante} si et seulement si :
\[ \boxed{\,v = \dfrac{E}{B}\,} \]
\end{propriete}

\begin{experience}[Démonstration guidée]
On raisonne dans le référentiel du laboratoire, supposé galiléen.
\begin{enumerate}
\item \textbf{Bilan des forces :} $\vec{F}_e = q\vec{E}$ et $\vec{F}_m = q\,\vec{v}\wedge\vec{B}$, colinéaires et de sens opposés par construction du dispositif. Le poids est négligeable.
\item \textbf{Deuxième loi de Newton :} $m\vec{a} = q\vec{E} + q\,\vec{v}\wedge\vec{B}$.
\item \textbf{Condition de mouvement rectiligne uniforme :} la trajectoire reste rectiligne le long de l'axe de $\vec{v}_0$ si l'accélération transversale est nulle, c'est-à-dire si la somme des forces s'annule :
\[ q\vec{E} + q\,\vec{v}\wedge\vec{B} = \vec{0} \quad\Longleftrightarrow\quad \vec{F}_e = -\vec{F}_m. \]
\item \textbf{Passage aux normes :} comme $\vec{v}\perp\vec{B}$, on a $F_m = |q|\,v\,B$, et $F_e = |q|\,E$. L'égalité des normes donne :
\[ |q|\,E = |q|\,v\,B \quad\Longleftrightarrow\quad v = \dfrac{E}{B}. \]
\item \textbf{Cohérence :} si les forces se compensent, $\vec{a} = \vec{0}$, donc $\vec{v}$ reste constante : la condition $v = E/B$, vraie à l'entrée, le demeure pendant toute la traversée. Le mouvement est bien rectiligne uniforme.
\end{enumerate}
\textbf{Analyse dimensionnelle :} $[E] = \mathrm{V\,m^{-1}} = \mathrm{kg\,m\,s^{-3}\,A^{-1}}$ et $[B] = \mathrm{T} = \mathrm{kg\,s^{-2}\,A^{-1}}$, donc $[E/B] = \mathrm{m\,s^{-1}}$ : c'est bien une vitesse.
\end{experience}

\textbf{Que deviennent les autres particules ?}
\begin{itemize}
\item Si $v > \dfrac{E}{B}$ : la force magnétique $|q|vB$ l'emporte sur la force électrique $|q|E$ ; la particule est déviée du côté de $\vec{F}_m$ et heurte une plaque ou la paroi.
\item Si $v < \dfrac{E}{B}$ : la force électrique l'emporte ; la particule est déviée du côté de $\vec{F}_e$ et est également interceptée.
\end{itemize}
Seules les particules dont la vitesse vaut exactement $E/B$ atteignent la fente de sortie. Le dispositif agit comme un \emph{filtre} : il « sélectionne » une vitesse. D'où son nom de \cle{sélecteur de vitesse} ou \cle{filtre de vitesse de Wien}.

\begin{attention}[La condition $v = E/B$ ne dépend ni de $q$ ni de $m$]
C'est le piège classique du baccalauréat : dans l'équation $|q|E = |q|vB$, la charge $q$ \textbf{se simplifie}, et la masse $m$ n'apparaît nulle part. Par conséquent :
\begin{itemize}
\item des particules de charges différentes (proton, particule $\alpha$, ion \ce{^{16}O+}...) ayant la \emph{même vitesse} $v = E/B$ traversent toutes sans déviation ;
\item le signe de $q$ ne change rien non plus : pour $q<0$, $\vec{F}_e$ et $\vec{F}_m$ s'inversent \emph{toutes les deux} et restent opposées.
\end{itemize}
Le filtre de Wien sélectionne une \textbf{vitesse}, pas une masse ni une charge. Ne réécris jamais la condition sous la forme $v = E/(qB)$ ou $v = mE/(qB)$ : ce serait faux.
\end{attention}

\begin{methode}[Orienter correctement les champs et les forces]
Face à un exercice où l'on demande le sens de $\vec{B}$ (ou de $\vec{E}$) pour que le filtre fonctionne, procède par étapes :
\begin{enumerate}
\item \textbf{Orienter la force électrique d'abord.} Le sens de $\vec{F}_e = q\vec{E}$ dépend du signe de $q$ : même sens que $\vec{E}$ si $q>0$, sens opposé si $q<0$.
\item \textbf{Imposer la force magnétique opposée.} Pour que la trajectoire soit rectiligne, $\vec{F}_m$ doit être exactement opposée à $\vec{F}_e$.
\item \textbf{En déduire le sens de $\vec{B}$} par la règle du produit vectoriel $\vec{F}_m = q\,\vec{v}\wedge\vec{B}$ (règle des trois doigts de la main droite, en tenant compte du signe de $q$).
\item \textbf{Vérifier} que $\vec{v}$, $\vec{E}$ et $\vec{B}$ sont bien deux à deux perpendiculaires, et conclure par la condition $v = E/B$.
\end{enumerate}
\end{methode}

\courssub{Exemple chiffré et rôle dans le spectrographe de masse}

\begin{exemplebox}[Vitesse sélectionnée par un filtre de Wien]
Un filtre de Wien est réalisé avec un champ électrique d'intensité $E = \num{1,0e4}\ \mathrm{V\,m^{-1}}$ (par exemple $U = 100\ \mathrm{V}$ entre deux plaques distantes de $d = 1{,}0\ \mathrm{cm}$) et un champ magnétique $B = \num{2,0e-2}\ \mathrm{T}$, perpendiculaires entre eux et à la vitesse des particules.

La vitesse sélectionnée est :
\[ v = \dfrac{E}{B} = \dfrac{\num{1,0e4}}{\num{2,0e-2}} = \num{5,0e5}\ \mathrm{m\,s^{-1}}. \]
Seules les particules animées de la vitesse $v = 5{,}0\times10^{5}\ \mathrm{m\,s^{-1}}$ traversent le filtre en ligne droite et franchissent la fente de sortie. Toutes les autres — plus lentes ou plus rapides — sont déviées et interceptées. Remarque : cette vitesse vaut environ $0{,}17\,\%$ de la célérité de la lumière ; l'approximation non relativiste reste donc excellente.
\end{exemplebox}

\textbf{Pourquoi ce filtre est-il indispensable en amont d'un spectrographe de masse ?} Rappelons l'équation du spectrographe : dans la chambre de déviation, le rayon de la trajectoire circulaire est $R = \dfrac{mv}{|q|B_0}$, où $B_0$ est le champ du spectrographe. Si le faisceau d'ions issu de la source possède une dispersion de vitesses, des ions de même rapport $m/q$ mais de vitesses différentes décrivent des cercles de rayons différents : les raies s'élargissent et se chevauchent, et la mesure des masses devient impossible. En plaçant un filtre de Wien entre la source et la chambre de déviation, on obtient un faisceau \cle{monocinétique} : tous les ions qui pénètrent dans le spectrographe ont la même vitesse $v = E/B$. Alors, à $v$ fixée, le rayon $R$ ne dépend plus que du rapport $m/q$ :
\[ R = \dfrac{m}{|q|}\times\dfrac{E}{B\,B_0}, \]
et chaque isotope donne une raie fine et bien séparée. C'est exactement le même rôle qu'un canon à électrons suivi d'une tension accélératrice bien contrôlée, mais le filtre de Wien présente un avantage décisif : il sélectionne la vitesse \emph{sans modifier l'énergie} du faisceau (les particules sélectionnées ne sont ni accélérées ni freinées, puisque les forces se compensent).

\begin{savaistu}[Le filtre de Wien dans les microscopes électroniques]
Le filtre de Wien, inventé par le physicien allemand Wilhelm Wien en 1898 (prix Nobel de physique en 1911), équipe aujourd'hui les \textbf{microscopes électroniques} les plus performants. Pour obtenir des images d'une résolution nanométrique — par exemple pour observer des virus ou des défauts cristallins dans les matériaux —, il faut un faisceau d'électrons de vitesse parfaitement homogène : toute dispersion de vitesse provoque un « flou chromatique » analogue à l'aberration chromatique d'une lentille optique. Les filtres de Wien modernes, à géométrie soigneusement compensée, permettent de sélectionner des électrons avec une dispersion d'énergie inférieure à $0{,}1\ \mathrm{eV}$. On les retrouve aussi dans les spectromètres de masse de haute précision utilisés en géochimie et en médecine nucléaire — des techniques dont bénéficient indirectement les laboratoires d'analyse et les hôpitaux, y compris au Cameroun, pour le contrôle qualité des échantillons.
\end{savaistu}

\courssub{Exercice résolu}

\begin{exoresolu}[Réglage d'un filtre de Wien pour des protons]
Une source produit des protons ($q = e = \num{1,6e-19}\ \mathrm{C}$, $m_p = \num{1,67e-27}\ \mathrm{kg}$) de vitesses comprises entre $\num{4,0e5}$ et $\num{6,0e5}\ \mathrm{m\,s^{-1}}$. On veut sélectionner les protons de vitesse $v_0 = \num{5,0e5}\ \mathrm{m\,s^{-1}}$ à l'aide d'un filtre de Wien dont le champ magnétique vaut $B = \num{2,5e-2}\ \mathrm{T}$.
\begin{enumerate}
\item Calculer la valeur du champ électrique $E$ à établir entre les plaques.
\item Les plaques du condensateur sont distantes de $d = 2{,}0\ \mathrm{cm}$. Quelle tension $U$ faut-il appliquer ?
\item Un proton de vitesse $v = \num{6,0e5}\ \mathrm{m\,s^{-1}}$ pénètre dans le filtre. Que se passe-t-il ? De quel côté est-il dévié si $\vec{F}_e$ est dirigée vers le bas ?
\end{enumerate}
\tcblower
\textbf{1. Valeur du champ électrique.} La condition de non-déviation s'écrit $v_0 = E/B$, d'où :
\[ E = v_0\,B = \num{5,0e5}\times\num{2,5e-2} = \num{1,25e4}\ \mathrm{V\,m^{-1}}. \]
\textbf{2. Tension entre les plaques.} Pour un champ uniforme entre plaques planes parallèles, $E = U/d$, donc :
\[ U = E\,d = \num{1,25e4}\times\num{2,0e-2} = 250\ \mathrm{V}. \]
\textbf{3. Sort du proton trop rapide.} Pour ce proton, la force magnétique $F_m = evB$ est \emph{plus grande} que la force électrique $F_e = eE$ car $v > E/B$. La force magnétique l'emporte : le proton est dévié du côté de $\vec{F}_m$, c'est-à-dire \textbf{vers le haut} (opposé à $\vec{F}_e$), et il heurte la paroi ou la plaque supérieure avant la fente de sortie. Il est donc éliminé du faisceau, comme prévu : seuls les protons de vitesse $\num{5,0e5}\ \mathrm{m\,s^{-1}}$ traversent.
\end{exoresolu}

\courssub{Synthèse de la leçon : actions comparées des champs $\vec{E}$ et $\vec{B}$}

Le filtre de Wien illustre à merveille la complémentarité des deux champs. Refermons cette leçon par un tableau comparatif des actions d'un champ électrique uniforme et d'un champ magnétique uniforme sur une particule chargée : c'est une synthèse précieuse pour le baccalauréat.

\begin{popfigure}
\centering
\begin{tabularx}{\linewidth}{|p{3.2cm}|X|X|}
\hline
\rowcolor{popblueL} \textbf{Critère} & \textbf{Champ électrique} $\vec{E}$ & \textbf{Champ magnétique} $\vec{B}$ \\
\hline
Force & $\vec{F}_e = q\vec{E}$, parallèle à $\vec{E}$, indépendante de $\vec{v}$ & $\vec{F}_m = q\,\vec{v}\wedge\vec{B}$, perpendiculaire à $\vec{v}$, nulle si $v=0$ \\
\hline
Travail & $W_{AB}(\vec{F}_e) = q\,U_{AB} \neq 0$ : la force électrique \textbf{travaille} & $W(\vec{F}_m) = 0$ en permanence : la force magnétique \textbf{ne travaille pas} \\
\hline
Énergie cinétique & Modifiée : $\vec{E}$ peut accélérer ou freiner ($\tfrac12 mv^2 = |q|U$) & Inchangée : la norme de $\vec{v}$ est conservée \\
\hline
Trajectoire ($\vec{v}_0\perp$ champ) & \textbf{Parabole} (déviation, oscilloscope) & \textbf{Cercle} de rayon $R = \dfrac{mv}{|q|B}$, parcouru uniformément \\
\hline
Effet sur la direction & Dévie la particule dans le plan contenant $\vec{E}$ & Courbe la trajectoire sans changer la vitesse \\
\hline
Applications & Canon à électrons, oscilloscope, accélérateurs & Spectrographe de masse, cyclotron, filtre de Wien (avec $\vec{E}$) \\
\hline
\end{tabularx}
\legende{Tableau comparatif des actions d'un champ électrique uniforme et d'un champ magnétique uniforme sur une particule chargée. Idée directrice : $\vec{E}$ travaille et modifie l'énergie cinétique ; $\vec{B}$ courbe la trajectoire sans jamais travailler.}
\end{popfigure}

\begin{aretenir}[L'essentiel du filtre de Wien]
\begin{itemize}
\item Dans une zone de \textbf{champs croisés} ($\vec{v}\perp\vec{E}$, $\vec{v}\perp\vec{B}$, $\vec{E}\perp\vec{B}$), la particule subit $\vec{F}_e = q\vec{E}$ et $\vec{F}_m = q\,\vec{v}\wedge\vec{B}$, orientées en sens opposés.
\item La trajectoire reste \textbf{rectiligne et uniforme} si et seulement si les forces se compensent, ce qui impose $\boxed{v = E/B}$.
\item Cette condition est \textbf{indépendante de la charge et de la masse} : le filtre de Wien sélectionne une vitesse, rien d'autre.
\item Placé en amont d'un spectrographe de masse, il fournit un faisceau \textbf{monocinétique}, condition d'une mesure précise des masses.
\item Bilan de la leçon : le champ électrique \emph{accélère, dévie et travaille} ; le champ magnétique \emph{courbe la trajectoire sans travailler}. Combinés, ils permettent de trier, guider et analyser les particules chargées — du tube cathodique de l'oscilloscope au cyclotron des hôpitaux.
\end{itemize}
\end{aretenir}

\newpage

\coursec{Activité d'intégration}

\coursec{Mouvement de particules chargées dans les champs électrique et magnétiques uniformes}

\courssub{Accélération d'une particule chargée par un champ électrique uniforme — Canon à électrons}

\begin{definition}[Champ électrique uniforme et force électrique]
Un champ électrique $\vec{E}$ est dit \textbf{uniforme} s'il a la même valeur (norme, direction, sens) en tout point de l'espace. Une particule de charge $q$ placée dans ce champ subit une force électrique constante :
\[ \vec{F}_e = q\,\vec{E} \]
Le travail de cette force lors d'un déplacement du point $A$ au point $C$ ne dépend que de la différence de potentiel $U = V_A - V_C$ :
\[ W_{A\to C}(\vec{F}_e) = q\,U \]
\end{definition}

\begin{propriete}[Théorème de l'énergie cinétique — Vitesse acquise]
Une particule de charge $q$ et de masse $m$, partie du repos (ou de vitesse négligeable) et accélérée par une tension $U$, acquiert une énergie cinétique égale au travail de la force électrique :
\[ \frac{1}{2}mv^2 = |q|\,U \quad\Rightarrow\quad \boxed{v = \sqrt{\frac{2|q|U}{m}}} \]
\emph{Condition d'application :} champ électrique uniforme, vitesse initiale nulle, effets relativistes négligeables ($v \ll c$).
\end{propriete}

\begin{methode}[Calcul de la vitesse à la sortie d'un canon à particules]
\begin{enumerate}
\item Identifier la charge $q$ (souvent $\pm e$), la masse $m$ et la tension d'accélération $U$.
\item Vérifier que $v \ll c$ (si $v > 0,1c \approx \SI{3e7}{m/s}$, passer en relativiste).
\item Appliquer $v = \sqrt{2|q|U/m}$.
\item Vérifier l'homogénéité : $\left[\sqrt{\dfrac{\mathrm{C}\cdot\mathrm{V}}{\mathrm{kg}}}\right] = \sqrt{\dfrac{\mathrm{J}}{\mathrm{kg}}} = \mathrm{m\,s^{-1}}$.
\end{enumerate}
\end{methode}

\begin{exemplebox}[Canon à électrons d'un oscilloscope]
Un canon à électrons accélère des électrons ($m_e = \SI{9,11e-31}{kg}$, $q=-e$) sous une tension $U = \SI{1,00}{kV}$.
\[ v = \sqrt{\frac{2\times\SI{1,60e-19}{}\times\SI{1,00e3}{}}{\SI{9,11e-31}{}}} = \SI{1,87e7}{m/s} \approx 0,062\,c \]
La vitesse reste non relativiste. Ces électrons forment un faisceau fin dirigé vers les plaques de déviation.
\end{exemplebox}

\courssub{Déviation électrique : trajectoire parabolique et oscilloscope}

\begin{propriete}[Mouvement entre deux plaques planes parallèles — Déviation électrique]
Un faisceau de particules de charge $q$, masse $m$, vitesse initiale $\vec{v}_0$ parallèle aux plaques, traverse un champ électrique uniforme $\vec{E} = \dfrac{U}{d}\,\vec{u}_y$ (perpendiculaire à $\vec{v}_0$) sur une longueur $L$.
\begin{itemize}
\item \textbf{Équations horaires} (origine à l'entrée, $x$ // $\vec{v}_0$, $y$ // $\vec{E}$) :
  \[ x(t) = v_0 t \quad;\quad y(t) = \frac{1}{2}\frac{|q|E}{m}t^2 = \frac{|q|U}{2md}\,t^2 \]
\item \textbf{Trajectoire} : en éliminant $t$, $y(x) = \dfrac{|q|U}{2md v_0^2}\,x^2$ \textbf{parabole} concave vers la plaque de signe opposé à $q$.
\item \textbf{Déflexion à la sortie des plaques} ($x=L$) :
  \[ y_L = \frac{|q|U L^2}{2md v_0^2} \quad;\quad \tan\theta = \frac{v_y}{v_x} = \frac{|q|U L}{md v_0^2} \]
\item \textbf{Sur l'écran} (distance $D$ après les plaques) : $y_{\text{écran}} = y_L + D\tan\theta$.
\end{itemize}
\end{propriete}

\begin{exemplebox}[Déflexion sur l'écran d'un oscilloscope]
Électrons ($v_0 = \SI{1,87e7}{m/s}$), plaques $L=\SI{2,0}{cm}$, $d=\SI{1,0}{cm}$, $U=\SI{50}{V}$, écran à $D=\SI{15}{cm}$.
\[ y_L = \frac{\SI{1,60e-19}{}\times\SI{50}{}\times(\SI{0,02}{})^2}{2\times\SI{9,11e-31}{}\times\SI{0,01}{}\times(\SI{1,87e7}{})^2} = \SI{2,5e-3}{m} = \SI{2,5}{mm} \]
\[ \tan\theta = \frac{2y_L}{L} = 0,25 \quad\Rightarrow\quad y_{\text{écran}} = \SI{2,5}{mm} + \SI{0,15}{}\times 0,25 = \SI{40}{mm} \]
La trace se déplace de \SI{4}{cm} sur l'écran : l'oscilloscope visualise la tension $U$.
\end{exemplebox}

\begin{popfigure}
\begin{center}
\begin{tikzpicture}[scale=1.2,>=Stealth]
\draw[fill=popcream] (-0.5,0) rectangle (0,1.5); \node at (-0.25,2) {\small Canon à $e^-$};
\draw[->,thick,poporange] (-0.25,0.75) -- (0,0.75);
\draw[thick,popblue] (0,1.5) -- (3,1.5) node[midway,above] {$+$};
\draw[thick,popred] (0,0) -- (3,0) node[midway,below] {$-$};
\draw[popblue,->] (1.5,1.2) -- (1.5,0.8) node[midway,right] {$\vec{E}$};
\draw[poporange,thick] (0,0.75) .. controls (1.5,1.1) and (3,1.3) .. (3,1.3);
\draw[poporange,thick] (3,1.3) -- (5,1.8);
\draw[popline] (5,-0.5) -- (5,2.5);
\foreach \y in {-0.5,0,0.5,1,1.5,2,2.5} \draw[popmuted] (5,\y) -- ++(0.1,0) node[right] {\tiny \y cm};
\fill[poporange] (5,1.8) circle (0.04);
\node at (2.5,-0.8) {Plaques de déviation ($L,d,U$)};
\node at (5.5,1.8) [right] {Écran ($D$)};
\end{tikzpicture}
\end{center}
\legende{Principe de la déviation électrique dans un oscilloscope : trajectoire parabolique dans les plaques, rectiligne ensuite.}
\end{popfigure}

\courssub{Mouvement dans un champ magnétique uniforme ($\vec{v} \perp \vec{B}$)}

\begin{definition}[Force de Lorentz magnétique]
Dans un champ magnétique $\vec{B}$, une particule de charge $q$ se déplaçant à la vitesse $\vec{v}$ subit la force :
\[ \vec{F}_m = q\,\vec{v}\wedge\vec{B} \]
Cette force est \textbf{toujours perpendiculaire} à $\vec{v}$ et à $\vec{B}$. Son travail est nul : $W(\vec{F}_m)=0$. Elle ne modifie que la \textbf{direction} de $\vec{v}$, pas sa norme.
\end{definition}

\begin{propriete}[Mouvement circulaire uniforme — Rayon et période]
Si $\vec{v} \perp \vec{B}$ et $\vec{E}=0$, la force de Lorentz est centripète. Le mouvement est un \textbf{cercle uniforme} :
\[ |q|vB = \frac{mv^2}{R} \quad\Rightarrow\quad \boxed{R = \frac{mv}{|q|B}} \]
La période $T$ et la pulsation cyclotron $\omega_c$ (indépendantes de $v$ !) sont :
\[ T = \frac{2\pi R}{v} = \frac{2\pi m}{|q|B} \quad;\quad \omega_c = \frac{2\pi}{T} = \frac{|q|B}{m} \]
Le sens de rotation est donné par la règle de la main droite (pour $q>0$) ou gauche (pour $q<0$) sur $\vec{v}\wedge\vec{B}$.
\end{propriete}

\begin{methode}[Analyse du mouvement dans $\vec{B}$ uniforme ($\vec{v}\perp\vec{B}$)]
\begin{enumerate}
\item Dessiner $\vec{B}$ (souvent $\odot$ sortant ou $\otimes$ entrant) et $\vec{v}_0$.
\item Déterminer le sens de $\vec{F}_m = q\vec{v}_0\wedge\vec{B}$ (règle des 3 doigts).
\item Écrire $R = mv/(|q|B)$ et $T = 2\pi m/(|q|B)$.
\item Vérifier que la particule reste dans la zone de champ (comparer $R$ aux dimensions de l'appareil).
\end{enumerate}
\end{methode}

\begin{exemplebox}[Séparation isotopique dans un champ magnétique]
Ions \ce{Cl+} ($q=+e$) de vitesse $v=\SI{1,5e5}{m/s}$ dans $B=\SI{0,40}{T}$.
$m(^{35}\text{Cl}) = \SI{5,81e-26}{kg}$, $m(^{37}\text{Cl}) = \SI{6,14e-26}{kg}$.
\[ R_{35} = \frac{\SI{5,81e-26}{}\times\SI{1,5e5}{}}{\SI{1,60e-19}{}\times\SI{0,40}{}} = \SI{13,6}{cm} \]
\[ R_{37} = \frac{\SI{6,14e-26}{}\times\SI{1,5e5}{}}{\SI{1,60e-19}{}\times\SI{0,40}{}} = \SI{14,4}{cm} \]
L'écart de \SI{8}{mm} sur le diamètre permet de distinguer les isotopes.
\end{exemplebox}

\courssub{Applications : Spectrographe de masse et Cyclotron}

\begin{propriete}[Spectrographe de masse (type Bainbridge)]
L'appareil combine un \textbf{filtre de vitesse} (champs $\vec{E}$ et $\vec{B}_1$ croisés) et une \textbf{chambre de déviation} (champ $\vec{B}_2$).
\begin{enumerate}
\item \textbf{Sélecteur de vitesse} : seules les particules de vitesse $v = E/B_1$ passent sans déviation.
\item \textbf{Déviation magnétique} : dans $\vec{B}_2$, rayon $R = \dfrac{mv}{qB_2} = \dfrac{mE}{qB_1B_2}$.
\item \textbf{Masse mesurée} : $m = \dfrac{qB_1B_2}{E}\,R$. À $E,B_1,B_2$ fixés, $R \propto m$.
\end{enumerate}
La résolution $\dfrac{\Delta m}{m} = \dfrac{\Delta R}{R}$ dépend de la précision sur $R$.
\end{propriete}

\begin{propriete}[Cyclotron — Accélérateur de particules]
Deux électrodes creuses (\og Dees \fg) sous champ $\vec{B}$ uniforme. Une tension alternative $U(t)=U_0\sin(\omega t)$ est appliquée entre les Dees.
\begin{itemize}
\item \textbf{Condition de résonance} : la période de la tension doit égaler la période de révolution : $\omega = \omega_c = \dfrac{qB}{m}$.
\item À chaque passage dans l'interstice, la particule gagne une énergie $qU_0$. Le rayon augmente : $R_n = \dfrac{m v_n}{qB}$.
\item \textbf{Énergie cinétique maximale} (rayon max $R_{\text{max}}$ imposé par la taille de l'appareil) :
  \[ E_{c,\text{max}} = \frac{1}{2}m v_{\text{max}}^2 = \frac{q^2 B^2 R_{\text{max}}^2}{2m} \]
\end{itemize}
\emph{Limite :} relativiste si $v \to c$ ($\omega_c$ diminue, désaccord). Solution : synchrocyclotron (fréquence modulée) ou synchrotron.
\end{propriete}

\begin{exemplebox}[Cyclotron de radiothérapie (contexte camerounais)]
Un cyclotron à protons ($q=+e$, $m_p=\SI{1,67e-27}{kg}$) pour la radiothérapie (hôpital de Yaoundé) : $B=\SI{2,5}{T}$, $R_{\text{max}}=\SI{0,60}{m}$.
\[ E_{c,\text{max}} = \frac{(\SI{1,60e-19}{})^2 \times (\SI{2,5}{})^2 \times (\SI{0,60}{})^2}{2 \times \SI{1,67e-27}{}} = \SI{1,73e-11}{J} \approx \SI{108}{MeV} \]
Fréquence RF : $f = \dfrac{qB}{2\pi m} = \dfrac{\SI{1,60e-19}{}\times\SI{2,5}{}}{2\pi\times\SI{1,67e-27}{}} = \SI{38,1}{MHz}$.
Ces protons traitent des tumeurs oculaires ou cérébrales par balayage de faisceau (pencil beam scanning).
\end{exemplebox}

\begin{popfigure}
\begin{center}
\begin{tikzpicture}[scale=0.9,>=Stealth]
% Cyclotron schematic
\draw[fill=popblue!10] (-2.5,-2.5) rectangle (2.5,2.5);
\foreach \angle in {0,30,...,330} \draw[popblue,->] (\angle:2.2) -- ++(\angle:0.3) node[pos=1.2,popblue] {\tiny $\otimes$};
\node at (0,2.8) {Champ $\vec{B}$ entrant ($\otimes$)};
% Dees
\draw[thick,popdark] (0,0) ellipse (1.8 and 1.2);
\draw[thick,popdark] (0,0) ellipse (1.8 and -1.2); % just border
\draw[thick,popdark] (-1.8,0) -- (1.8,0); % Gap
\node at (-2.2,0) [left] {Dee 1};
\node at (2.2,0) [right] {Dee 2};
% Spiral trajectory
\draw[poporange,thick] (0,0) 
  foreach \r/\a in {0.1/0, 0.3/180, 0.5/360, 0.7/540, 0.9/720, 1.1/900, 1.3/1080, 1.5/1260, 1.7/1440} {
    .. controls ++(\a+45:0.15) and ++(\a+135:0.15) .. (\a+180:\r)
  };
\node[poporange] at (1.5,1.5) {Trajectoire spirale};
% RF source
\draw[popred,thick] (-2.5,0) -- (-3,0) node[left] {$U(t)=U_0\sin(\omega t)$};
\draw[popred,thick] (2.5,0) -- (3,0);
\end{tikzpicture}
\end{center}
\legende{Principe du cyclotron : les particules spiralent sous l'effet de $\vec{B}$ et sont accélérées à chaque passage dans l'interstice par la tension alternative de fréquence $\omega = qB/m$.}
\end{popfigure}

\courssub{Filtre de Wien : champs électrique et magnétique croisés}

\begin{propriete}[Filtre de Wien — Sélecteur de vitesse]
Une particule traverse une zone où $\vec{E}$, $\vec{B}$ et $\vec{v}$ sont deux à deux perpendiculaires ($\vec{E} \perp \vec{B}$, $\vec{v} \perp \vec{E}$, $\vec{v} \perp \vec{B}$).
Les forces électrique $\vec{F}_e = q\vec{E}$ et magnétique $\vec{F}_m = q\vec{v}\wedge\vec{B}$ sont \textbf{opposées}.
\begin{itemize}
\item \textbf{Condition de traversée non déviée} : $\vec{F}_e + \vec{F}_m = \vec{0}$ $\Rightarrow$ $qE = qvB$ $\Rightarrow$ \boxed{v = \dfrac{E}{B}}.
\item Cette vitesse ne dépend \textbf{ni de la masse $m$, ni de la charge $q$} (signe de $q$ détermine seulement le sens des forces).
\item Si $v > E/B$ : déviation magnétique domine. Si $v < E/B$ : déviation électrique domine.
\end{itemize}
\end{propriete}

\begin{methode}[Dimensionnement d'un filtre de Wien]
\begin{enumerate}
\item Choisir la vitesse $v$ à sélectionner (ex: sortie canon à électrons).
\item Choisir une valeur réaliste pour $B$ (aimant permanent ou électroaimant, typ. \SI{0,1}{T} à \SI{1}{T}).
\item Calculer $E = vB$.
\item Si $E$ est créé par deux plaques planes distance $d$ sous tension $U'$ : $U' = E\,d = vBd$.
\item Vérifier que $U'$ est compatible avec l'isolement (pas d'effet de pointe, claquage air $\approx \SI{3e6}{V/m}$).
\end{enumerate}
\end{methode}

\begin{exemplebox}[Filtre de Wien en amont d'un spectrographe]
On veut sélectionner $v = \SI{1,5e5}{m/s}$ avec $B' = \SI{0,10}{T}$.
$E = vB' = \SI{1,5e4}{V/m}$. Pour $d = \SI{5,0}{cm}$, $U' = \SI{750}{V}$. Tension aisément réalisable.
\end{exemplebox}

\begin{popfigure}
\begin{center}
\begin{tikzpicture}[scale=1.1,>=Stealth]
\draw[fill=popgreen!10] (0,0) rectangle (4,2);
\draw[popblue,thick,->] (0.5,1) -- (3.5,1) node[midway,above] {$\vec{v}$};
\draw[popred,thick,->] (2,1.5) -- (2,0.5) node[midway,right] {$\vec{E}$};
\foreach \x in {0.5,1,1.5,2,2.5,3,3.5} \draw[popblue] (\x,1.8) -- ++(0,-0.2) node[above,popblue] {$\odot$};
\node at (2,2.2) {$\vec{B}$ sortant ($\odot$)};
\draw[poporange,thick,->] (2,1) -- (2,1.5) node[above] {$\vec{F}_e$};
\draw[poppurple,thick,->] (2,1) -- (2,0.5) node[below] {$\vec{F}_m$};
\node at (2,-0.5) {$\vec{F}_e + \vec{F}_m = \vec{0}$ si $v=E/B$};
\end{tikzpicture}
\end{center}
\legende{Filtre de Wien : forces électrique et magnétique opposées. Seule la vitesse $v=E/B$ passe sans déviation.}
\end{popfigure}

\courssub{Activité d'intégration : Conception d'un analyseur d'isotopes du chlore}

\courssub{Objectifs de l'activité}
À l'issue de cette activité d'intégration, tu seras capable de :
\begin{itemize}
\item mobiliser le \cle{théorème de l'énergie cinétique} pour déterminer la vitesse acquise par une particule chargée accélérée par une tension $U$ ;
\item exploiter l'expression du \cle{rayon de la trajectoire circulaire} $R = \dfrac{mv}{|q|B}$ d'une particule dans un champ magnétique uniforme ;
\item combiner accélération électrique et déviation magnétique pour \cle{dimensionner un spectrographe de masse} (méthode de Bainbridge simplifiée) ;
\item exploiter le principe du \cle{filtre de Wien} (champs $\vec{E}$ et $\vec{B}$ croisés) comme sélecteur de vitesse ;
\item travailler en équipe, rédiger une note de calcul argumentée et produire un schéma coté, comme le ferait un ingénieur.
\end{itemize}

\courssub{Situation}
Le laboratoire de physico-chimie d'une grande entreprise agro-industrielle de la région du Littoral souhaite contrôler la composition isotopique du chlore présent dans des échantillons d'eau de forage destinés à l'irrigation. Le chlore naturel est un mélange de deux isotopes stables : \ce{^{35}_{17}Cl} (environ 75 \%) et \ce{^{37}_{17}Cl} (environ 25 \%). Pour les distinguer, l'entreprise fait appel à ton groupe d'élèves de Terminale C, transformé pour l'occasion en bureau d'études : il s'agit de \textbf{concevoir un analyseur d'isotopes}, c'est-à-dire un mini-spectrographe de masse.

\textbf{Principe de l'appareil.} Les atomes de chlore sont d'abord ionisés en ions \ce{Cl+} de charge $q = +e$. Les ions, supposés émis avec une vitesse initiale négligeable, traversent une \textbf{chambre d'accélération} où règne une tension $U = \SI{4,0}{kV}$. Ils pénètrent ensuite dans une \textbf{chambre de déviation} où règne un champ magnétique uniforme $\vec{B}$, perpendiculaire à leur vitesse, de valeur $B = \SI{0,40}{T}$. Chaque ion décrit un demi-cercle puis frappe une \textbf{plaque de détection} photographique de longueur $L = \SI{10}{cm}$. Les deux isotopes, de masses différentes, doivent laisser deux impacts distincts.

\textbf{Données numériques :}
\begin{itemize}
\item charge élémentaire : $e = \SI{1,60 \times 10^{-19}}{C}$ ;
\item unité de masse atomique : $\SI{1}{u} = \SI{1,66 \times 10^{-27}}{kg}$ ;
\item masses des isotopes : $m_1 = \SI{35}{u}$ pour \ce{^{35}_{17}Cl} et $m_2 = \SI{37}{u}$ pour \ce{^{37}_{17}Cl} ;
\item tension d'accélération : $U = \SI{4,0 \times 10^{3}}{V}$ ; champ magnétique : $B = \SI{0,40}{T}$ ;
\item longueur de la plaque de détection : $L = \SI{10}{cm}$.
\end{itemize}

\textbf{Livrables attendus :} une note de calcul complète et un schéma coté de l'analyseur, présentés au « client » (le reste de la classe) en cinq minutes.

\courssub{Consignes}
Par groupes de quatre, répondez aux tâches suivantes, dans l'ordre. Chaque résultat numérique sera donné avec son unité SI et trois chiffres significatifs.

\begin{enumerate}
\item \textbf{Accélération.} En appliquant le théorème de l'énergie cinétique à un ion entre les deux électrodes de la chambre d'accélération, établir l'expression littérale de la vitesse $v$ d'un ion à la sortie, en fonction de $e$, $U$ et $m$. Calculer $v_1$ et $v_2$ pour les deux isotopes. Lequel est le plus rapide ? Justifier sans calcul supplémentaire.
\item \textbf{Déviation magnétique.} Montrer que, dans la chambre de déviation, le mouvement de chaque ion est circulaire uniforme et établir l'expression du rayon $R = \dfrac{mv}{eB}$. Calculer les rayons $R_1$ et $R_2$ des deux demi-cercles.
\item \textbf{Séparation.} Les ions entrent dans la chambre de déviation au même point $O$ et la plaque de détection est placée perpendiculairement à la direction d'injection, de sorte que chaque isotope frappe la plaque à une distance $2R$ de $O$. Calculer l'écart $\Delta d$ entre les deux points d'impact. La plaque de $\SI{10}{cm}$ suffit-elle à recueillir les deux impacts ?
\item \textbf{Optimisation.} Le client souhaite \textbf{doubler l'écart} $\Delta d$ pour faciliter la lecture. En exprimant $\Delta d$ en fonction de $U$, $B$, $e$ et des masses, proposer \emph{deux} réglages indépendants (l'un portant sur $U$, l'autre sur $B$) permettant d'obtenir $\Delta d' = 2\,\Delta d$, et calculer les nouvelles valeurs.
\item \textbf{Bonus — filtre de Wien.} En amont de l'accélérateur, on veut sélectionner uniquement les ions de vitesse $v_1$ (isotope 35) à l'aide d'un filtre de vitesse : champ électrique $\vec{E}$ et champ magnétique $\vec{B'}$ uniformes, croisés et perpendiculaires à la vitesse. Établir la condition de traversée sans déviation, puis calculer $E$ pour $B' = \SI{0,10}{T}$. Quelle tension $U'$ faut-il appliquer entre deux plaques planes distantes de $d = \SI{5,0}{cm}$ pour créer ce champ ?
\end{enumerate}

\courssub{Ressources à mobiliser}

\begin{aretenir}[Boîte à outils de l'ingénieur]
\begin{itemize}
\item \textbf{Théorème de l'énergie cinétique} entre deux points $A$ et $C$ : $\Delta E_c = W_{A\to C}(\vec{F})$, avec le travail de la force électrique $W = |q|\,U$ pour une charge accélérée sous la tension $U$.
\item \textbf{Vitesse acquise} (particule initialement au repos) : $\dfrac{1}{2}mv^2 = |q|U \;\Rightarrow\; v = \sqrt{\dfrac{2|q|U}{m}}$.
\item \textbf{Champ magnétique uniforme}, $\vec{v} \perp \vec{B}$ : la force de Lorentz $\vec{F} = q\,\vec{v}\wedge\vec{B}$ est centripète ; le mouvement est circulaire uniforme de rayon $R = \dfrac{mv}{|q|B}$ et de période $T = \dfrac{2\pi m}{|q|B}$.
\item \textbf{Spectrographe de masse} : en combinant les deux résultats, $R = \dfrac{1}{B}\sqrt{\dfrac{2mU}{|q|}}$ : à $U$ et $B$ fixés, le rayon ne dépend que de la masse.
\item \textbf{Filtre de Wien} : traversée rectiligne non déviée si $eE = evB'$, soit $v = \dfrac{E}{B'}$ ; champ entre plaques planes : $E = \dfrac{U'}{d}$.
\end{itemize}
\end{aretenir}

\courssub{Démarche et corrigé commenté}

\begin{exoresolu}[Note de calcul — Analyseur d'isotopes du chlore]
Énoncé : dimensionner le spectrographe de masse décrit dans la situation (tâches 1 à 5).
\tcblower
\textbf{Étape 1 — Vitesse après accélération.}

Entre la source (potentiel $V_A$) et la sortie du canon à ions (potentiel $V_C$), seule la force électrique travaille. Le théorème de l'énergie cinétique appliqué à un ion \ce{Cl+} s'écrit :
\[ \frac{1}{2}mv^2 - 0 = eU \quad\Rightarrow\quad \boxed{v = \sqrt{\frac{2eU}{m}}} \]
Vérification dimensionnelle : $\left[\sqrt{\dfrac{\mathrm{C}\cdot\mathrm{V}}{\mathrm{kg}}}\right] = \sqrt{\dfrac{\mathrm{J}}{\mathrm{kg}}} = \sqrt{\mathrm{m^2\,s^{-2}}} = \mathrm{m\,s^{-1}}$. Correct.

Masses : $m_1 = 35\times\SI{1,66 \times 10^{-27}}{} = \SI{5,81 \times 10^{-26}}{kg}$ et $m_2 = 37\times\SI{1,66 \times 10^{-27}}{} = \SI{6,14 \times 10^{-26}}{kg}$.
\begin{align*}
v_1 &= \sqrt{\frac{2\times\SI{1,60 \times 10^{-19}}{}\times\SI{4,0 \times 10^{3}}{}}{\SI{5,81 \times 10^{-26}}{}}} = \SI{1,48 \times 10^{5}}{m/s} \\
v_2 &= \sqrt{\frac{2\times\SI{1,60 \times 10^{-19}}{}\times\SI{4,0 \times 10^{3}}{}}{\SI{6,14 \times 10^{-26}}{}}} = \SI{1,44 \times 10^{5}}{m/s}
\end{align*}
L'isotope léger \ce{^{35}Cl} est le plus rapide : à énergie cinétique $eU$ identique, $v$ varie en $1/\sqrt{m}$.

\textbf{Étape 2 — Rayons des trajectoires.}

Dans la chambre de déviation, $\vec{v}\perp\vec{B}$ : la force de Lorentz, perpendiculaire à la vitesse, ne travaille pas ; le mouvement est circulaire uniforme. La deuxième loi de Newton projetée sur la normale donne $evB = \dfrac{mv^2}{R}$, d'où $R = \dfrac{mv}{eB}$.
\begin{align*}
R_1 &= \frac{\SI{5,81 \times 10^{-26}}{}\times\SI{1,48 \times 10^{5}}{}}{\SI{1,60 \times 10^{-19}}{}\times\SI{0,40}{}} = \SI{0,135}{m} = \SI{13,5}{cm} \\
R_2 &= \frac{\SI{6,14 \times 10^{-26}}{}\times\SI{1,44 \times 10^{5}}{}}{\SI{1,60 \times 10^{-19}}{}\times\SI{0,40}{}} = \SI{0,139}{m} = \SI{13,9}{cm}
\end{align*}

\textbf{Étape 3 — Écart entre les impacts.}

Chaque ion décrit un demi-cercle : il frappe la plaque à la distance $2R$ du point d'injection $O$. L'écart entre les deux impacts est :
\[ \Delta d = 2R_2 - 2R_1 = 2\times(\SI{13,9}{cm} - \SI{13,5}{cm}) = \SI{7,8 \times 10^{-3}}{m} = \SI{7,8}{mm}. \]
Les deux impacts sont situés à $2R_1 = \SI{27,0}{cm}$ et $2R_2 = \SI{27,8}{cm}$ de $O$ : ils tiennent largement sur la plaque de $\SI{10}{cm}$ (à condition de la positionner correctement, centrée vers \SI{27,4}{cm}). L'écart de $\SI{7,8}{mm}$ est nettement supérieur à la largeur des traces (fraction de millimètre) : la séparation est \textbf{réussie}.

\textbf{Étape 4 — Doubler l'écart.}

En injectant $v = \sqrt{2eU/m}$ dans $R$ : $R = \dfrac{1}{B}\sqrt{\dfrac{2mU}{e}}$, donc
\[ \Delta d = \frac{2}{B}\sqrt{\frac{2U}{e}}\left(\sqrt{m_2}-\sqrt{m_1}\right) \;\propto\; \frac{\sqrt{U}}{B}. \]
Deux solutions indépendantes :
\begin{itemize}
\item multiplier $U$ par $4$ : $U' = \SI{16}{kV}$ (car $\sqrt{4}=2$) ;
\item diviser $B$ par $2$ : $B'' = \SI{0,20}{T}$.
\end{itemize}
Dans les deux cas, $\Delta d' = 2\,\Delta d \approx \SI{15,6}{mm}$. Le réglage sur $B$ est techniquement le plus simple (on agit sur le courant de l'électroaimant).

\textbf{Étape 5 — Filtre de Wien.}

Dans le filtre, l'ion subit $\vec{F}_e = e\vec{E}$ et $\vec{F}_m = e\,\vec{v}\wedge\vec{B'}$ en sens opposés. La traversée rectiligne impose $eE = ev_1B'$, soit $v_1 = \dfrac{E}{B'}$, indépendamment de la masse et de la charge : seule la vitesse est filtrée.
\[ E = v_1 B' = \SI{1,48 \times 10^{5}}{}\times\SI{0,10}{} = \SI{1,48 \times 10^{4}}{V/m}. \]
Entre deux plaques distantes de $d = \SI{5,0 \times 10^{-2}}{m}$ : $U' = E\,d = \SI{1,48 \times 10^{4}}{}\times\SI{5,0 \times 10^{-2}}{} = \SI{7,4 \times 10^{2}}{V}$, soit environ $\SI{740}{V}$ : tension tout à fait réalisable en laboratoire.

\textbf{Schéma coté (restitution).}
\begin{popfigure}
\begin{center}
\begin{tikzpicture}[scale=0.85,>=Stealth]
% Acceleration chamber
\draw[fill=popcream] (0,0) rectangle (2,2.4);
\node[align=center] at (1,2.7) {\small Accélération $U=\SI{4,0}{kV}$};
\draw[->,thick,poporange] (0.3,1.2) -- (1.7,1.2);
\fill[poppurple] (2,1.2) circle (0.06) node[above left]{\small $O$};
% Deflection chamber
\draw[fill=popblue!10] (2,-0.2) rectangle (7.6,-3.4);
\node at (4.8,-3.7) {\small Chambre de déviation : $\vec{B}$ sortant, $B=\SI{0,40}{T}$};
\foreach \x in {2.5,3.5,4.5,5.5,6.5,7.3} \foreach \y in {-0.6,-1.4,-2.2,-3.0} \node[popblue] at (\x,\y) {\scriptsize $\odot$};
% Trajectories
\draw[popink,thick] (2,1.2) -- (2,-1.35);
\draw[popink,thick] (2,-1.35) arc (90:-90:1.35);
\draw[popgreen,thick] (2,1.2) -- (2,-1.39);
\draw[popgreen,thick] (2,-1.39) arc (90:-90:1.39);
% Detector plate
\draw[thick,popdark] (2,-0.4) -- (2,-4.3);
\fill[popink] (2,-4.05) circle (0.05); \fill[popgreen] (2,-4.18) circle (0.05);
\node[right,align=left] at (2.2,-4.35) {\small Impacts séparés de $\Delta d \approx \SI{7,8}{mm}$};
\node[popink] at (3.1,-2.4) {\small $R_1=\SI{13,5}{cm}$};
\node[popgreen] at (3.2,-3.0) {\small $R_2=\SI{13,9}{cm}$};
\end{tikzpicture}
\end{center}
\legende{Schéma de principe coté de l'analyseur : accélération puis demi-cercles de rayons $R_1 = \SI{13,5}{cm}$ (\ce{^{35}Cl}) et $R_2 = \SI{13,9}{cm}$ (\ce{^{37}Cl}).}
\end{popfigure}
\end{exoresolu}

\begin{attention}[Pièges rencontrés pendant l'activité]
\begin{itemize}
\item Ne pas oublier de convertir les masses en kilogrammes ($\times\SI{1,66 \times 10^{-27}}{}$) et la tension en volts.
\item L'écart cherché est $2R_2 - 2R_1$ (diamètres), pas $R_2 - R_1$ : erreur classique d'un facteur 2.
\item Dans le filtre de Wien, la condition $v = E/B'$ ne dépend ni de $m$ ni de $q$ : inutile de chercher à faire intervenir la masse.
\item La force de Lorentz magnétique ne travaille jamais : elle modifie la direction de $\vec{v}$, pas sa norme.
\item Vérifier toujours la non-relativité : ici $v \sim \SI{1,5e5}{m/s} \ll c$, approximation newtonienne validée.
\end{itemize}
\end{attention}

\courssub{Bilan de l'activité}

Cette activité t'a placé dans la peau d'un ingénieur confronté à un cahier des charges réel. Elle a permis de mettre en évidence que :

\begin{itemize}
\item l'\textbf{accélération électrique} fixe l'énergie cinétique des ions : $\frac{1}{2}mv^2 = eU$, donc la vitesse dépend de la masse en $1/\sqrt{m}$ ;
\item la \textbf{déviation magnétique} transforme cette différence de vitesse (et de masse) en différence de rayons : $R = \dfrac{1}{B}\sqrt{\dfrac{2mU}{e}}$ ; c'est le principe même de la \textbf{spectrométrie de masse}, utilisée en contrôle qualité, en médecine et en datation ;
\item la résolution de l'appareil se règle facilement : l'écart entre impacts varie comme $\sqrt{U}/B$, ce qui offre deux leviers de réglage indépendants ;
\item le \textbf{filtre de Wien} sélectionne une vitesse unique quelles que soient la masse et la charge, et constitue un complément précieux en amont de l'analyseur ;
\item enfin, la démarche complète — modéliser, calculer, vérifier la faisabilité (plaque de $\SI{10}{cm}$), optimiser, rédiger une note de calcul et un schéma coté — est exactement celle du physicien ou de l'ingénieur, et mobilise l'ensemble des savoir-faire de la leçon.
\end{itemize}

\begin{savaistu}[Et au quotidien ?]
Les spectromètres de masse modernes, descendants directs de ton analyseur, pèsent parfois quelques kilogrammes seulement et équipent les laboratoires d'analyse des eaux et des denrées alimentaires (contrôle qualité SONARA, analyses LNAN). Le même principe physique — champ $\vec{B}$ et trajectoires circulaires — est à l'œuvre dans le cyclotron des centres de radiothérapie, comme ceux que le Cameroun développe pour le traitement des cancers : les particules y sont accélérées en spirale avant d'être dirigées vers la tumeur. La physique des particules chargées n'est pas qu'un chapitre du Bac : c'est un outil au service de la santé et de l'industrie.
\end{savaistu}

\newpage

\coursec{Méthodes et exercices résolus}

\courssub{Méthode 1 — Accélération d'une particule par un champ électrique uniforme (canon à électrons)}

\begin{methode}[Accélération d'une particule chargée]
Pour déterminer la vitesse acquise par une particule chargée accélérée entre deux plaques soumises à une tension $U$ :
\begin{enumerate}
\item \textbf{Identifier} la particule (charge $q$, masse $m$) et la tension accélératrice $U$ (en valeur absolue). Vérifier que la vitesse initiale est nulle ou négligeable, sinon le préciser.
\item \textbf{Appliquer le théorème de l'énergie cinétique} entre la cathode C et l'anode A :
\[ E_{c,A}-E_{c,C}=W_{C\to A}(\vec{F})=q\,U_{CA}=q(V_C-V_A). \]
\item \textbf{Justifier} que le poids de la particule est négligeable devant la force électrique (comparer $mg$ et $|q|E$ : rapport typique de $10^{-14}$ pour un électron).
\item \textbf{Isoler la vitesse} : si $v_C\approx 0$ et si la particule est effectivement accélérée ($qU_{CA}>0$),
\[ v_A=\sqrt{\dfrac{2\,|q|\,U}{m}}. \]
\item \textbf{Convertir} si besoin l'énergie cinétique en électronvolts : $\SI{1}{eV}=\num{1,60e-19}{J}$.
\item \textbf{Vérifier} l'ordre de grandeur : un électron accéléré sous quelques centaines de volts atteint $10^6$ à $10^7\ \mathrm{m\cdot s^{-1}}$.
\end{enumerate}
\end{methode}

\begin{exoresolu}[Canon à électrons d'un oscilloscope]
Dans le canon à électrons d'un oscilloscope, un électron (charge $q=-e$ avec $e=\num{1,60e-19}{C}$, masse $m=\num{9,11e-31}{kg}$) est émis sans vitesse initiale par la cathode C, puis accéléré jusqu'à l'anode A par une tension $U_{AC}=\num{1500}{V}$ (l'anode est au potentiel le plus élevé).
\begin{enumerate}
\item Montrer que le poids de l'électron est négligeable devant la force électrique, sachant que la distance entre les plaques vaut $d=\SI{3,0}{cm}$.
\item Calculer la vitesse $v_A$ de l'électron à la sortie de l'anode.
\item Exprimer puis calculer son énergie cinétique en joules puis en électronvolts.
\end{enumerate}
On donne $g=\SI{9,8}{N.kg^{-1}}$.
\tcblower
\textbf{1. Comparaison des forces.} Le champ électrique entre les plaques vaut
\[ E=\dfrac{U}{d}=\dfrac{1500}{\num{3,0e-2}}=\SI{5,0e4}{V.m^{-1}}. \]
Force électrique : $F_e=eE=\num{1,60e-19}\times\num{5,0e4}=\SI{8,0e-15}{N}$.
Poids : $P=mg=\num{9,11e-31}\times 9,8=\SI{8,9e-30}{N}$.
Le rapport vaut $\dfrac{P}{F_e}\approx\num{1,1e-15}$ : le poids est totalement négligeable devant la force électrique. Seule $\vec{F}=q\vec{E}$ sera prise en compte.

\textbf{2. Vitesse en A.} Théorème de l'énergie cinétique entre C et A :
\[ E_{c,A}-E_{c,C}=W(\vec{F})=q\,U_{CA}=-e(V_C-V_A)=e\,U_{AC}. \]
Comme $v_C=0$ : $\dfrac{1}{2}mv_A^2=eU_{AC}$, d'où
\[ v_A=\sqrt{\dfrac{2eU_{AC}}{m}}=\sqrt{\dfrac{2\times\num{1,60e-19}\times 1500}{\num{9,11e-31}}}
=\sqrt{\num{5,27e14}}\approx\SI{2,3e7}{m.s^{-1}}. \]

\textbf{3. Énergie cinétique.}
\[ E_{c,A}=eU_{AC}=\num{1,60e-19}\times 1500=\SI{2,4e-16}{J}. \]
En électronvolts : $E_{c,A}=\SI{1500}{eV}=\SI{1,5}{keV}$, car un électron accéléré sous $\SI{1}{V}$ gagne $\SI{1}{eV}$.

\textbf{Conclusion :} l'électron quitte l'anode avec la vitesse $v_A\approx\SI{2,3e7}{m.s^{-1}}$ et une énergie cinétique de $\SI{1,5}{keV}$. L'énergie acquise ne dépend que de la tension accélératrice, pas de la géométrie du canon.
\end{exoresolu}

\courssub{Méthode 2 — Déviation électrique : mise en équation du mouvement}

\begin{methode}[Trajectoire parabolique entre les plaques]
Pour établir et exploiter le mouvement d'une particule pénétrant avec une vitesse $\vec{v}_0$ horizontale dans un champ $\vec{E}$ uniforme perpendiculaire à $\vec{v}_0$ :
\begin{enumerate}
\item \textbf{Choisir le repère} : origine O à l'entrée des plaques, $(Ox)$ selon $\vec{v}_0$, $(Oy)$ selon la direction de $\vec{E}$ (sens choisi selon le signe de $q$).
\item \textbf{Faire le bilan des forces} : poids négligeable (à justifier une fois), donc $\vec{F}=q\vec{E}$ seule.
\item \textbf{Appliquer la deuxième loi de Newton} : $\vec{a}=\dfrac{q\vec{E}}{m}$, d'où $a_x=0$ et $a_y=\dfrac{qE}{m}$ (constante).
\item \textbf{Intégrer deux fois} avec les conditions initiales $x(0)=y(0)=0$, $v_x(0)=v_0$, $v_y(0)=0$ :
\[ x(t)=v_0\,t \qquad ; \qquad y(t)=\dfrac{qE}{2m}\,t^2. \]
\item \textbf{Éliminer le temps} ($t=x/v_0$) pour obtenir la trajectoire :
\[ y(x)=\dfrac{qE}{2mv_0^2}\,x^2 \quad\text{(parabole).} \]
\item \textbf{Exploiter} : à la sortie des plaques ($x=\ell$), calculer la déviation $Y=y(\ell)$, l'angle de sortie $\tan\alpha=\dfrac{v_y}{v_x}=\dfrac{qE\ell}{mv_0^2}$, et vérifier que la particule ne heurte pas une plaque ($|Y|<d/2$).
\end{enumerate}
\end{methode}

\begin{exoresolu}[Déviation d'un électron entre deux plaques]
Un électron ($m=\num{9,11e-31}{kg}$, $q=-e$, $e=\num{1,60e-19}{C}$) pénètre en O avec la vitesse $v_0=\SI{2,0e7}{m.s^{-1}}$ entre deux plaques horizontales de longueur $\ell=\SI{4,0}{cm}$, distantes de $d=\SI{2,0}{cm}$, entre lesquelles règne un champ $\vec{E}$ vertical ascendant de norme $E=\SI{5,0e3}{V.m^{-1}}$.
\begin{enumerate}
\item Établir les équations horaires du mouvement et l'équation de la trajectoire.
\item Calculer l'ordonnée $Y_S$ du point de sortie S et vérifier que l'électron ne touche pas les plaques.
\item Calculer l'angle $\alpha$ de déviation à la sortie.
\end{enumerate}
\tcblower
\textbf{1. Équations du mouvement.} Le poids ($\approx\SI{9e-30}{N}$) est négligeable devant $eE=\SI{8,0e-16}{N}$. Deuxième loi de Newton : $\vec{a}=\dfrac{q\vec{E}}{m}$. L'électron étant chargé négativement, son accélération est dirigée vers le bas (opposée à $\vec{E}$) ; en choisissant $(Oy)$ vers le haut :
\[ a_x=0 \quad ; \quad a_y=-\dfrac{eE}{m}. \]
Par intégrations successives avec les conditions initiales :
\[ x(t)=v_0\,t \quad ; \quad y(t)=-\dfrac{eE}{2m}t^2. \]
Trajectoire : $t=\dfrac{x}{v_0}$, donc
\[ y(x)=-\dfrac{eE}{2mv_0^2}x^2. \]
C'est une branche de parabole : le mouvement est \cle{parabolique}.

\textbf{2. Déviation à la sortie.} Application numérique du coefficient :
\[ \dfrac{eE}{2mv_0^2}=\dfrac{\num{1,60e-19}\times\num{5,0e3}}{2\times\num{9,11e-31}\times(\num{2,0e7})^2}
=\dfrac{\num{8,0e-16}}{\num{7,29e-16}}\approx\SI{1,10}{m^{-1}}. \]
À la sortie, $x=\ell=\num{4,0e-2}\ \mathrm{m}$ :
\[ |Y_S|=1{,}10\times(\num{4,0e-2})^2=1{,}10\times\num{1,6e-3}=\SI{1,76e-3}{m}\approx\SI{1,8}{mm}. \]
Or $\dfrac{d}{2}=\SI{10}{mm}$ : $|Y_S|<d/2$, l'électron sort bien des plaques sans les toucher.

\textbf{3. Angle de déviation.} À la sortie :
\[ v_y=-\dfrac{eE}{m}\,t_S=-\dfrac{eE\ell}{mv_0}, \qquad \tan\alpha=\dfrac{|v_y|}{v_x}=\dfrac{eE\ell}{mv_0^2}. \]
\[ \tan\alpha=\dfrac{\num{1,60e-19}\times\num{5,0e3}\times\num{4,0e-2}}{\num{9,11e-31}\times\num{4,0e14}}
=\dfrac{\num{3,2e-17}}{\num{3,64e-16}}\approx\num{8,8e-2}. \]
Donc $\alpha\approx 5,0^{\circ}$.

\textbf{Conclusion :} la trajectoire est parabolique d'équation $y=-1{,}10\,x^2$ (en mètres) ; l'électron est dévié de $\SI{1,8}{mm}$ à la sortie et quitte les plaques sous un angle d'environ $5^{\circ}$.
\end{exoresolu}

\courssub{Méthode 3 — Déflexion sur l'écran (oscilloscope)}

\begin{methode}[Calcul de la déflexion sur l'écran]
Après la sortie des plaques, la particule n'est plus soumise à aucune force : son mouvement est \cle{rectiligne uniforme} (principe d'inertie) selon la tangente à la parabole au point de sortie S.
\begin{enumerate}
\item \textbf{Modéliser} : la tangente en S coupe l'axe $(Ox)$ en son milieu I ($x_I=\ell/2$) — propriété de la parabole $y=kx^2$ : la pente en $\ell$ vaut $2k\ell$, soit le double de la pente de la corde joignant O à S.
\item \textbf{Écrire} la déflexion sur l'écran placé à la distance $D$ du centre I des plaques :
\[ Y_{\text{écran}}=D\tan\alpha=\dfrac{qE\ell}{mv_0^2}\,D. \]
\item \textbf{Exprimer en fonction des tensions} : avec $E=\dfrac{U}{d}$ et $\dfrac{1}{2}mv_0^2=|q|U_A$ (tension accélératrice), on obtient
\[ Y_{\text{écran}}=\dfrac{U\,\ell\,D}{2\,d\,U_A} \quad\text{(indépendant de } q \text{ et } m\text{)}. \]
\item \textbf{Conclure} : la déflexion est proportionnelle à la tension $U$ appliquée aux plaques — c'est le principe de la mesure de tension à l'oscilloscope.
\end{enumerate}
\end{methode}

\begin{exoresolu}[Oscilloscope : déflexion d'un spot]
Dans un oscilloscope, des électrons accélérés sous $U_A=\SI{1500}{V}$ traversent les plaques de déviation verticale ($\ell=\SI{4,0}{cm}$, $d=\SI{2,0}{cm}$) soumises à la tension $U=\SI{100}{V}$. L'écran est situé à $D=\SI{20}{cm}$ du milieu des plaques.
\begin{enumerate}
\item Calculer la vitesse $v_0$ des électrons à l'entrée des plaques ($e=\num{1,60e-19}{C}$, $m=\num{9,11e-31}{kg}$).
\item Calculer la déflexion $Y$ du spot sur l'écran.
\item Quelle tension faut-il appliquer pour obtenir une déflexion de $\SI{3,0}{cm}$ ?
\end{enumerate}
\tcblower
\textbf{1. Vitesse d'entrée.} Théorème de l'énergie cinétique dans le canon : $\dfrac{1}{2}mv_0^2=eU_A$, donc
\[ v_0=\sqrt{\dfrac{2eU_A}{m}}=\sqrt{\dfrac{2\times\num{1,60e-19}\times 1500}{\num{9,11e-31}}}\approx\SI{2,3e7}{m.s^{-1}}. \]

\textbf{2. Déflexion sur l'écran.} Le champ entre les plaques vaut $E=\dfrac{U}{d}=\dfrac{100}{\num{2,0e-2}}=\SI{5,0e3}{V.m^{-1}}$.
Angle de sortie :
\[ \tan\alpha=\dfrac{eE\ell}{mv_0^2}=\dfrac{\num{1,60e-19}\times\num{5,0e3}\times\num{4,0e-2}}{\num{9,11e-31}\times(\num{2,3e7})^2}
=\dfrac{\num{3,2e-17}}{\num{4,82e-16}}\approx\num{6,6e-2}. \]
Après les plaques, le mouvement est rectiligne uniforme (plus aucune force) ; la droite suivie passe par le milieu I des plaques :
\[ Y=D\tan\alpha=0{,}20\times\num{6,6e-2}=\SI{1,3e-2}{m}=\SI{1,3}{cm}. \]
Vérification par la formule directe : $Y=\dfrac{U\ell D}{2dU_A}=\dfrac{100\times 4{,}0\times 20}{2\times 2{,}0\times 1500}\ \mathrm{cm}=\dfrac{8000}{6000}\ \mathrm{cm}\approx\SI{1,3}{cm}$. Cohérent.

\textbf{3. Tension pour $Y'=\SI{3,0}{cm}$.} La déflexion est proportionnelle à $U$ :
\[ U'=U\times\dfrac{Y'}{Y}=100\times\dfrac{3{,}0}{1{,}3}\approx\SI{2,3e2}{V}. \]
(Le calcul exact avec $Y=\SI{1,33}{cm}$ donne $U'\approx\SI{225}{V}$.)

\textbf{Conclusion :} le spot est dévié de $\SI{1,3}{cm}$ sous $\SI{100}{V}$ ; il faut environ $\SI{225}{V}$ pour une déflexion de $\SI{3,0}{cm}$. La proportionnalité $Y\propto U$ fait de l'oscilloscope un voltmètre.
\end{exoresolu}

\courssub{Méthode 4 — Mouvement dans un champ magnétique uniforme}

\begin{methode}[Mouvement circulaire uniforme dans $\vec{B}$]
Pour une particule de vitesse $\vec{v}_0$ \textbf{perpendiculaire} à un champ magnétique uniforme $\vec{B}$ :
\begin{enumerate}
\item \textbf{Écrire la force de Lorentz} : $\vec{F}=q\,\vec{v}\wedge\vec{B}$, toujours perpendiculaire à $\vec{v}$.
\item \textbf{Justifier l'uniformité du mouvement} : la puissance $P=\vec{F}\cdot\vec{v}=0$, donc l'énergie cinétique et la norme $v$ sont constantes (la force magnétique ne travaille pas).
\item \textbf{Montrer que le mouvement est circulaire} : la deuxième loi de Newton projetée sur la normale donne une accélération normale constante $\dfrac{|q|vB}{m}=\dfrac{v^2}{R}$, d'où un rayon constant :
\[ R=\dfrac{mv}{|q|B}. \]
\item \textbf{Calculer la période} : $T=\dfrac{2\pi R}{v}=\dfrac{2\pi m}{|q|B}$ — remarquer que $T$ ne dépend pas de $v$.
\item \textbf{Orientations} : le sens de rotation dépend du signe de $q$ (règle du tire-bouchon ou des trois doigts avec $\vec{v}\wedge\vec{B}$, en inversant si $q<0$).
\item \textbf{Vérifier les conditions d'application} : $\vec{v}_0\perp\vec{B}$ ; sinon le mouvement est hélicoïdal (hors programme, mais le signaler).
\end{enumerate}
\end{methode}

\begin{exoresolu}[Proton dans le champ d'un électroaimant]
Un proton ($q=e=\num{1,60e-19}{C}$, $m=\num{1,67e-27}{kg}$) pénètre avec la vitesse $v_0=\SI{1,0e6}{m.s^{-1}}$, perpendiculairement à un champ magnétique uniforme $B=\SI{0,20}{T}$.
\begin{enumerate}
\item Montrer que le mouvement du proton est circulaire uniforme et calculer le rayon $R$ de sa trajectoire.
\item Calculer la période $T$ du mouvement. Que devient $T$ si la vitesse double ?
\item Le poids du proton est-il négligeable ?
\end{enumerate}
\tcblower
\textbf{1. Nature du mouvement et rayon.} La force magnétique $\vec{F}=q\,\vec{v}\wedge\vec{B}$ est perpendiculaire à $\vec{v}$ : sa puissance est nulle, donc $v=v_0$ est constante (mouvement uniforme). La deuxième loi de Newton, projetée sur la normale à la trajectoire (le mouvement a lieu dans le plan perpendiculaire à $\vec{B}$), s'écrit :
\[ ev_0B=m\dfrac{v_0^2}{R} \quad\Longrightarrow\quad R=\dfrac{mv_0}{eB}. \]
$R$ est constant : la trajectoire est un cercle. Application numérique :
\[ R=\dfrac{\num{1,67e-27}\times\num{1,0e6}}{\num{1,60e-19}\times 0{,}20}
=\dfrac{\num{1,67e-21}}{\num{3,2e-20}}\approx\SI{5,2e-2}{m}=\SI{5,2}{cm}. \]

\textbf{2. Période.}
\[ T=\dfrac{2\pi R}{v_0}=\dfrac{2\pi m}{eB}=\dfrac{2\pi\times\num{1,67e-27}}{\num{1,60e-19}\times 0{,}20}
=\dfrac{\num{1,05e-26}}{\num{3,2e-20}}\approx\SI{3,3e-7}{s}. \]
$T$ ne dépend pas de la vitesse : si $v_0$ double, $R$ double mais $T$ reste égale à $\SI{3,3e-7}{s}$. C'est la propriété fondamentale exploitée par le cyclotron.

\textbf{3. Poids.} $P=mg=\num{1,67e-27}\times 9,8\approx\SI{1,6e-26}{N}$, contre $F=ev_0B=\num{1,60e-19}\times\num{1,0e6}\times 0{,}20=\SI{3,2e-14}{N}$. Le poids est environ $10^{12}$ fois plus faible : il est négligeable.

\textbf{Conclusion :} le proton décrit un cercle de rayon $\SI{5,2}{cm}$ à vitesse constante, avec une période $T\approx\SI{0,33}{\micro s}$ indépendante de sa vitesse.
\end{exoresolu}

\courssub{Méthode 5 — Spectrographe de masse et cyclotron}

\begin{methode}[Exploiter le spectrographe de masse et le cyclotron]
\textbf{Spectrographe de masse} (séparation d'isotopes) :
\begin{enumerate}
\item \textbf{Chambre d'ionisation} : production d'ions de charge $q$ connue.
\item \textbf{Accélération} sous la tension $U$ : $\dfrac{1}{2}mv^2=|q|U$ donne $v=\sqrt{\dfrac{2|q|U}{m}}$.
\item \textbf{Chambre de déviation} ($\vec{B}\perp\vec{v}$) : demi-cercle de rayon $R=\dfrac{mv}{|q|B}=\dfrac{1}{B}\sqrt{\dfrac{2mU}{|q|}}$.
\item \textbf{Exploitation} : les ions frappent la plaque photographique à la distance $2R$ de la fente d'entrée ; à $q$, $U$, $B$ fixés, $R\propto\sqrt{m}$ : les isotopes sont séparés. Rapport des rayons : $\dfrac{R_2}{R_1}=\sqrt{\dfrac{m_2}{m_1}}$.
\end{enumerate}
\textbf{Cyclotron} (accélérateur de particules) :
\begin{enumerate}
\item La particule décrit des demi-cercles dans les « dés » où règne $\vec{B}$ ; entre les dés, une tension alternative l'accélère à chaque passage.
\item La fréquence de la tension doit égaler la fréquence cyclotron $f=\dfrac{|q|B}{2\pi m}$, constante car $T$ est indépendante de $v$.
\item L'énergie cinétique augmente de $|q|U$ à chaque traversée de l'intervalle ; après $n$ passages : $E_c=n\,|q|U$ (si $E_{c,0}\approx 0$).
\end{enumerate}
\end{methode}

\begin{exoresolu}[Séparation des isotopes de l'uranium]
Dans un spectrographe de masse, des ions \ce{^{235}_{92}U+} et \ce{^{238}_{92}U+} (charge $q=e=\num{1,60e-19}{C}$) sont accélérés sous $U=\SI{4,0e3}{V}$ puis pénètrent dans un champ $B=\SI{0,10}{T}$ perpendiculaire à leur vitesse. On donne $m_1=235\ u$, $m_2=238\ u$, $1\ u=\num{1,66e-27}{kg}$.
\begin{enumerate}
\item Calculer les vitesses $v_1$ et $v_2$ à l'entrée de la chambre de déviation.
\item Calculer les rayons $R_1$ et $R_2$ des trajectoires.
\item En déduire la distance $d$ séparant les deux points d'impact sur la plaque.
\end{enumerate}
\tcblower
\textbf{1. Vitesses.} Théorème de l'énergie cinétique : $\dfrac{1}{2}mv^2=eU$, donc $v=\sqrt{\dfrac{2eU}{m}}$.
\[ v_1=\sqrt{\dfrac{2\times\num{1,60e-19}\times\num{4,0e3}}{235\times\num{1,66e-27}}}
=\sqrt{\dfrac{\num{1,28e-15}}{\num{3,90e-25}}}=\sqrt{\num{3,28e9}}\approx\SI{5,73e4}{m.s^{-1}}. \]
\[ v_2=\sqrt{\dfrac{\num{1,28e-15}}{238\times\num{1,66e-27}}}
=\sqrt{\num{3,24e9}}\approx\SI{5,69e4}{m.s^{-1}}. \]
L'isotope le plus lourd est, comme prévu, légèrement plus lent ($v\propto 1/\sqrt{m}$).

\textbf{2. Rayons.} $R=\dfrac{mv}{eB}$ :
\[ R_1=\dfrac{235\times\num{1,66e-27}\times\num{5,73e4}}{\num{1,60e-19}\times 0{,}10}
=\dfrac{\num{2,24e-17}}{\num{1,60e-20}}\approx\SI{1,40}{m}. \]
\[ R_2=\dfrac{238\times\num{1,66e-27}\times\num{5,69e4}}{\num{1,60e-19}\times 0{,}10}
=\dfrac{\num{2,25e-17}}{\num{1,60e-20}}\approx\SI{1,41}{m}. \]
Vérification par le rapport : $\dfrac{R_2}{R_1}=\sqrt{\dfrac{238}{235}}=\sqrt{1{,}0128}\approx 1{,}0064$, donc $R_2\approx 1{,}40\times 1{,}0064\approx\SI{1,41}{m}$. Cohérent.

\textbf{3. Distance entre les impacts.} Chaque ion décrit un demi-cercle : l'impact est à la distance $2R$ de la fente. Donc
\[ d=2R_2-2R_1=2(R_2-R_1)\approx 2\times(\num{1,406}-\num{1,396})\approx\SI{1,8e-2}{m}\approx\SI{1,8}{cm}. \]

\textbf{Conclusion :} les deux isotopes, de masses très proches, sont séparés d'environ $\SI{1,8}{cm}$ sur la plaque : le spectrographe résout efficacement des masses voisines car $R\propto\sqrt{m}$.
\end{exoresolu}

\courssub{Méthode 6 — Filtre de Wien : champs électrique et magnétique croisés}

\begin{methode}[Filtre de vitesse (sélecteur de Wien)]
Des champs $\vec{E}$ et $\vec{B}$ uniformes, perpendiculaires entre eux et à la vitesse $\vec{v}$ de la particule, sont orientés de sorte que les forces électrique et magnétique soient \textbf{opposées}.
\begin{enumerate}
\item \textbf{Écrire les deux forces} : $\vec{F}_e=q\vec{E}$ et $\vec{F}_m=q\,\vec{v}\wedge\vec{B}$, de sens opposés par construction.
\item \textbf{Condition de trajectoire rectiligne} (non déviée) : $\vec{F}_e+\vec{F}_m=\vec{0}$, soit en normes
\[ |q|E=|q|vB \quad\Longrightarrow\quad v=\dfrac{E}{B}. \]
\item \textbf{Remarquer} que la condition est indépendante de $q$ (signe et valeur) et de $m$ : le filtre sélectionne une \emph{vitesse}, pas une particule.
\item \textbf{Prévoir le comportement des autres particules} : si $v>E/B$, la force magnétique l'emporte ; si $v<E/B$, c'est la force électrique. Seules les particules de vitesse $E/B$ traversent le filtre sans déviation.
\end{enumerate}
\end{methode}

\begin{exoresolu}[Sélecteur de vitesse]
Entre deux plaques distantes de $d=\SI{5,0}{cm}$ soumises à la tension $U=\SI{500}{V}$, règne un champ magnétique uniforme $B=\SI{2,0e-2}{T}$, perpendiculaire au champ électrique et à la vitesse des particules. Un faisceau d'ions de vitesses diverses traverse le dispositif.
\begin{enumerate}
\item Calculer la norme $E$ du champ électrique.
\item Déterminer la vitesse $v_0$ des ions qui traversent le filtre sans déviation.
\item Un ion de vitesse $v>v_0$ est-il dévié du côté de la force électrique ou de la force magnétique ? Justifier.
\end{enumerate}
\tcblower
\textbf{1. Champ électrique.}
\[ E=\dfrac{U}{d}=\dfrac{500}{\num{5,0e-2}}=\SI{1,0e4}{V.m^{-1}}. \]

\textbf{2. Vitesse sélectionnée.} Les champs sont orientés de sorte que $\vec{F}_e=q\vec{E}$ et $\vec{F}_m=q\vec{v}\wedge\vec{B}$ soient opposées. Trajectoire rectiligne non déviée si et seulement si les normes sont égales :
\[ |q|E=|q|v_0B \quad\Longrightarrow\quad v_0=\dfrac{E}{B}=\dfrac{\num{1,0e4}}{\num{2,0e-2}}=\SI{5,0e5}{m.s^{-1}}. \]
Cette condition ne dépend ni de la charge ni de la masse des ions : tous les ions de vitesse $\SI{5,0e5}{m.s^{-1}}$ traversent le filtre en ligne droite.

\textbf{3. Ion plus rapide.} La force électrique $F_e=|q|E$ ne dépend pas de la vitesse, tandis que $F_m=|q|vB$ croît avec $v$. Si $v>v_0$, alors $F_m>F_e$ : la résultante est du côté de la force magnétique, et l'ion est dévié dans le sens de $\vec{F}_m$. (Symétriquement, un ion de vitesse $v<v_0$ est dévié du côté de $\vec{F}_e$.)

\textbf{Conclusion :} le filtre laisse passer sans déviation les ions de vitesse $v_0=\dfrac{E}{B}=\SI{5,0e5}{m.s^{-1}}$ ; les autres sont écartés, d'où le nom de \cle{sélecteur de vitesse}.
\end{exoresolu}

\begin{attention}[Les 5 erreurs qui coûtent des points au Bac]
\begin{enumerate}
\item \textbf{Oublier de justifier que le poids est négligeable.} La deuxième loi de Newton $\vec{a}=q\vec{E}/m$ n'est valable que si l'on a comparé $mg$ à $|q|E$ (ou énoncé l'hypothèse). Une ligne de calcul suffit, mais elle est exigée.
\item \textbf{Se tromper de signe dans le travail électrique.} Le travail vaut $W=q\,U_{CA}=q(V_C-V_A)$ : pour un électron ($q=-e$) accéléré vers l'anode, c'est $W=eU_{AC}>0$. Une erreur de signe conduit à une vitesse imaginaire — vérifier toujours que $E_c$ finale $>0$.
\item \textbf{Confondre déviation dans les plaques et déflexion sur l'écran.} Dans les plaques : trajectoire parabolique $y=\dfrac{qE}{2mv_0^2}x^2$. Après les plaques : mouvement rectiligne uniforme (plus de champ !), la droite passant par le milieu des plaques. Ne pas prolonger la parabole jusqu'à l'écran.
\item \textbf{Croire que la force magnétique modifie la vitesse.} $\vec{F}=q\vec{v}\wedge\vec{B}$ est perpendiculaire à $\vec{v}$ : elle ne travaille pas, la norme de la vitesse est constante. Écrire « la particule accélère » dans un champ $\vec{B}$ seul est une faute grave ; dire « mouvement circulaire \emph{uniforme} ».
\item \textbf{Oublier les conditions d'application des formules.} $R=\dfrac{mv}{|q|B}$ exige $\vec{v}_0\perp\vec{B}$ ; $v=\dfrac{E}{B}$ exige $\vec{E}\perp\vec{B}\perp\vec{v}$ avec forces opposées ; $T=\dfrac{2\pi m}{|q|B}$ suppose le champ uniforme. Citer la condition avant d'appliquer la formule, et toujours convertir les données en unités SI (cm $\to$ m, eV $\to$ J) avant tout calcul numérique.
\end{enumerate}
\end{attention}

\newpage

\coursec{Exercices}

\courssub{Je vérifie mes connaissances}

\begin{exercice}[QCM : une seule réponse exacte]
Pour chaque question, choisir la bonne réponse en justifiant brièvement.
\begin{enumerate}
\item Un électron de charge $-e$ et de masse $m$, initialement au repos, est accéléré par une tension $U$. Sa vitesse finale vaut :
\begin{enumerate}
\item $v = \sqrt{\dfrac{2eU}{m}}$ \quad (b) $v = \dfrac{2eU}{m}$ \quad (c) $v = \sqrt{\dfrac{eU}{2m}}$ \quad (d) $v = \dfrac{eU}{m}$
\end{enumerate}
\item Une particule chargée pénètre dans un champ magnétique uniforme $\vec{B}$ avec une vitesse $\vec{v}_0$ parallèle à $\vec{B}$. Sa trajectoire est :
\begin{enumerate}
\item un cercle \quad (b) une parabole \quad (c) une droite \quad (d) une hélice
\end{enumerate}
\item Le rayon de la trajectoire circulaire d'une particule de charge $q$, de masse $m$, de vitesse $v$ perpendiculaire à un champ magnétique $B$ est :
\begin{enumerate}
\item $R = \dfrac{|q|B}{mv}$ \quad (b) $R = \dfrac{mv}{|q|B}$ \quad (c) $R = \dfrac{mB}{|q|v}$ \quad (d) $R = \dfrac{|q|v}{mB}$
\end{enumerate}
\item Dans un filtre de Wien (champs $\vec{E}$ et $\vec{B}$ croisés), les particules non déviées ont une vitesse :
\begin{enumerate}
\item $v = \dfrac{B}{E}$ \quad (b) $v = EB$ \quad (c) $v = \dfrac{E}{B}$ \quad (d) $v = \sqrt{\dfrac{E}{B}}$
\end{enumerate}
\end{enumerate}
\end{exercice}

\begin{exercice}[Vrai ou faux ? Justifier]
Répondre par vrai ou faux en justifiant chaque réponse par un argument physique.
\begin{enumerate}
\item La force magnétique de Lorentz $\vec{F} = q\,\vec{v} \wedge \vec{B}$ peut modifier la valeur de la vitesse d'une particule chargée.
\item Le travail de la force électrique subie par une particule de charge $q$ traversant une tension $U$ est $W = |q|U$, quel que soit le signe de $q$.
\item Entre les plaques d'un condensateur plan, un électron animé d'une vitesse initiale horizontale décrit une trajectoire parabolique.
\item La période du mouvement circulaire d'une particule chargée dans un champ magnétique uniforme dépend de sa vitesse.
\end{enumerate}
\end{exercice}

\begin{exercice}[Définitions et formules]
\begin{enumerate}
\item Définir : champ électrique uniforme ; force de Lorentz magnétique.
\item Énoncer le théorème de l'énergie cinétique.
\item Donner l'expression de la force électrique $\vec{F}_e$ subie par une charge $q$ placée dans un champ $\vec{E}$, puis celle de la force magnétique $\vec{F}_m$ dans un champ $\vec{B}$. Préciser les unités SI de chaque grandeur.
\item Vérifier par analyse dimensionnelle que $R = \dfrac{mv}{|q|B}$ est homogène à une longueur. On rappelle qu'une force de $1~\mathrm{N}$ s'exprime par $1~\mathrm{kg\,m\,s^{-2}}$ et que $F_m = |q|vB$.
\end{enumerate}
\end{exercice}

\begin{exercice}[Calculs directs]
Données : $e = \num{1,60e-19}$ C ; masse de l'électron $m_e = \num{9,11e-31}$ kg ; masse du proton $m_p = \num{1,67e-27}$ kg.
\begin{enumerate}
\item Calculer la vitesse acquise par un électron initialement au repos, accéléré sous une tension $U = 500$ V.
\item Calculer le rayon de la trajectoire d'un proton de vitesse $v = \num{1,0e6}$ m/s pénétrant perpendiculairement dans un champ magnétique uniforme $B = \num{0,10}$ T.
\item Calculer la période de ce mouvement circulaire. Conclusion sur la dépendance de $T$ vis-à-vis de $v$ ?
\end{enumerate}
\end{exercice}

\courssub{Je m'entraîne}

\begin{exercice}[Canon à électrons d'un oscilloscope]
Dans l'oscilloscope du laboratoire d'un lycée de Yaoundé, les électrons sont émis par la cathode avec une vitesse négligeable puis accélérés vers l'anode par une tension $U = 1500$ V.
Données : $e = \num{1,60e-19}$ C ; $m_e = \num{9,11e-31}$ kg.
\begin{enumerate}
\item Représenter sur un schéma les plaques, le champ $\vec{E}$ et la force $\vec{F}$ subie par l'électron.
\item En appliquant le théorème de l'énergie cinétique, établir l'expression de la vitesse $v_0$ de l'électron à la sortie de l'anode.
\item Calculer $v_0$.
\item Quelle tension faudrait-il appliquer pour doubler cette vitesse ?
\end{enumerate}
\end{exercice}

\begin{exercice}[Proton dans un champ magnétique]
Un proton pénètre en un point $O$ dans une zone où règne un champ magnétique uniforme $\vec{B}$ de norme $B = \num{0,20}$ T, avec une vitesse $\vec{v}_0$ de norme $v_0 = \num{2,0e6}$ m/s perpendiculaire à $\vec{B}$.
Données : $e = \num{1,60e-19}$ C ; $m_p = \num{1,67e-27}$ kg.
\begin{enumerate}
\item Montrer que le poids du proton est négligeable devant la force magnétique.
\item En appliquant la deuxième loi de Newton dans le repère de Frenet, démontrer que le mouvement est circulaire et uniforme.
\item Calculer le rayon $R$ de la trajectoire.
\item Calculer la période $T$ du mouvement.
\item Représenter sur un schéma $\vec{v}_0$, $\vec{B}$ (sortant du plan de la feuille) et la force magnétique en $O$ ; en déduire le sens de rotation du proton.
\end{enumerate}
\end{exercice}

\begin{exercice}[Filtre de Wien]
Un filtre de vitesse est constitué d'un condensateur plan créant un champ électrique uniforme $E = \num{2,0e4}$ V/m et d'un champ magnétique uniforme $\vec{B}$ perpendiculaire à $\vec{E}$ et à la vitesse $\vec{v}$ des ions. On souhaite sélectionner des ions de vitesse $v_0 = \num{4,0e5}$ m/s.
\begin{enumerate}
\item Faire un schéma montrant $\vec{E}$, $\vec{B}$, $\vec{v}_0$ et les forces $\vec{F}_e$ et $\vec{F}_m$ pour un ion positif.
\item Établir la condition pour que l'ion traverse le filtre sans déviation.
\item Calculer la valeur $B$ du champ magnétique nécessaire.
\item Que devient un ion de vitesse $v > v_0$ ? De vitesse $v < v_0$ ? Justifier.
\end{enumerate}
\end{exercice}

\begin{exercice}[Déviation d'un électron entre deux plaques]
Un électron pénètre avec une vitesse horizontale $v_0 = \num{2,0e7}$ m/s entre les plaques horizontales d'un condensateur plan de longueur $L = 5{,}0$ cm, distantes de $d = 4{,}0$ cm, soumises à une tension $U' = 200$ V.
Données : $e = \num{1,60e-19}$ C ; $m_e = \num{9,11e-31}$ kg.
\begin{enumerate}
\item Calculer l'intensité du champ électrique $E'$ entre les plaques.
\item Établir les équations horaires $x(t)$ et $y(t)$ du mouvement de l'électron entre les plaques.
\item En déduire l'équation de la trajectoire $y(x)$ et préciser sa nature.
\item Calculer la durée de passage entre les plaques et la déviation verticale $y_S$ à la sortie.
\end{enumerate}
\end{exercice}

\courssub{Je me prépare au Bac}

\begin{exercice}[Type Bac : oscilloscope cathodique]
Un oscilloscope, utilisé au laboratoire de physique d'un lycée de Douala, comporte un canon à électrons, des plaques de déviation verticale et un écran fluorescent. Un électron, émis par la cathode $C$ avec une vitesse négligeable, est accéléré par une tension $U = 1000$ V appliquée entre $C$ et l'anode $A$. Il pénètre ensuite en $O$, avec la vitesse $\vec{v}_0$ horizontale, entre les plaques horizontales $P_1$ et $P_2$ d'un condensateur plan de longueur $L = 6{,}0$ cm, distantes de $d = 3{,}0$ cm, soumises à la tension $U' = 150$ V ($P_1$ étant la plaque positive, située en haut). L'écran est placé à la distance $D = 25$ cm du centre $I$ du condensateur.
Données : $e = \num{1,60e-19}$ C ; $m_e = \num{9,11e-31}$ kg. On néglige le poids de l'électron devant les forces électriques.
\begin{enumerate}
\item \textbf{Accélération.}
\begin{enumerate}
\item Représenter le champ $\vec{E}$ entre $C$ et $A$ et la force accélératrice.
\item Établir l'expression de $v_0$ en fonction de $e$, $U$ et $m_e$ ; calculer $v_0$.
\end{enumerate}
\item \textbf{Déviation entre les plaques.}
\begin{enumerate}
\item Indiquer le sens de déviation de l'électron entre les plaques. Justifier.
\item Établir les équations horaires du mouvement et l'équation cartésienne $y(x)$ de la trajectoire entre les plaques.
\item Calculer la déviation $y_S$ à la sortie des plaques et l'angle $\alpha$ que fait la vitesse de sortie avec l'horizontale ($\tan\alpha = \dfrac{v_{yS}}{v_0}$).
\end{enumerate}
\item \textbf{Déflexion sur l'écran.}
\begin{enumerate}
\item Justifier qu'à la sortie des plaques le mouvement est rectiligne uniforme.
\item Montrer que la déflexion sur l'écran s'écrit $Y = \dfrac{eU'L}{m_e v_0^2 d}\left(\dfrac{L}{2} + D\right)$, puis calculer $Y$.
\item On double la tension accélératrice $U$. Que devient $Y$ ? Commenter l'intérêt pratique.
\end{enumerate}
\end{enumerate}
\end{exercice}

\begin{exercice}[Type Bac : spectrographe de masse et isotopes du brome]
Dans un laboratoire d'analyses, on utilise un spectrographe de masse pour séparer les isotopes \ce{^{79}Br} et \ce{^{81}Br} du brome. Les ions \ce{Br+} de charge $q = +e$, produits dans une chambre d'ionisation avec une vitesse négligeable, sont accélérés par une tension $U = 4000$ V. Ils pénètrent ensuite dans une chambre de déviation où règne un champ magnétique uniforme $B = \num{0,50}$ T perpendiculaire à leur vitesse, décrivent un demi-cercle et viennent frapper une plaque photographique.
Données : $e = \num{1,60e-19}$ C ; unité de masse atomique $u = \num{1,66e-27}$ kg ; $m(\ce{^{79}Br}) = 79\,u$ ; $m(\ce{^{81}Br}) = 81\,u$.
\begin{enumerate}
\item Établir l'expression de la vitesse $v$ d'un ion de masse $m$ à l'entrée de la chambre de déviation, en fonction de $e$, $U$ et $m$.
\item Démontrer que dans la chambre de déviation le mouvement de chaque ion est circulaire et uniforme ; établir l'expression du rayon $R = \dfrac{1}{B}\sqrt{\dfrac{2mU}{e}}$.
\item Calculer les rayons $R_1$ et $R_2$ des trajectoires des ions \ce{^{79}Br+} et \ce{^{81}Br+}.
\item Les impacts sont situés aux diamètres des demi-cercles. Calculer l'écart $\Delta = 2(R_2 - R_1)$ entre les deux impacts sur la plaque.
\item Comment évolue l'écart $\Delta$ si l'on augmente la tension $U$ ? Si l'on augmente $B$ ? Conclure sur les réglages favorisant une bonne séparation.
\end{enumerate}
\end{exercice}

\begin{exercice}[Type Bac : cyclotron et filtre de Wien en médecine]
Un centre hospitalier universitaire envisage la production locale d'isotopes médicaux à l'aide d'un cyclotron accélérant des protons. Le cyclotron est constitué de deux demi-cylindres creux (les « dees ») placés dans un champ magnétique uniforme $B = 1{,}0$ T perpendiculaire au plan des dees. Une tension alternative de fréquence $f$ est appliquée entre les dees. Le rayon maximal des dees est $R_{\max} = 50$ cm.
Données : $e = \num{1,60e-19}$ C ; $m_p = \num{1,67e-27}$ kg ; $1~\mathrm{eV} = \num{1,60e-19}$ J.
\begin{enumerate}
\item \textbf{Étude du cyclotron.}
\begin{enumerate}
\item Expliquer qualitativement le rôle du champ magnétique et celui de la tension alternative.
\item Montrer que la période du mouvement circulaire d'un proton dans un dee est $T = \dfrac{2\pi m_p}{eB}$. Que constate-t-on ?
\item En déduire la fréquence $f$ que doit avoir la tension alternative pour que les protons soient accélérés à chaque passage entre les dees. Calculer $f$.
\item Établir l'expression de l'énergie cinétique maximale des protons $E_{c,\max} = \dfrac{e^2 B^2 R_{\max}^2}{2m_p}$ ; la calculer en joules puis en MeV.
\end{enumerate}
\item \textbf{Sélection par filtre de Wien.} À la sortie du cyclotron, on veut ne conserver que les protons de vitesse $v_0 = \num{4,0e6}$ m/s. On utilise un filtre de Wien : champ électrique $E = \num{2,0e5}$ V/m et champ magnétique $B'$ croisés.
\begin{enumerate}
\item Calculer la valeur $B'$ du champ magnétique du filtre.
\item Un proton arrive avec une vitesse $v > v_0$. Quelle force l'emporte ? Vers quelle plaque est-il dévié si $\vec{F}_e$ est dirigée vers le bas ? Justifier.
\end{enumerate}
\end{enumerate}
\end{exercice}

\newpage

\coursec{Corrigés des exercices}

\begin{corrige}[Exercice 1 — QCM : une seule réponse exacte]
\begin{enumerate}
\item \textbf{Réponse (a)} : $v = \sqrt{\dfrac{2eU}{m}}$.\\
\emph{Justification :} le théorème de l'énergie cinétique appliqué entre la cathode ($v_i \approx 0$) et l'anode ($v_f = v$) donne $\frac{1}{2}mv^2 - 0 = W(\vec{F}_e) = eU$, d'où $v = \sqrt{2eU/m}$. Vérification dimensionnelle : $[eU] = \mathrm{J} = \mathrm{kg\,m^2\,s^{-2}}$, donc $\sqrt{[eU]/[m]} = \mathrm{m\,s^{-1}}$ : homogène à une vitesse. Les réponses (b), (c), (d) ne sont pas homogènes à une vitesse ou ne satisfont pas le bilan énergétique.
\item \textbf{Réponse (c)} : une droite.\\
\emph{Justification :} si $\vec{v}_0 \parallel \vec{B}$, alors $\vec{F}_m = q\,\vec{v}_0\wedge\vec{B} = \vec{0}$ (produit vectoriel de vecteurs colinéaires, $\sin 0 = 0$). Le poids étant négligeable, la somme des forces est nulle : d'après le principe d'inertie, le mouvement est rectiligne uniforme.
\item \textbf{Réponse (b)} : $R = \dfrac{mv}{|q|B}$.\\
\emph{Justification :} la force magnétique joue le rôle de force centripète : $|q|vB = m\dfrac{v^2}{R}$, d'où $R = \dfrac{mv}{|q|B}$. Vérification dimensionnelle : $[B] = \mathrm{kg\,C^{-1}\,s^{-1}}$ (car $F=qvB$), donc $\dfrac{[m][v]}{[q][B]} = \dfrac{\mathrm{kg\,m\,s^{-1}}}{\mathrm{C\,kg\,C^{-1}\,s^{-1}}} = \mathrm{m}$. Homogène à une longueur.
\item \textbf{Réponse (c)} : $v = \dfrac{E}{B}$.\\
\emph{Justification :} non-déviation $\Leftrightarrow \vec{F}_e + \vec{F}_m = \vec{0} \Leftrightarrow |q|E = |q|vB \Leftrightarrow v = E/B$. Analyse dimensionnelle : $[E/B] = \dfrac{\mathrm{V/m}}{\mathrm{T}} = \dfrac{\mathrm{N/C}}{\mathrm{N/(A\,m)}} = \mathrm{m/s}$.
\end{enumerate}
\end{corrige}

\begin{corrige}[Exercice 2 — Vrai ou faux ? Justifier]
\begin{enumerate}
\item \textbf{Faux.} La force magnétique de Lorentz est à chaque instant perpendiculaire à la vitesse : $\vec{F}_m \perp \vec{v}$. Sa puissance est donc nulle : $\mathcal{P} = \vec{F}_m\cdot\vec{v} = 0$, et son travail aussi. D'après le théorème de l'énergie cinétique, $E_c = \frac{1}{2}mv^2$ est constante : la force magnétique \textbf{courbe la trajectoire sans modifier la norme de la vitesse}.
\item \textbf{Faux.} Le travail de la force électrique est $W_{A\to B} = q(V_A - V_B) = qU_{AB}$ : il dépend du \textbf{signe de $q$} et du \textbf{sens de la tension}. Pour un électron ($q=-e$) accéléré de la cathode $C$ vers l'anode $A$ avec $U = V_A - V_C > 0$ : $W_{C\to A} = q(V_C - V_A) = (-e)(-U) = +eU > 0$ (travail moteur). La formule $W = |q|U$ ne donne que la valeur absolue du transfert d'énergie, pas le travail algébrique.
\item \textbf{Vrai.} Entre les plaques, l'électron ne subit que la force électrique constante $\vec{F}_e = -e\vec{E}$ (poids négligeable). Le mouvement est uniforme selon $Ox$ ($x = v_0 t$) et uniformément accéléré selon $Oy$ ($y = \frac{1}{2}\frac{eE}{m_e}t^2$). En éliminant le temps : $y = \dfrac{eE}{2m_e v_0^2}x^2$ : trajectoire \textbf{parabolique}.
\item \textbf{Faux.} La période vaut $T = \dfrac{2\pi R}{v} = \dfrac{2\pi m}{|q|B}$ : elle ne dépend que de $m$, $|q|$ et $B$, et \textbf{pas de la vitesse} (en régime non relativiste). C'est d'ailleurs la propriété d'isochronisme qui rend le cyclotron possible.
\end{enumerate}
\end{corrige}

\begin{corrige}[Exercice 3 — Définitions et formules]
\begin{enumerate}
\item \textbf{Champ électrique uniforme :} champ électrique dont le vecteur $\vec{E}$ garde la même norme, la même direction et le même sens en tout point de la région considérée (réalisé entre les armatures d'un condensateur plan, loin des bords).\\
\textbf{Force magnétique de Lorentz :} force subie par une particule de charge $q$ animée de la vitesse $\vec{v}$ dans un champ magnétique $\vec{B}$ : $\vec{F}_m = q\,\vec{v}\wedge\vec{B}$. Elle est perpendiculaire à $\vec{v}$ et à $\vec{B}$, de norme $F_m = |q|vB\sin\theta$, où $\theta$ est l'angle entre $\vec{v}$ et $\vec{B}$.
\item \textbf{Théorème de l'énergie cinétique :} dans un référentiel galiléen, la variation d'énergie cinétique d'un solide (ou d'une particule) entre deux positions est égale à la somme des travaux de toutes les forces appliquées : $\Delta E_c = E_{c,B} - E_{c,A} = \sum W_{A\to B}(\vec{F})$.
\item \textbf{Force électrique :} $\vec{F}_e = q\vec{E}$, avec $[F_e] = \mathrm{N}$, $[q] = \mathrm{C}$, $[E] = \mathrm{V/m} = \mathrm{N/C}$.\\
\textbf{Force magnétique :} $\vec{F}_m = q\,\vec{v}\wedge\vec{B}$, avec $[v] = \mathrm{m/s}$, $[B] = \mathrm{T}$, et $1~\mathrm{T} = 1~\mathrm{N/(A\,m)}$.
\item On a $[R] = \dfrac{[m][v]}{[q][B]}$. De $F_m = |q|vB$ on tire $[B] = \dfrac{[F]}{[q][v]} = \dfrac{\mathrm{kg\,m\,s^{-2}}}{\mathrm{C\,m\,s^{-1}}} = \mathrm{kg\,C^{-1}\,s^{-1}}$. Donc :
\[
[R] = \frac{\mathrm{kg}\times\mathrm{m\,s^{-1}}}{\mathrm{C}\times\mathrm{kg\,C^{-1}\,s^{-1}}} = \mathrm{m}.
\]
L'expression $R = \dfrac{mv}{|q|B}$ est bien homogène à une longueur.
\end{enumerate}
\end{corrige}

\begin{corrige}[Exercice 4 — Calculs directs]
\begin{enumerate}
\item Théorème de l'énergie cinétique : $\frac{1}{2}m_e v^2 = eU$, d'où :
\[
v = \sqrt{\frac{2eU}{m_e}} = \sqrt{\frac{2\times\num{1,60e-19}\times 500}{\num{9,11e-31}}} = \sqrt{\num{1,756e14}} \approx \num{1,33e7}~\mathrm{m/s}.
\]
\item Le proton entre avec $\vec{v}\perp\vec{B}$ : trajectoire circulaire de rayon
\[
R = \frac{m_p v}{eB} = \frac{\num{1,67e-27}\times\num{1,0e6}}{\num{1,60e-19}\times 0,10} = \frac{\num{1,67e-21}}{\num{1,60e-20}} \approx \num{0,104}~\mathrm{m} = 10,4~\mathrm{cm}.
\]
\item Période :
\[
T = \frac{2\pi R}{v} = \frac{2\pi m_p}{eB} = \frac{2\pi\times\num{1,67e-27}}{\num{1,60e-19}\times 0,10} \approx \num{6,57e-7}~\mathrm{s}.
\]
\textbf{Conclusion :} l'expression $T = \dfrac{2\pi m_p}{eB}$ ne contient pas $v$ : la période est \textbf{indépendante de la vitesse} du proton (régime non relativiste). Si $v$ double, $R$ double et $T$ reste inchangée.
\end{enumerate}
\end{corrige}

\begin{corrige}[Exercice 5 — Canon à électrons d'un oscilloscope]
\begin{enumerate}
\item \textbf{Schéma :} cathode $C$ à gauche (potentiel bas), anode $A$ à droite (potentiel haut, trouée pour laisser passer le faisceau). Le champ $\vec{E}$ pointe des hauts vers les bas potentiels : de $A$ vers $C$ (vers la gauche). L'électron portant la charge $q=-e$, la force $\vec{F}_e = -e\vec{E}$ est dirigée de $C$ vers $A$ (vers la droite) : c'est elle qui accélère l'électron.
\item Entre $C$ ($v_C \approx 0$) et $A$ ($v_A = v_0$), seule la force électrique travaille (poids négligeable) :
\[
\Delta E_c = W_{C\to A}(\vec{F}_e) \;\Longrightarrow\; \frac{1}{2}m_e v_0^2 - 0 = q(V_C - V_A) = (-e)(-U) = eU,
\]
d'où
\[
\boxed{v_0 = \sqrt{\frac{2eU}{m_e}}}.
\]
\item Application numérique :
\[
v_0 = \sqrt{\frac{2\times\num{1,60e-19}\times 1500}{\num{9,11e-31}}} = \sqrt{\num{5,27e14}} \approx \num{2,30e7}~\mathrm{m/s}.
\]
Soit environ $8\,\%$ de la vitesse de la lumière : le cadre non relativiste reste acceptable.
\item Comme $v_0 \propto \sqrt{U}$, doubler la vitesse ($v_0' = 2v_0$) exige $U' = 4U$ :
\[
U' = 4\times 1500 = 6000~\mathrm{V} = 6,0~\mathrm{kV}.
\]
\end{enumerate}
\textbf{Point de vigilance :} ne pas écrire $W = qU$ avec $U = V_C - V_A$ sans préciser le sens ; pour l'électron, c'est le signe de $q$ qui rend le travail moteur quand il se déplace vers le potentiel le plus élevé.
\end{corrige}

\begin{corrige}[Exercice 6 — Proton dans un champ magnétique]
\begin{enumerate}
\item \textbf{Comparaison des forces :}
\[
P = m_p g = \num{1,67e-27}\times 9,81 \approx \num{1,64e-26}~\mathrm{N},
\]
\[
F_m = ev_0B = \num{1,60e-19}\times\num{2,0e6}\times 0,20 = \num{6,4e-14}~\mathrm{N}.
\]
Rapport : $\dfrac{P}{F_m} \approx \num{2,6e-13} \ll 1$ : le poids est \textbf{négligeable} devant la force magnétique.
\item \textbf{Démonstration (exigible).} Référentiel terrestre supposé galiléen ; système : proton. Seule force : $\vec{F}_m = e\,\vec{v}\wedge\vec{B}$, perpendiculaire à $\vec{v}$. Dans le repère de Frenet, $\vec{a} = \dfrac{dv}{dt}\vec{\tau} + \dfrac{v^2}{R}\vec{n}$. La deuxième loi de Newton $m_p\vec{a} = \vec{F}_m$ projetée :
\begin{itemize}
\item sur $\vec{\tau}$ : $m_p\dfrac{dv}{dt} = 0 \Rightarrow v = v_0 = \text{cste}$ : mouvement \textbf{uniforme} ;
\item sur $\vec{n}$ : $m_p\dfrac{v_0^2}{R} = ev_0B \Rightarrow R = \dfrac{m_p v_0}{eB} = \text{cste}$ : trajectoire \textbf{circulaire}.
\end{itemize}
Le mouvement est donc circulaire uniforme.
\item \[
R = \frac{\num{1,67e-27}\times\num{2,0e6}}{\num{1,60e-19}\times 0,20} = \frac{\num{3,34e-21}}{\num{3,2e-20}} \approx \num{0,104}~\mathrm{m} = 10,4~\mathrm{cm}.
\]
\item \[
T = \frac{2\pi R}{v_0} = \frac{2\pi m_p}{eB} = \frac{2\pi\times\num{1,67e-27}}{\num{1,60e-19}\times 0,20} \approx \num{3,28e-7}~\mathrm{s}.
\]
\item \textbf{Schéma et sens de rotation :} $\vec{v}_0$ horizontal vers la droite, $\vec{B}$ sortant de la feuille ($\odot$). Règle de la main droite pour $\vec{v}_0\wedge\vec{B}$ : le résultat pointe \textbf{vers le haut}. Le proton étant positif ($q = +e$), $\vec{F}_m$ est aussi vers le haut : c'est la force centripète, donc le centre du cercle se trouve au-dessus de $O$. Le proton tourne dans le \textbf{sens trigonométrique} (inverse des aiguilles d'une montre).
\end{enumerate}
\textbf{Point de vigilance :} pour une charge \emph{négative}, il faut inverser le sens de la force obtenu par la règle de la main droite, donc le sens de rotation.
\end{corrige}

\begin{corrige}[Exercice 7 — Filtre de Wien]
\begin{enumerate}
\item \textbf{Schéma :} $\vec{v}_0$ horizontal vers la droite ; $\vec{E}$ vertical vers le haut ; $\vec{B}$ perpendiculaire au plan, \textbf{entrant} ($\otimes$), choisi de sorte que $\vec{v}_0\wedge\vec{B}$ pointe vers le bas. Pour un ion positif : $\vec{F}_e = q\vec{E}$ vers le haut, $\vec{F}_m = q\,\vec{v}_0\wedge\vec{B}$ vers le bas. Les deux forces sont opposées.
\item L'ion traverse le filtre sans déviation si la somme des forces est nulle (mouvement rectiligne uniforme) :
\[
\vec{F}_e + \vec{F}_m = \vec{0} \;\Longleftrightarrow\; qE = qv_0B \;\Longleftrightarrow\; \boxed{v_0 = \frac{E}{B}}.
\]
\item \[
B = \frac{E}{v_0} = \frac{\num{2,0e4}}{\num{4,0e5}} = \num{5,0e-2}~\mathrm{T} = 0,050~\mathrm{T}.
\]
\item \textbf{Si $v > v_0$ :} $F_m = qvB > qE = F_e$ : la force magnétique l'emporte ; la résultante est dirigée dans le sens de $\vec{F}_m$ (vers le \textbf{bas} sur le schéma) : l'ion est dévié vers le bas.\\
\textbf{Si $v < v_0$ :} $F_e > F_m$ : la force électrique l'emporte ; la résultante est dirigée dans le sens de $\vec{F}_e$ (vers le \textbf{haut}) : l'ion est dévié vers le haut.\\
Seuls les ions de vitesse exactement $v_0 = E/B$ traversent en ligne droite : d'où le nom de \emph{filtre de vitesse}.
\end{enumerate}
\textbf{Point de vigilance :} la condition $v_0 = E/B$ est indépendante de la charge $q$ et de la masse $m$ de l'ion : le filtre sélectionne une \emph{vitesse}, pas une espèce ionique.
\end{corrige}

\begin{corrige}[Exercice 8 — Déviation d'un électron entre deux plaques]
\begin{enumerate}
\item Champ entre les plaques (condensateur plan) :
\[
E' = \frac{U'}{d} = \frac{200}{\num{4,0e-2}} = \num{5,0e3}~\mathrm{V/m}.
\]
\item \textbf{Mise en équations.} Repère $(O,x,y)$ : $O$ à l'entrée des plaques, $Ox$ selon $\vec{v}_0$, $Oy$ vertical ascendant. La plaque du haut étant positive, $\vec{E}$ pointe vers le bas ; l'électron ($q=-e$) subit $\vec{F}_e = -e\vec{E}$ dirigée vers le haut, de norme $eE'$. Deuxième loi de Newton (poids négligeable) : $m_e\vec{a} = \vec{F}_e$, soit
\[
a_x = 0, \qquad a_y = \frac{eE'}{m_e}.
\]
Conditions initiales : $x(0)=y(0)=0$, $v_x(0)=v_0$, $v_y(0)=0$. Par intégrations successives :
\[
\boxed{x(t) = v_0\,t}, \qquad \boxed{y(t) = \frac{1}{2}\,\frac{eE'}{m_e}\,t^2}.
\]
\item En éliminant le temps ($t = x/v_0$) :
\[
y(x) = \frac{eE'}{2m_e v_0^2}\,x^2.
\]
C'est une \textbf{parabole} de sommet $O$, d'axe $Oy$, de concavité tournée vers le haut (vers la plaque positive).
\item Durée de passage :
\[
\Delta t = \frac{L}{v_0} = \frac{\num{5,0e-2}}{\num{2,0e7}} = \num{2,5e-9}~\mathrm{s}.
\]
Déviation à la sortie :
\[
y_S = \frac{1}{2}\,\frac{eE'}{m_e}\,(\Delta t)^2 = \frac{1}{2}\times\frac{\num{1,60e-19}\times\num{5,0e3}}{\num{9,11e-31}}\times(\num{2,5e-9})^2.
\]
\[
y_S = \frac{1}{2}\times\num{8,78e14}\times\num{6,25e-18} \approx \num{2,7e-3}~\mathrm{m} = 2,7~\mathrm{mm}.
\]
Vérification de cohérence : $y_S = 2,7~\mathrm{mm} < d/2 = 20~\mathrm{mm}$ : l'électron sort bien des plaques sans les toucher.
\end{enumerate}
\textbf{Point de vigilance :} toujours vérifier que $y_S < d/2$ ; sinon l'hypothèse « la particule ressort des plaques » est fausse et l'exercice change de nature.
\end{corrige}

\begin{corrige}[Exercice 9 — Type Bac : oscilloscope cathodique]
\textbf{Barème indicatif (sur 10 pts) :} 1.a : 0,5 ; 1.b : 1,5 ; 2.a : 1 ; 2.b : 2 ; 2.c : 1,5 ; 3.a : 1 ; 3.b : 1,5 ; 3.c : 1.

\begin{enumerate}
\item \textbf{Accélération.}
\begin{enumerate}
\item Entre $C$ et $A$, le champ $\vec{E}$ est dirigé de l'anode (potentiel le plus élevé) vers la cathode. La force sur l'électron $\vec{F}_e = -e\vec{E}$ est dirigée de $C$ vers $A$ : force accélératrice.
\item Théorème de l'énergie cinétique entre $C$ ($v_C \approx 0$) et $A$ ($v_A = v_0$) :
\[
\frac{1}{2}m_e v_0^2 = eU \;\Longrightarrow\; v_0 = \sqrt{\frac{2eU}{m_e}} = \sqrt{\frac{2\times\num{1,60e-19}\times 1000}{\num{9,11e-31}}} \approx \num{1,87e7}~\mathrm{m/s}.
\]
\end{enumerate}
\item \textbf{Déviation entre les plaques.}
\begin{enumerate}
\item $P_1$ (plaque du haut) est positive : $\vec{E}'$ pointe de $P_1$ vers $P_2$, donc vers le bas. L'électron, de charge négative, subit $\vec{F}_e = -e\vec{E}'$ dirigée \textbf{vers le haut} : le faisceau est dévié vers la plaque positive $P_1$.
\item Avec $a_x = 0$ et $a_y = \dfrac{eE'}{m_e} = \dfrac{eU'}{m_e d}$, et les conditions initiales en $O$ :
\[
x(t) = v_0 t, \qquad y(t) = \frac{1}{2}\,\frac{eU'}{m_e d}\,t^2, \qquad y(x) = \frac{eU'}{2m_e d\, v_0^2}\,x^2 \quad (\text{parabole}).
\]
\item Durée de transit : $\Delta t = L/v_0 = \dfrac{0,060}{\num{1,87e7}} \approx \num{3,2e-9}$ s.
\[
y_S = \frac{eU'L^2}{2m_e d\, v_0^2} = \frac{\num{1,60e-19}\times 150\times(0,060)^2}{2\times\num{9,11e-31}\times 0,030\times(\num{1,87e7})^2} \approx \num{4,5e-3}~\mathrm{m} = 4,5~\mathrm{mm}.
\]
Composante verticale de la vitesse de sortie :
\[
v_{yS} = \frac{eU'L}{m_e d\, v_0} \approx \num{2,8e6}~\mathrm{m/s}, \qquad \tan\alpha = \frac{v_{yS}}{v_0} = \frac{eU'L}{m_e d\, v_0^2} \approx 0,150,
\]
soit $\alpha \approx 8,5^\circ$.
\end{enumerate}
\item \textbf{Déflexion sur l'écran.}
\begin{enumerate}
\item À la sortie des plaques, le champ électrique est nul (on néglige l'effet de bord), le champ magnétique terrestre est négligé et le poids est négligeable : $\sum\vec{F}_{\text{ext}} = \vec{0}$. D'après le \textbf{principe d'inertie} (première loi de Newton), le mouvement de l'électron est rectiligne uniforme, à la vitesse de sortie $\vec{v}_S$.
\item La tangente à la trajectoire en $S$ coupe l'axe $Ox$ au point $I$, milieu des plaques (propriété de la parabole : la tangente en $x=L$ coupe l'axe en $x = L/2$). Entre $I$ et l'écran, distance horizontale $D$, l'électron monte de $D\tan\alpha$. D'où :
\[
Y = \left(\frac{L}{2} + D\right)\tan\alpha = \frac{eU'L}{m_e d\, v_0^2}\left(\frac{L}{2} + D\right).
\]
Application numérique :
\[
Y = \frac{\num{1,60e-19}\times 150\times 0,060}{\num{9,11e-31}\times(\num{1,87e7})^2\times 0,030}\times(0,030 + 0,25) \approx \num{4,2e-2}~\mathrm{m} = 4,2~\mathrm{cm}.
\]
\item Comme $v_0^2 = \dfrac{2eU}{m_e}$, on a $Y \propto \dfrac{1}{v_0^2} \propto \dfrac{1}{U}$. Si $U$ double, $Y$ est \textbf{divisé par deux} : $Y' \approx 2,1~\mathrm{cm}$.\\
\textbf{Commentaire :} une tension accélératrice élevée donne un faisceau rapide et énergétique (spot lumineux brillant, bien piqué) mais une \emph{sensibilité} de déviation plus faible ; une tension faible donne une grande déflexion mais un spot moins lumineux. Le constructeur cherche un compromis.
\end{enumerate}
\end{enumerate}
\textbf{Points de vigilance (rédaction Bac) :} citer le référentiel (terrestre, supposé galiléen) ; justifier la négligence du poids ; préciser les conditions initiales avant d'intégrer ; citer le principe d'inertie pour la phase hors champ ; donner chaque résultat avec son unité et un nombre cohérent de chiffres significatifs.
\end{corrige}

\begin{corrige}[Exercice 10 — Type Bac : spectrographe de masse et isotopes du brome]
\textbf{Barème indicatif (sur 10 pts) :} 1 : 1,5 ; 2 : 3 ; 3 : 2 ; 4 : 1,5 ; 5 : 2.

\begin{enumerate}
\item \textbf{Vitesse à l'entrée de la chambre de déviation.} Théorème de l'énergie cinétique entre la chambre d'ionisation ($v_i \approx 0$) et l'entrée de la chambre de déviation, l'ion \ce{Br+} ($q = +e$) traversant la tension $U$ :
\[
\frac{1}{2}mv^2 = eU \;\Longrightarrow\; \boxed{v = \sqrt{\frac{2eU}{m}}}.
\]
\item \textbf{Nature du mouvement.} Dans la chambre, l'ion n'est soumis qu'à la force magnétique $\vec{F}_m = e\,\vec{v}\wedge\vec{B}$ (poids négligeable), perpendiculaire à $\vec{v}$. Son travail est nul : la vitesse garde une norme constante (mouvement uniforme). Comme $\vec{v}\perp\vec{B}$, la force est centripète, de norme constante $evB$ ; la deuxième loi de Newton projetée sur la normale de Frenet donne :
\[
evB = m\frac{v^2}{R} \;\Longrightarrow\; R = \frac{mv}{eB} = \text{cste} : \text{ trajectoire circulaire.}
\]
En injectant l'expression de $v$ :
\[
R = \frac{m}{eB}\sqrt{\frac{2eU}{m}} = \boxed{\frac{1}{B}\sqrt{\frac{2mU}{e}}}.
\]
\item \textbf{Rayons.} Masses : $m_1 = 79u = 79\times\num{1,66e-27} = \num{1,311e-25}$ kg ; $m_2 = 81u = \num{1,345e-25}$ kg.
\[
R_1 = \frac{1}{0,50}\sqrt{\frac{2\times\num{1,311e-25}\times 4000}{\num{1,60e-19}}} = 2\sqrt{\num{6,56e-3}} \approx \num{0,162}~\mathrm{m},
\]
\[
R_2 = \frac{1}{0,50}\sqrt{\frac{2\times\num{1,345e-25}\times 4000}{\num{1,60e-19}}} = 2\sqrt{\num{6,72e-3}} \approx \num{0,164}~\mathrm{m}.
\]
On vérifie que $R_2 > R_1$ : l'isotope le plus lourd est le moins dévié.
\item \textbf{Écart entre les impacts.} Chaque ion décrit un demi-cercle : l'impact se situe à un diamètre ($2R$) du point d'entrée. Donc :
\[
\Delta = 2(R_2 - R_1) = 2\times(0,164 - 0,162) \approx \num{4,0e-3}~\mathrm{m} = 4,0~\mathrm{mm}.
\]
\item \textbf{Influence des réglages.} De $R \propto \sqrt{mU}/B$, on tire $\Delta \propto \dfrac{\sqrt{U}}{B}$ :
\begin{itemize}
\item augmenter $U$ \textbf{augmente} $\Delta$ (en $\sqrt{U}$) : meilleure séparation ;
\item augmenter $B$ \textbf{diminue} $\Delta$ (en $1/B$) : séparation moins bonne.
\end{itemize}
\textbf{Conclusion :} pour bien séparer les isotopes, on privilégie une tension d'accélération élevée et un champ magnétique modéré, tout en gardant des rayons compatibles avec les dimensions de la chambre et une détection précise des impacts.
\end{enumerate}
\textbf{Points de vigilance (rédaction Bac) :} la démonstration du caractère \emph{uniforme} du mouvement (travail nul de $\vec{F}_m$) est attendue explicitement ; ne pas oublier le facteur 2 entre rayon et diamètre dans le calcul de $\Delta$ ; convertir les masses en kilogrammes avant tout calcul numérique.
\end{corrige}

\begin{corrige}[Exercice 11 — Type Bac : cyclotron et filtre de Wien en médecine]
\textbf{Barème indicatif (sur 10 pts) :} 1.a : 1 ; 1.b : 2 ; 1.c : 1,5 ; 1.d : 2 ; 2.a : 1,5 ; 2.b : 2.

\begin{enumerate}
\item \textbf{Étude du cyclotron.}
\begin{enumerate}
\item \textbf{Rôle du champ magnétique :} la force de Lorentz $\vec{F}_m = e\,\vec{v}\wedge\vec{B}$, perpendiculaire à la vitesse, courbe la trajectoire des protons en demi-cercles à l'intérieur des dees (où le champ électrique est nul) : elle les ramène périodiquement dans l'intervalle entre les dees sans modifier leur énergie.\\
\textbf{Rôle de la tension alternative :} elle crée dans l'intervalle entre les dees un champ électrique dont le sens s'inverse périodiquement ; si sa fréquence est bien choisie, le champ est toujours accélérateur à chaque passage du proton : l'énergie cinétique (donc la vitesse et le rayon) augmente à chaque traversée — la trajectoire est une spirale.
\item Dans un dee, $\vec{E} = \vec{0}$ : seule la force magnétique agit, centripète. La deuxième loi de Newton projetée sur la normale de Frenet :
\[
evB = m_p\frac{v^2}{R} \;\Longrightarrow\; R = \frac{m_p v}{eB}.
\]
La période de révolution vaut alors :
\[
T = \frac{2\pi R}{v} = \boxed{\frac{2\pi m_p}{eB}}.
\]
\textbf{Constat :} $T$ ne dépend ni de $v$ ni de $R$ (isochronisme des révolutions, en régime non relativiste) : c'est ce qui permet une tension alternative de fréquence fixe.
\item Pour que le proton soit accéléré à chaque passage, la tension doit changer de signe à chaque demi-tour : sa période propre doit égaler $T$, donc
\[
f = \frac{1}{T} = \frac{eB}{2\pi m_p} = \frac{\num{1,60e-19}\times 1,0}{2\pi\times\num{1,67e-27}} \approx \num{1,52e7}~\mathrm{Hz} = 15,2~\mathrm{MHz}.
\]
\item Au rayon maximal, $v_{\max} = \dfrac{eBR_{\max}}{m_p}$, d'où :
\[
E_{c,\max} = \frac{1}{2}m_p v_{\max}^2 = \boxed{\frac{e^2B^2R_{\max}^2}{2m_p}}.
\]
Application numérique :
\[
E_{c,\max} = \frac{(\num{1,60e-19})^2\times(1,0)^2\times(0,50)^2}{2\times\num{1,67e-27}} = \frac{\num{6,4e-39}}{\num{3,34e-27}} \approx \num{1,9e-12}~\mathrm{J}.
\]
Conversion : $E_{c,\max} = \dfrac{\num{1,9e-12}}{\num{1,60e-19}} \approx \num{1,2e7}~\mathrm{eV} = 12~\mathrm{MeV}$, énergie typique des cyclotrons produisant des radioisotopes pour l'imagerie médicale (TEP).
\end{enumerate}
\item \textbf{Sélection par filtre de Wien.}
\begin{enumerate}
\item Condition de traversée sans déviation : $eE = ev_0B'$, soit
\[
B' = \frac{E}{v_0} = \frac{\num{2,0e5}}{\num{4,0e6}} = \num{5,0e-2}~\mathrm{T} = 0,050~\mathrm{T}.
\]
\item Si $v > v_0$ : $F_m = evB' > eE = F_e$ : \textbf{la force magnétique l'emporte}. Pour qu'à $v = v_0$ les forces se compensent avec $\vec{F}_e$ dirigée vers le bas, $\vec{F}_m$ est dirigée vers le haut. Pour $v > v_0$, la force résultante $\vec{F}_e + \vec{F}_m$ est donc dirigée \textbf{vers le haut} : le proton est dévié vers la plaque du haut et ne franchit pas le diaphragme de sortie. Seuls les protons de vitesse $v_0$ sont transmis.
\end{enumerate}
\end{enumerate}
\textbf{Points de vigilance (rédaction Bac) :} distinguer clairement les deux zones (dans un dee : $\vec{E}=\vec{0}$, trajectoire circulaire ; entre les dees : accélération électrique) ; la fréquence de la tension alternative est la fréquence cyclotron $f = eB/(2\pi m_p)$, indépendante de $v$ ; dans la conversion J $\to$ MeV, utiliser $1~\mathrm{MeV} = \num{1,60e-13}$ J ; pour le filtre, raisonner sur les \emph{normes} des forces puis conclure sur le sens de déviation à partir du schéma.
\end{corrige}

\newpage

\coursec{Fiche bilan}

\begin{popfigure}
\begin{tikzpicture}[mindmap, grow cyclic, every node/.style={concept}, concept color=popblue,
  level 1/.append style={level distance=3.6cm, sibling angle=60},
  level 2/.append style={level distance=2.2cm},
  text width=2.6cm, align=center, font=\small]
\node[concept color=popdark, text=white, font=\small\bfseries]{Particule chargée dans $\vec{E}$ et $\vec{B}$}
  child[concept color=popblue] { node {Accélération par $\vec{E}$ uniforme}
    child { node[concept color=popblueL, text width=2.4cm] {$\tfrac12 mv^2 = |q|U$} }
  }
  child[concept color=poporange] { node {Déviation électrique}
    child { node[concept color=poporangeL, text width=2.4cm] {trajectoire parabolique} }
  }
  child[concept color=popgold] { node {Déflexion sur écran}
    child { node[concept color=popgoldL, text width=2.4cm] {$Y = \dfrac{q\ell L}{mv_0^2 d}\,U$} }
  }
  child[concept color=popgreen] { node {Champ $\vec{B}$ uniforme, $\vec{v}_0 \perp \vec{B}$}
    child { node[concept color=popgreenL, text width=2.4cm] {MCU : $R=\dfrac{mv}{|q|B}$} }
  }
  child[concept color=poppurple] { node {Applications}
    child { node[concept color=poppurpleL, text width=2.4cm] {spectrographe de masse, cyclotron} }
  }
  child[concept color=poppink] { node {Filtre de Wien}
    child { node[concept color=poppinkL, text width=2.4cm] {$v = \dfrac{E}{B}$} }
  };
\end{tikzpicture}
\legende{Carte mentale de la leçon : les six situations clés et leurs formules phares.}
\end{popfigure}

\begin{bilan}
\textbf{Définitions clés}
\begin{itemize}
  \item \cle{Force électrique} : $\vec{F}_e = q\vec{E}$ (parallèle à $\vec{E}$ si $q>0$, opposée si $q<0$).
  \item \cle{Force de Lorentz magnétique} : $\vec{F}_m = q\,\vec{v}\wedge\vec{B}$, toujours perpendiculaire à $\vec{v}$ : elle ne \textbf{travaille pas} et ne modifie pas la norme de la vitesse.
  \item \cle{Champ uniforme} entre deux armatures planes : $E = \dfrac{U}{d}$.
\end{itemize}

\textbf{Formules à connaître par c\oe ur}
\begin{center}
\begin{tabularx}{\linewidth}{|l|X|X|}
\hline
\rowcolor{popblueL} \textbf{Situation} & \textbf{Résultat} & \textbf{Condition} \\
\hline
Accélération (canon à électrons) & $\dfrac{1}{2}mv^2 - \dfrac{1}{2}mv_0^2 = qU$ & théorème de l'énergie cinétique \\
\hline
Déviation électrique & $y = \dfrac{qE}{2mv_0^2}\,x^2$ (parabole) & $\vec{v}_0 \perp \vec{E}$, poids négligé \\
\hline
Déflexion sur écran & $Y = \dfrac{qE\ell L}{mv_0^2 d}$ ou $Y=\dfrac{q\ell L}{mv_0^2 d}U$ & $L$ : distance plaques--écran \\
\hline
Champ magnétique & $R = \dfrac{mv}{|q|B}$ ; $T = \dfrac{2\pi m}{|q|B}$ & $\vec{v}_0 \perp \vec{B}$, MCU \\
\hline
Filtre de Wien & $v = \dfrac{E}{B}$ & $\vec{E}$, $\vec{B}$, $\vec{v}$ croisés \\
\hline
Cyclotron & $f = \dfrac{|q|B}{2\pi m}$ (indépendante de $v$) & dés, tension alternative \\
\hline
\end{tabularx}
\end{center}

\textbf{Méthodes express}
\begin{enumerate}
  \item \textbf{Accélération} : écrire le théorème de l'énergie cinétique entre les deux plaques $\Rightarrow$ $v = \sqrt{2|q|U/m}$ (si $v_0 \approx 0$).
  \item \textbf{Déviation} : appliquer la 2\up{e} loi de Newton, projeter sur $(Ox)$ et $(Oy)$, intégrer deux fois, éliminer $t$ $\Rightarrow$ équation $y(x)$.
  \item \textbf{Magnétique} : montrer que $\vec{F}_m \perp \vec{v}$ $\Rightarrow$ $v$ constante, mouvement circulaire ; identifier la force à la force centripète : $|q|vB = \dfrac{mv^2}{R}$.
  \item \textbf{Vérifier} toujours : signe de $q$ (sens de la déviation), homogénéité des formules, poids négligeable devant $|q|E$ ou $|q|vB$.
\end{enumerate}
\end{bilan}

\begin{aretenir}[Checklist avant le Bac]
\begin{itemize}
  \item $\square$ Je sais exprimer la vitesse acquise par une particule accélérée sous une tension $U$.
  \item $\square$ Je sais établir les équations horaires $x(t)$ et $y(t)$ dans un champ $\vec{E}$ uniforme.
  \item $\square$ Je sais obtenir l'équation de la trajectoire parabolique $y(x)$ par élimination du temps.
  \item $\square$ Je sais calculer la déflexion $Y$ sur l'écran d'un oscilloscope.
  \item $\square$ Je sais justifier que le mouvement dans $\vec{B}$ uniforme (avec $\vec{v}_0\perp\vec{B}$) est circulaire et uniforme.
  \item $\square$ Je sais démontrer $R = \dfrac{mv}{|q|B}$ et $T = \dfrac{2\pi m}{|q|B}$.
  \item $\square$ Je sais que la force magnétique ne travaille pas : l'énergie cinétique est conservée.
  \item $\square$ Je sais expliquer le principe du spectrographe de masse (séparation d'isotopes selon $R$).
  \item $\square$ Je sais expliquer le fonctionnement du cyclotron (fréquence indépendante de $v$).
  \item $\square$ Je sais trouver la condition du filtre de Wien : $v = E/B$.
  \item $\square$ Je vérifie le signe de la charge et le sens des vecteurs sur mon schéma.
  \item $\square$ Je contrôle l'homogénéité et j'écris le résultat avec son unité SI et des chiffres significatifs cohérents.
\end{itemize}
\end{aretenir}

\findecours

\end{document}
